% ============================================================
%  Topologie Générale — Notes de Cours
%  Licence L2–L3, Mathématiques
%  Partie 1 : Chapitres 1–2
% ============================================================
\documentclass[12pt,a4paper,openany]{book}

% --- Langues et encodage ---
\usepackage[utf8]{inputenc}
\usepackage[T1]{fontenc}
\usepackage[french]{babel}

% --- Géométrie ---
\usepackage[margin=2.5cm,top=3cm,bottom=3cm]{geometry}

% --- Mathématiques ---
\usepackage{amsmath,amssymb,amsthm,mathrsfs}

% --- Graphiques ---
\usepackage{tikz}
\usetikzlibrary{arrows.meta,calc,positioning,patterns,decorations.markings,cd}

% --- Boîtes colorées ---
\usepackage[most]{tcolorbox}
\tcbuselibrary{theorems,skins,breakable}

% --- Liens hypertexte ---
\usepackage[colorlinks=true,linkcolor=blue!60!black,citecolor=green!50!black,urlcolor=blue!70!black]{hyperref}

% --- Index ---
\usepackage{makeidx}
\makeindex

% --- Listes ---
\usepackage{enumitem}

% --- En-têtes ---
\usepackage{fancyhdr}
\pagestyle{fancy}
\fancyhf{}
\fancyhead[LE]{\leftmark}
\fancyhead[RO]{\rightmark}
\fancyfoot[C]{\thepage}
\renewcommand{\headrulewidth}{0.4pt}

% --- Divers ---
\usepackage{xcolor}
\usepackage{microtype}
\usepackage{caption}

% ============================================================
%  Environnements de théorèmes
% ============================================================

% Définition (amsthm + tcolorbox wrapping)
\theoremstyle{definition}
\newtheorem{definition}{Définition}[chapter]

\tcolorboxenvironment{definition}{%
  breakable,
  enhanced,
  colback=blue!5!white,
  colframe=blue!60!black,
  fonttitle=\bfseries,
  before skip=12pt,
  after skip=12pt,
  left=8pt,right=8pt,top=6pt,bottom=6pt
}

\newtheoremstyle{cours}
  {10pt}{10pt}{\itshape}{}{\bfseries}{.}{.5em}{}

\theoremstyle{cours}
\newtheorem{theoreme}{Théorème}[chapter]
\newtheorem{proposition}[theoreme]{Proposition}
\newtheorem{lemme}[theoreme]{Lemme}
\newtheorem{corollaire}[theoreme]{Corollaire}

\newtheoremstyle{coursnormal}
  {10pt}{10pt}{}{}{\bfseries}{.}{.5em}{}

\theoremstyle{coursnormal}
\newtheorem{remarque}[theoreme]{Remarque}
\newtheorem{exemple}[theoreme]{Exemple}
\newtheorem{exercice}{Exercice}[chapter]

% ============================================================
%  Commandes personnalisées
% ============================================================
\newcommand{\N}{\mathbb{N}}
\newcommand{\Z}{\mathbb{Z}}
\newcommand{\Q}{\mathbb{Q}}
\newcommand{\R}{\mathbb{R}}
\newcommand{\C}{\mathbb{C}}
\newcommand{\Int}[1]{\mathring{#1}}
\newcommand{\Cl}[1]{\overline{#1}}
\newcommand{\Fr}[1]{\partial #1}
\newcommand{\Vois}[1]{\mathcal{V}(#1)}
\newcommand{\Ouv}[1]{\mathcal{O}_{#1}}
\newcommand{\tomark}[1]{\texorpdfstring{#1}{}}
\newcommand{\T}{\mathcal{T}}
\newcommand{\B}{\mathcal{B}}

% ============================================================
\begin{document}

% ============================================================
%  Page de titre
% ============================================================
\begin{titlepage}
\begin{tikzpicture}[remember picture, overlay]
  \fill[left color=blue!15!white, right color=blue!5!white]
    (current page.south west) rectangle (current page.north east);
  \fill[color=orange!70!yellow]
    ([yshift=-2cm]current page.north west) rectangle ([yshift=-2.3cm]current page.north east);
  \fill[color=blue!80!black, rounded corners=8pt]
    ([xshift=3cm, yshift=-4cm]current page.north west) rectangle
    ([xshift=-3cm, yshift=-10cm]current page.north east);
  \node[anchor=center, text=white, font=\Huge\bfseries]
    at ([yshift=-6cm]current page.north) {Topologie G\'en\'erale};
  \node[anchor=center, text=orange!80!yellow, font=\Large]
    at ([yshift=-7.5cm]current page.north) {Notes de Cours};
  \node[anchor=center, text=white!80, font=\large]
    at ([yshift=-8.8cm]current page.north) {L2--L3 --- 2025--2026};
  \node[anchor=center, text=blue!80!black, font=\Large\itshape]
    at ([yshift=-12cm]current page.north) {Ya\'e Ulrich Gaba};
  \draw[orange!70!yellow, line width=2pt]
    ([xshift=5cm, yshift=-13.5cm]current page.north west) --
    ([xshift=-5cm, yshift=-13.5cm]current page.north east);
  \node[anchor=center, text width=12cm, align=center, text=blue!60!black, font=\itshape]
    at ([yshift=-16cm]current page.north) {<<~La g\'eom\'etrie est l'art de raisonner juste sur des figures fausses.~>> --- Henri Poincar\'e};
  \node[anchor=south, text=gray, font=\small]
    at ([yshift=1.5cm]current page.south) {\today};
\end{tikzpicture}
\end{titlepage}

% ============================================================
%  Table des matières
% ============================================================
\tableofcontents

% ============================================================
%  Préface
% ============================================================
\chapter*{Préface}
\addcontentsline{toc}{chapter}{Préface}
\markboth{Préface}{Préface}

L'étude de la topologie générale constitue un passage essentiel dans la formation mathématique. En analyse, l'étudiant rencontre d'abord la topologie de~$\R$ puis celle de~$\R^n$, à travers les notions de suites convergentes, d'ouverts et de fermés définis au moyen de la distance euclidienne. Rapidement, il s'aperçoit que nombre de raisonnements --- continuité, compacité, connexité --- ne reposent pas sur la structure métrique elle-même, mais sur la seule donnée des ouverts. C'est cette observation qui motive le passage à l'abstraction.

\medskip

La topologie générale propose un cadre unifié dans lequel les espaces métriques ne sont qu'un cas particulier. Ce cadre permet de traiter sur un pied d'égalité des espaces aussi différents que les espaces de fonctions munis de topologies faibles, les espaces de Zariski en géométrie algébrique, ou les espaces de Stone--Čech en analyse fonctionnelle. La généralité du formalisme, loin d'être un exercice d'abstraction gratuite, se révèle indispensable dès que l'on souhaite donner des démonstrations véritablement intrinsèques.

\medskip

Ces notes s'inscrivent dans la tradition initiée par Nicolas Bourbaki dans les \emph{Éléments de mathématique}, dont le chapitre consacré à la topologie générale demeure une référence incontournable. Nous adoptons les conventions de Bourbaki : les définitions sont précises, les démonstrations complètes, et les notations systématiques. Cependant, nous avons fait le choix d'accompagner la théorie de nombreux exemples, contre-exemples et exercices, afin de rendre le texte accessible à un public de Licence L2--L3.

\medskip

Le plan du cours suit une progression classique. Nous partons de la définition d'un espace topologique et de ses premières conséquences (ouverts, fermés, intérieur, adhérence), pour construire progressivement les outils fondamentaux : bases et sous-bases, continuité, topologie produit, séparation, compacité, connexité, et enfin espaces complets et théorèmes de Baire. Chaque chapitre se termine par un résumé et une série d'exercices classés par difficulté.

\medskip

Le lecteur est supposé maîtriser la théorie des ensembles (relations, applications, cardinalité) ainsi que les bases de l'analyse réelle (suites, séries, continuité dans~$\R$). Aucune connaissance préalable de topologie abstraite n'est requise.

\vfill

\hfill\textit{L'auteur}

% ============================================================
%  Notations
% ============================================================
\chapter*{Notations}
\addcontentsline{toc}{chapter}{Notations}
\markboth{Notations}{Notations}

Nous rassemblons ici les notations utilisées tout au long de ce cours.

\bigskip

\begin{tabular}{ll}
\hline
\textbf{Notation} & \textbf{Signification} \\
\hline
$(X,\tau)$ & Espace topologique \\
$\tau$ ou $\Ouv{X}$ & Topologie (ensemble des ouverts) de $X$ \\
$\Int{A}$ ou $A^{\circ}$ & Intérieur de $A$ \\
$\Cl{A}$ & Adhérence (fermeture) de $A$ \\
$\Fr{A}$ & Frontière de $A$ \\
$A'$ & Ensemble dérivé de $A$ \\
$\Vois{x}$ & Système de voisinages de $x$ \\
$B(x,r)$ & Boule ouverte de centre $x$ et de rayon $r$ \\
$\overline{B}(x,r)$ & Boule fermée de centre $x$ et de rayon $r$ \\
$\N, \Z, \Q, \R, \C$ & Ensembles de nombres usuels \\
$\mathcal{P}(X)$ & Ensemble des parties de $X$ \\
$A^c$ ou $X \setminus A$ & Complémentaire de $A$ dans $X$ \\
$\tau_1 \subset \tau_2$ & $\tau_1$ est moins fine que $\tau_2$ \\
$\tau_d$ & Topologie discrète \\
$\tau_i$ & Topologie indiscrète (grossière) \\
$\tau_{\mathrm{cof}}$ & Topologie cofinie \\
$\tau_{\mathrm{cod}}$ & Topologie codénombrable \\
\hline
\end{tabular}

% ============================================================
%  CHAPITRE 1
% ============================================================
\chapter{Espaces topologiques --- Définitions et axiomes}
\label{chap:espaces-topologiques}

\section*{Introduction}
\addcontentsline{toc}{section}{Introduction}

Dans un cours d'analyse, on définit les ouverts de~$\R^n$ à l'aide de la distance euclidienne~$d$ : une partie $U \subset \R^n$ est ouverte si, pour tout $x \in U$, il existe $r > 0$ tel que la boule ouverte
\[
  B(x,r) = \{y \in \R^n : d(x,y) < r\}
\]
soit contenue dans $U$. On vérifie alors que la famille~$\tau_d$ de tous les ouverts satisfait trois propriétés remarquables :
\begin{enumerate}[label=(\roman*)]
  \item $\varnothing$ et $\R^n$ appartiennent à $\tau_d$ ;
  \item toute réunion (même infinie) d'éléments de $\tau_d$ appartient à $\tau_d$ ;
  \item toute intersection \emph{finie} d'éléments de $\tau_d$ appartient à $\tau_d$.
\end{enumerate}

Ce sont ces trois propriétés, et elles seules, qui sous-tendent la plupart des raisonnements topologiques. L'idée fondatrice de la topologie générale est de prendre ces propriétés comme \emph{axiomes}, en oubliant la distance qui leur a donné naissance. On obtient ainsi un cadre beaucoup plus souple, applicable à des situations où aucune distance naturelle n'est disponible.

% ============================================================
\section{Définition d'un espace topologique}
\label{sec:def-espace-topo}

\begin{definition}[Topologie]\label{def:topologie}
  Soit $X$ un ensemble. On appelle \textbf{topologie}\index{topologie} sur $X$ une famille $\tau \subset \mathcal{P}(X)$ vérifiant les trois axiomes suivants :
  \begin{enumerate}[label=\textbf{(O\arabic*)}]
    \item\label{ax:O1} $\varnothing \in \tau$ et $X \in \tau$.
    \item\label{ax:O2} Si $(U_i)_{i \in I}$ est une famille quelconque d'éléments de $\tau$, alors $\displaystyle\bigcup_{i \in I} U_i \in \tau$.
    \item\label{ax:O3} Si $U_1, U_2, \ldots, U_n \in \tau$ (famille finie), alors $\displaystyle\bigcap_{k=1}^{n} U_k \in \tau$.
  \end{enumerate}
  Le couple $(X,\tau)$ est appelé \textbf{espace topologique}\index{espace topologique}. Les éléments de $\tau$ sont appelés les \textbf{ouverts}\index{ouvert} de $(X,\tau)$.
\end{definition}

\begin{remarque}
  L'axiome~\ref{ax:O2} autorise des réunions indexées par un ensemble $I$ quelconque, éventuellement non dénombrable, voire vide (auquel cas la réunion vaut $\varnothing$). En revanche, l'axiome~\ref{ax:O3} ne porte que sur des intersections \emph{finies}. Cette dissymétrie est essentielle : nous verrons (exemple~\ref{ex:intersection-infinie-ouverts}) que l'intersection dénombrable d'ouverts n'est pas nécessairement ouverte.
\end{remarque}

% ============================================================
\section{Exemples fondamentaux}
\label{sec:exemples-topologies}

\begin{exemple}[Topologie discrète]\label{ex:topo-discrete}
  \index{topologie!discrète}
  Soit $X$ un ensemble quelconque. La famille $\tau_d = \mathcal{P}(X)$ (toutes les parties de $X$) est une topologie sur $X$, appelée \textbf{topologie discrète}. En effet, les trois axiomes sont trivialement satisfaits puisque toute partie est ouverte.

  Dans cette topologie, tout singleton $\{x\}$ est ouvert, et donc toute partie de $X$ est à la fois ouverte et fermée.
\end{exemple}

\begin{exemple}[Topologie indiscrète]\label{ex:topo-indiscrete}
  \index{topologie!indiscrète}
  La famille $\tau_i = \{\varnothing, X\}$ est une topologie sur $X$, appelée \textbf{topologie indiscrète} (ou \textbf{grossière}). C'est la topologie la plus pauvre possible : les seuls ouverts sont l'ensemble vide et $X$ tout entier.
\end{exemple}

\begin{exemple}[Topologie cofinie]\label{ex:topo-cofinie}
  \index{topologie!cofinie}
  Soit $X$ un ensemble. On pose
  \[
    \tau_{\mathrm{cof}} = \{\varnothing\} \cup \{U \subset X : X \setminus U \text{ est fini}\}.
  \]
  Vérifions que $\tau_{\mathrm{cof}}$ est une topologie.
  \begin{proof}
    \ref{ax:O1} : $\varnothing \in \tau_{\mathrm{cof}}$ par définition, et $X \in \tau_{\mathrm{cof}}$ car $X \setminus X = \varnothing$ est fini.

    \ref{ax:O2} : Soit $(U_i)_{i \in I}$ une famille d'éléments de $\tau_{\mathrm{cof}}$. Si tous les $U_i$ sont vides, leur réunion est $\varnothing \in \tau_{\mathrm{cof}}$. Sinon, il existe $i_0 \in I$ tel que $U_{i_0} \neq \varnothing$, et alors $X \setminus U_{i_0}$ est fini. Or
    \[
      X \setminus \bigcup_{i \in I} U_i = \bigcap_{i \in I}(X \setminus U_i) \subset X \setminus U_{i_0},
    \]
    donc $X \setminus \bigcup_{i \in I} U_i$ est fini, et la réunion est dans $\tau_{\mathrm{cof}}$.

    \ref{ax:O3} : Soient $U_1, \ldots, U_n \in \tau_{\mathrm{cof}}$. Si l'un des $U_k$ est vide, l'intersection est vide, donc dans $\tau_{\mathrm{cof}}$. Sinon, chaque $X \setminus U_k$ est fini, et
    \[
      X \setminus \bigcap_{k=1}^n U_k = \bigcup_{k=1}^n (X \setminus U_k)
    \]
    est une réunion finie d'ensembles finis, donc finie. Ainsi $\bigcap_{k=1}^n U_k \in \tau_{\mathrm{cof}}$.
  \end{proof}

  Lorsque $X$ est infini, cette topologie n'est pas métrisable. Lorsque $X$ est fini, $\tau_{\mathrm{cof}} = \mathcal{P}(X) = \tau_d$ coïncide avec la topologie discrète.
\end{exemple}

\begin{exemple}[Topologie codénombrable]\label{ex:topo-codenombrable}
  \index{topologie!codénombrable}
  Sur un ensemble $X$ non dénombrable, on peut définir
  \[
    \tau_{\mathrm{cod}} = \{\varnothing\} \cup \{U \subset X : X \setminus U \text{ est dénombrable}\}.
  \]
  La vérification des axiomes est analogue à celle de la topologie cofinie. Si $X = \R$, cette topologie est strictement plus fine que $\tau_{\mathrm{cof}}$ et strictement moins fine que $\tau_d$.
\end{exemple}

\begin{exemple}[Topologie métrique]\label{ex:topo-metrique}
  \index{topologie!métrique}
  Soit $(X,d)$ un espace métrique. La topologie $\tau_d$ engendrée par les boules ouvertes
  \[
    B(x,r) = \{y \in X : d(x,y) < r\}, \quad x \in X,\; r > 0,
  \]
  est définie par : $U \in \tau_d$ si et seulement si pour tout $x \in U$, il existe $r > 0$ tel que $B(x,r) \subset U$. On vérifie les trois axiomes :
  \begin{proof}
    \ref{ax:O1} : $\varnothing$ satisfait la condition vacuement. Pour $X$, tout point $x$ admet $B(x,1) \subset X$.

    \ref{ax:O2} : Si chaque $U_i$ est ouvert et $x \in \bigcup_i U_i$, alors $x \in U_{i_0}$ pour un certain $i_0$, et il existe $r > 0$ tel que $B(x,r) \subset U_{i_0} \subset \bigcup_i U_i$.

    \ref{ax:O3} : Si $x \in U_1 \cap \cdots \cap U_n$, pour chaque $k$ il existe $r_k > 0$ tel que $B(x,r_k) \subset U_k$. En posant $r = \min(r_1,\ldots,r_n) > 0$, on a $B(x,r) \subset U_1 \cap \cdots \cap U_n$.
  \end{proof}
  La topologie usuelle de $\R^n$ est la topologie métrique associée à la distance euclidienne.
\end{exemple}

\begin{exemple}[Topologie de Zariski sur $\R$]\label{ex:topo-zariski}
  \index{topologie!de Zariski}
  Sur $\R$ (ou plus généralement sur un corps $k$), on définit la \textbf{topologie de Zariski} en déclarant fermés les ensembles de zéros de familles de polynômes. Dans le cas de $\R$ (vu comme droite affine), les fermés de Zariski sont les ensembles finis et $\R$ lui-même. Ainsi, la topologie de Zariski sur $\R$ coïncide avec la topologie cofinie.

  En dimension supérieure, sur $k^n$ avec $n \geq 2$, la topologie de Zariski ne coïncide plus avec la topologie cofinie. Par exemple, dans $\R^2$, l'ensemble $\{(x,y) : y = 0\}$ est un fermé de Zariski (zéros du polynôme $y$) qui n'est pas fini.
\end{exemple}

% ============================================================
\section{Comparaison de topologies}
\label{sec:comparaison-topologies}

\begin{definition}[Topologies comparables]\label{def:topo-comparables}
  \index{topologie!plus fine}\index{topologie!moins fine}
  Soient $\tau_1$ et $\tau_2$ deux topologies sur un même ensemble $X$. On dit que $\tau_1$ est \textbf{moins fine} (ou \textbf{plus grossière}) que $\tau_2$, et on note $\tau_1 \subset \tau_2$, si tout ouvert de $\tau_1$ est aussi un ouvert de $\tau_2$. On dit alors aussi que $\tau_2$ est \textbf{plus fine} que $\tau_1$.

  Lorsque $\tau_1 \subset \tau_2$ ou $\tau_2 \subset \tau_1$, on dit que $\tau_1$ et $\tau_2$ sont \textbf{comparables}.
\end{definition}

\begin{remarque}
  La relation $\subset$ définit un ordre partiel sur l'ensemble de toutes les topologies sur~$X$. La topologie indiscrète $\tau_i$ est le plus petit élément, et la topologie discrète $\tau_d$ est le plus grand élément. On a donc, pour toute topologie $\tau$ sur $X$ :
  \[
    \tau_i \subset \tau \subset \tau_d.
  \]
\end{remarque}

\begin{exemple}\label{ex:comparaison-R}
  Sur $\R$, on a la chaîne d'inclusions strictes suivante :
  \[
    \tau_i \subsetneq \tau_{\mathrm{cof}} \subsetneq \tau_{\mathrm{cod}} \subsetneq \tau_{\mathrm{us}} \subsetneq \tau_d,
  \]
  où $\tau_{\mathrm{us}}$ désigne la topologie usuelle (métrique euclidienne).
\end{exemple}

\begin{proposition}\label{prop:intersection-topologies}
  \index{topologie!intersection de topologies}
  Soit $({\tau_j})_{j \in J}$ une famille de topologies sur un ensemble $X$. Alors
  \[
    \tau = \bigcap_{j \in J} \tau_j
  \]
  est une topologie sur $X$.
\end{proposition}

\begin{proof}
  Vérifions les trois axiomes.

  \ref{ax:O1} : Pour tout $j \in J$, $\varnothing \in \tau_j$ et $X \in \tau_j$, donc $\varnothing \in \tau$ et $X \in \tau$.

  \ref{ax:O2} : Soit $(U_i)_{i \in I}$ une famille d'éléments de $\tau$. Pour chaque $j \in J$ et chaque $i \in I$, on a $U_i \in \tau_j$. Comme $\tau_j$ est une topologie, $\bigcup_{i \in I} U_i \in \tau_j$. Ceci valant pour tout $j$, on a $\bigcup_{i \in I} U_i \in \tau$.

  \ref{ax:O3} : Le raisonnement est identique pour les intersections finies.
\end{proof}

\begin{remarque}
  En revanche, la réunion de deux topologies n'est en général pas une topologie (voir exercice~\ref{exo:reunion-topologies}).
\end{remarque}

La proposition~\ref{prop:intersection-topologies} montre que l'ensemble des topologies sur $X$, ordonné par inclusion, est un treillis complet. L'infimum d'une famille de topologies est leur intersection. Le supremum est la topologie la moins fine contenant leur réunion.

% ============================================================
\section{Voisinages}
\label{sec:voisinages}

\begin{definition}[Voisinage]\label{def:voisinage}
  \index{voisinage}
  Soit $(X,\tau)$ un espace topologique et $x \in X$. On appelle \textbf{voisinage} de $x$ toute partie $V \subset X$ telle qu'il existe un ouvert $U \in \tau$ vérifiant $x \in U \subset V$.

  L'ensemble de tous les voisinages de $x$ est noté $\Vois{x}$ et s'appelle le \textbf{système de voisinages}\index{système de voisinages} de $x$.
\end{definition}

\begin{proposition}\label{prop:proprietes-voisinages}
  Soit $(X,\tau)$ un espace topologique et $x \in X$. Le système de voisinages $\Vois{x}$ possède les propriétés suivantes :
  \begin{enumerate}[label=\textbf{(V\arabic*)}]
    \item\label{ax:V1} $\Vois{x} \neq \varnothing$, et pour tout $V \in \Vois{x}$, $x \in V$.
    \item\label{ax:V2} Si $V \in \Vois{x}$ et $V \subset W$, alors $W \in \Vois{x}$.
    \item\label{ax:V3} Si $V_1, V_2 \in \Vois{x}$, alors $V_1 \cap V_2 \in \Vois{x}$.
    \item\label{ax:V4} Pour tout $V \in \Vois{x}$, il existe $W \in \Vois{x}$ tel que $W \subset V$ et $V \in \Vois{y}$ pour tout $y \in W$.
  \end{enumerate}
\end{proposition}

\begin{proof}
  \ref{ax:V1} : $X$ est un ouvert contenant $x$, donc $X \in \Vois{x}$. De plus, si $V \in \Vois{x}$, il existe un ouvert $U$ avec $x \in U \subset V$, d'où $x \in V$.

  \ref{ax:V2} : Si $V \in \Vois{x}$, il existe $U \in \tau$ avec $x \in U \subset V$. Si $V \subset W$, alors $x \in U \subset W$, donc $W \in \Vois{x}$.

  \ref{ax:V3} : Soient $U_1, U_2 \in \tau$ avec $x \in U_k \subset V_k$ pour $k = 1,2$. Alors $U_1 \cap U_2 \in \tau$ (intersection finie d'ouverts) et $x \in U_1 \cap U_2 \subset V_1 \cap V_2$.

  \ref{ax:V4} : Si $V \in \Vois{x}$, soit $U \in \tau$ avec $x \in U \subset V$. Posons $W = U$. Pour tout $y \in W = U$, comme $U$ est ouvert et $U \subset V$, on a $V \in \Vois{y}$. De plus $W \subset V$ et $W \in \Vois{x}$.
\end{proof}

\begin{remarque}
  La propriété~\ref{ax:V4} est subtile. Elle exprime que tout voisinage contient un voisinage <<\,ouvert\,>>, c'est-à-dire un voisinage de chacun de ses points. Réciproquement, étant donné pour chaque $x \in X$ un filtre $\Vois{x}$ satisfaisant \ref{ax:V1}--\ref{ax:V4}, il existe une unique topologie sur $X$ dont les systèmes de voisinages sont exactement les $\Vois{x}$. Cette approche dite <<\,par les voisinages\,>> est celle adoptée par Bourbaki.
\end{remarque}

\begin{proposition}\label{prop:ouvert-voisinage}
  \index{ouvert!caractérisation par voisinages}
  Soit $(X,\tau)$ un espace topologique. Une partie $U \subset X$ est ouverte si et seulement si $U$ est voisinage de chacun de ses points, c'est-à-dire :
  \[
    U \in \tau \iff \forall x \in U,\; U \in \Vois{x}.
  \]
\end{proposition}

\begin{proof}
  Si $U \in \tau$ et $x \in U$, alors $x \in U \subset U$ avec $U$ ouvert, donc $U \in \Vois{x}$.

  Réciproquement, supposons que pour tout $x \in U$, $U \in \Vois{x}$. Pour chaque $x \in U$, il existe $U_x \in \tau$ avec $x \in U_x \subset U$. Alors $U = \bigcup_{x \in U} U_x$ est une réunion d'ouverts, donc $U \in \tau$.
\end{proof}

% ============================================================
\section{Illustrations}
\label{sec:illustrations-chap1}

\begin{center}
\begin{tikzpicture}[scale=1.0]
  % Venn diagram of open sets
  \draw[thick,rounded corners=15pt,fill=blue!8] (-3.5,-2.5) rectangle (3.5,2.5);
  \node[anchor=north east] at (3.4,2.4) {$X$};

  \draw[thick,fill=green!15,draw=green!60!black] (0,0) ellipse (2.5cm and 1.5cm);
  \node at (2.0,1.1) {$U_1$};

  \draw[thick,fill=orange!15,draw=orange!70!black] (-0.8,0.2) ellipse (1.5cm and 1.0cm);
  \node at (-1.8,0.9) {$U_2$};

  \draw[thick,fill=red!10,draw=red!60!black] (0.5,-0.1) circle (0.7cm);
  \node at (0.5,-0.1) {$U_3$};

  \node[below] at (0,-2.8) {Ouverts emboîtés : $U_3 \subset U_2 \cap U_1 \subset U_1 \subset X$};
\end{tikzpicture}
\end{center}

\begin{center}
\begin{tikzpicture}[scale=0.85,
  topo/.style={draw,thick,rounded corners=3pt,minimum width=2cm,minimum height=1cm,align=center}]
  % Lattice of topologies
  \node[topo,fill=gray!10] (ti) at (0,0) {$\tau_i = \{\varnothing, X\}$};
  \node[topo,fill=yellow!15] (tcof) at (0,2) {$\tau_{\mathrm{cof}}$};
  \node[topo,fill=orange!12] (tcod) at (-2.5,4) {$\tau_{\mathrm{cod}}$};
  \node[topo,fill=green!12] (tus) at (2.5,4) {$\tau_{\mathrm{us}}$};
  \node[topo,fill=blue!10] (td) at (0,6) {$\tau_d = \mathcal{P}(X)$};

  \draw[-{Stealth},thick] (ti) -- (tcof) node[midway,right] {\small plus fine};
  \draw[-{Stealth},thick] (tcof) -- (tcod);
  \draw[-{Stealth},thick] (tcof) -- (tus);
  \draw[-{Stealth},thick] (tcod) -- (td);
  \draw[-{Stealth},thick] (tus) -- (td);

  \node[below] at (0,-0.8) {Treillis des topologies sur $\R$ (diagramme partiel)};
\end{tikzpicture}
\end{center}

% ============================================================
\section{Contre-exemple : intersection infinie d'ouverts}
\label{sec:contre-exemple-intersection}

\begin{exemple}\label{ex:intersection-infinie-ouverts}
  \index{intersection infinie d'ouverts}
  Dans $\R$ muni de la topologie usuelle, considérons la famille d'ouverts
  \[
    U_n = \left(-\frac{1}{n},\, \frac{1}{n}\right), \quad n \geq 1.
  \]
  Chaque $U_n$ est ouvert. Cependant,
  \[
    \bigcap_{n=1}^{\infty} U_n = \{0\},
  \]
  qui est un singleton, donc un fermé non ouvert de $\R$. Cet exemple montre que l'axiome~\ref{ax:O3} ne peut pas être étendu aux intersections infinies.
\end{exemple}

\begin{proof}
  Soit $x \in \bigcap_{n \geq 1} U_n$. Pour tout $n \geq 1$, $|x| < 1/n$, donc $|x| \leq 0$, d'où $x = 0$. Réciproquement, $0 \in U_n$ pour tout $n$. Ainsi $\bigcap_{n \geq 1} U_n = \{0\}$.

  Le singleton $\{0\}$ n'est pas ouvert dans $\R$ : pour tout $r > 0$, l'intervalle $(-r,r)$ n'est pas contenu dans $\{0\}$.
\end{proof}

% ============================================================
\section{Exemples travaillés}
\label{sec:exemples-travailles-chap1}

\begin{exemple}[Topologie sur un ensemble à trois éléments]\label{ex:topo-trois-elements}
  Soit $X = \{a,b,c\}$. Considérons $\tau = \{\varnothing, \{a\}, \{a,b\}, X\}$. Vérifions que $\tau$ est une topologie.

  \ref{ax:O1} : $\varnothing \in \tau$ et $X \in \tau$. \checkmark

  \ref{ax:O2} : Vérifions toutes les réunions possibles :
  $\{a\} \cup \{a,b\} = \{a,b\} \in \tau$,
  $\{a\} \cup X = X \in \tau$,
  $\{a,b\} \cup X = X \in \tau$,
  $\{a\} \cup \{a,b\} \cup X = X \in \tau$. \checkmark

  \ref{ax:O3} : Vérifions les intersections finies :
  $\{a\} \cap \{a,b\} = \{a\} \in \tau$,
  $\{a\} \cap X = \{a\} \in \tau$,
  $\{a,b\} \cap X = \{a,b\} \in \tau$. \checkmark

  Ainsi $\tau$ est bien une topologie. On constate que $\{b\}$ n'est pas ouvert et que $\{a,c\}$ n'est pas ouvert.
\end{exemple}

\begin{exemple}[Voisinages dans la topologie cofinie]\label{ex:voisinages-cofinie}
  Soit $X$ un ensemble infini muni de la topologie cofinie $\tau_{\mathrm{cof}}$. Soit $x \in X$. Les voisinages de $x$ sont les parties $V$ de $X$ telles qu'il existe un ouvert $U$ avec $x \in U \subset V$, c'est-à-dire les parties $V$ contenant un ouvert non vide passant par $x$.

  Un ouvert non vide $U$ de $\tau_{\mathrm{cof}}$ est un ensemble dont le complémentaire est fini. Ainsi, $V \in \Vois{x}$ si et seulement si $V$ contient une partie $U$ de complémentaire fini avec $x \in U$. En particulier, $V$ est voisinage de $x$ si et seulement si $X \setminus V$ est fini ou, de façon équivalente, $V$ est cofiniment contenu dans $X$.

  Plus précisément, $V \in \Vois{x}$ si et seulement si $x \in V$ et $X \setminus V$ est fini. Cela revient à dire que chaque voisinage de $x$ ne diffère de $X$ que par un nombre fini de points.
\end{exemple}

\begin{exemple}[Deux distances, même topologie]\label{ex:distances-meme-topo}
  Sur $\R^n$, considérons les trois distances suivantes :
  \[
    d_1(x,y) = \sum_{i=1}^n |x_i - y_i|, \quad
    d_2(x,y) = \sqrt{\sum_{i=1}^n (x_i - y_i)^2}, \quad
    d_\infty(x,y) = \max_{1 \leq i \leq n} |x_i - y_i|.
  \]
  Ces trois distances sont équivalentes au sens où il existe des constantes $c, C > 0$ telles que
  \[
    c \, d_\infty(x,y) \leq d_2(x,y) \leq d_1(x,y) \leq C \, d_\infty(x,y)
  \]
  pour tous $x, y \in \R^n$. Plus précisément, $d_\infty \leq d_2 \leq d_1 \leq n \, d_\infty$ et $d_2 \leq \sqrt{n}\, d_\infty$.

  En conséquence, les trois distances engendrent la même topologie sur $\R^n$ : toute boule ouverte pour l'une contient une boule ouverte pour l'autre, et réciproquement.
\end{exemple}

% ============================================================
\section{Exercices}
\label{sec:exercices-chap1}

\begin{exercice}[$\star$]\label{exo:verif-topologie}
  Soit $X = \{1,2,3,4\}$. Parmi les familles suivantes, lesquelles sont des topologies sur $X$ ?
  \begin{enumerate}[label=(\alph*)]
    \item $\tau_1 = \{\varnothing, \{1\}, \{1,2\}, \{1,2,3\}, X\}$
    \item $\tau_2 = \{\varnothing, \{1\}, \{2,3\}, \{1,2,3\}, X\}$
    \item $\tau_3 = \{\varnothing, \{1,2\}, \{3,4\}, X\}$
    \item $\tau_4 = \{\varnothing, \{1\}, \{2\}, \{1,2\}, X\}$
  \end{enumerate}
\end{exercice}

\begin{exercice}[$\star$]\label{exo:cofinie-T1}
  \index{séparation!$T_1$}
  Montrer que la topologie cofinie sur un ensemble $X$ est $T_1$, c'est-à-dire que pour tous $x \neq y$ dans $X$, il existe un ouvert contenant $x$ mais pas $y$.
\end{exercice}

\begin{exercice}[$\star$]\label{exo:cofinie-non-hausdorff}
  Montrer que si $X$ est infini, la topologie cofinie sur $X$ n'est pas séparée (non $T_2$, non Hausdorff) : deux ouverts non vides quelconques ont une intersection non vide.
\end{exercice}

\begin{exercice}[$\star\star$]\label{exo:reunion-topologies}
  Donner un exemple de deux topologies $\tau_1$ et $\tau_2$ sur un ensemble $X$ telles que $\tau_1 \cup \tau_2$ ne soit pas une topologie.
\end{exercice}

\begin{exercice}[$\star\star$]\label{exo:nombre-topologies}
  Déterminer le nombre de topologies distinctes sur un ensemble à $2$ éléments, puis sur un ensemble à $3$ éléments.
\end{exercice}

\begin{exercice}[$\star\star$]\label{exo:topo-engendree}
  \index{topologie engendrée}
  Soit $X$ un ensemble et $\mathcal{S} \subset \mathcal{P}(X)$. On définit la \emph{topologie engendrée par $\mathcal{S}$} comme la plus petite topologie contenant $\mathcal{S}$. Montrer que cette topologie existe et la décrire explicitement à l'aide d'intersections finies et de réunions quelconques d'éléments de $\mathcal{S}$.
\end{exercice}

\begin{exercice}[$\star\star$]\label{exo:codenombrable-non-metrisable}
  Montrer que la topologie codénombrable sur $\R$ n'est pas métrisable. \emph{Indication : montrer que dans cette topologie, la seule suite convergente est la suite stationnaire.}
\end{exercice}

\begin{exercice}[$\star\star\star$]\label{exo:sup-topologies}
  Soient $\tau_1$ et $\tau_2$ deux topologies sur $X$. On note $\tau_1 \vee \tau_2$ la borne supérieure de $\tau_1$ et $\tau_2$ dans le treillis des topologies sur $X$. Montrer que $\tau_1 \vee \tau_2$ est la topologie engendrée par $\tau_1 \cup \tau_2$ et décrire ses ouverts.
\end{exercice}

\begin{exercice}[$\star\star\star$]\label{exo:topologie-particuliere}
  Soit $p$ un nombre premier. Sur $\Z$, on définit $U \subset \Z$ comme ouvert si pour tout $n \in U$, il existe $k \geq 0$ tel que $n + p^k \Z \subset U$ (c'est-à-dire la classe de $n$ modulo $p^k$ est contenue dans $U$). Montrer que cela définit une topologie sur $\Z$ (topologie $p$-adique). Décrire les ouverts de cette topologie.
\end{exercice}

% ============================================================
\section*{Résumé du chapitre}
\addcontentsline{toc}{section}{Résumé du chapitre}

\begin{itemize}
  \item Un \textbf{espace topologique} $(X,\tau)$ est un ensemble $X$ muni d'une famille $\tau$ de parties (les ouverts) stable par réunion quelconque et intersection finie, et contenant $\varnothing$ et $X$.
  \item Exemples fondamentaux : topologie discrète, indiscrète, cofinie, codénombrable, métrique, de Zariski.
  \item Les topologies sur $X$ forment un \textbf{treillis complet} pour l'inclusion, avec la topologie indiscrète comme minimum et la discrète comme maximum.
  \item L'intersection d'une famille de topologies est une topologie ; la réunion ne l'est pas en général.
  \item Un \textbf{voisinage} de $x$ est une partie contenant un ouvert passant par $x$. Le système de voisinages satisfait les axiomes \ref{ax:V1}--\ref{ax:V4} et caractérise la topologie.
  \item L'axiome d'intersection finie ne s'étend pas aux intersections infinies.
\end{itemize}

% ============================================================
%  CHAPITRE 2
% ============================================================
\chapter{Ouverts, fermés, intérieur, adhérence, frontière}
\label{chap:ouverts-fermes}

\section*{Introduction}
\addcontentsline{toc}{section}{Introduction}

Dans un espace métrique $(X,d)$, un ensemble $F$ est fermé lorsque son complémentaire $X \setminus F$ est ouvert, ou de manière équivalente, lorsque $F$ contient toutes les limites de suites d'éléments de $F$. L'intérieur de $A$ est le plus grand ouvert contenu dans $A$, l'adhérence de $A$ le plus petit fermé contenant $A$, et la frontière est la <<\,pellicule\,>> qui les sépare.

En topologie générale, ces notions se définissent de manière purement ensembliste, sans recours à une distance ni aux suites. Ce chapitre développe systématiquement ces concepts et établit leurs propriétés fondamentales. Nous verrons en particulier que la caractérisation séquentielle de l'adhérence, valable dans les espaces métriques, échoue dans le cadre général.

% ============================================================
\section{Ensembles fermés}
\label{sec:fermes}

\begin{definition}[Ensemble fermé]\label{def:ferme}
  \index{fermé}
  Soit $(X,\tau)$ un espace topologique. Une partie $F \subset X$ est dite \textbf{fermée} si son complémentaire $X \setminus F$ est ouvert, c'est-à-dire $X \setminus F \in \tau$.
\end{definition}

\begin{proposition}\label{prop:proprietes-fermes}
  Soit $(X,\tau)$ un espace topologique. La famille $\mathcal{F}$ des fermés de $X$ vérifie :
  \begin{enumerate}[label=\textbf{(F\arabic*)}]
    \item\label{ax:F1} $\varnothing \in \mathcal{F}$ et $X \in \mathcal{F}$.
    \item\label{ax:F2} Toute intersection (même infinie) d'éléments de $\mathcal{F}$ appartient à $\mathcal{F}$.
    \item\label{ax:F3} Toute réunion finie d'éléments de $\mathcal{F}$ appartient à $\mathcal{F}$.
  \end{enumerate}
\end{proposition}

\begin{proof}
  \ref{ax:F1} : $X \setminus \varnothing = X \in \tau$ et $X \setminus X = \varnothing \in \tau$, donc $\varnothing$ et $X$ sont fermés.

  \ref{ax:F2} : Soit $(F_i)_{i \in I}$ une famille de fermés. Par les lois de De Morgan,
  \[
    X \setminus \bigcap_{i \in I} F_i = \bigcup_{i \in I} (X \setminus F_i).
  \]
  Chaque $X \setminus F_i$ est ouvert, donc leur réunion est ouverte par~\ref{ax:O2}. Ainsi $\bigcap_{i \in I} F_i$ est fermé.

  \ref{ax:F3} : De même, si $F_1, \ldots, F_n$ sont fermés,
  \[
    X \setminus \bigcup_{k=1}^n F_k = \bigcap_{k=1}^n (X \setminus F_k)
  \]
  est une intersection finie d'ouverts, donc un ouvert par~\ref{ax:O3}.
\end{proof}

\begin{remarque}
  La dualité ouverts/fermés par passage au complémentaire permet de définir une topologie de manière équivalente par la donnée des fermés satisfaisant \ref{ax:F1}--\ref{ax:F3}. Les rôles des réunions et intersections sont échangés : les fermés sont stables par intersection quelconque mais seulement par réunion finie.
\end{remarque}

\begin{exemple}\label{ex:fermes-cofinie}
  Dans la topologie cofinie sur un ensemble infini $X$, les fermés sont $X$ et les parties finies de $X$. En effet, $F$ est fermé si et seulement si $X \setminus F$ est ouvert, c'est-à-dire $X \setminus F = \varnothing$ (donc $F = X$) ou $X \setminus (X \setminus F) = F$ est fini.
\end{exemple}

\begin{exemple}\label{ex:ouverts-et-fermes}
  \index{ensemble!ouvert et fermé}
  Dans tout espace topologique $(X,\tau)$, les ensembles $\varnothing$ et $X$ sont à la fois ouverts et fermés. Un ensemble qui est simultanément ouvert et fermé est parfois dit \textbf{clopen}\index{clopen} (de l'anglais \emph{closed-open}). Dans la topologie discrète, toute partie est clopen.
\end{exemple}

% ============================================================
\section{Intérieur}
\label{sec:interieur}

\begin{definition}[Intérieur]\label{def:interieur}
  \index{intérieur}
  Soit $(X,\tau)$ un espace topologique et $A \subset X$. L'\textbf{intérieur} de $A$, noté $\Int{A}$ ou $A^\circ$, est la réunion de tous les ouverts contenus dans $A$ :
  \[
    \Int{A} = \bigcup_{\substack{U \in \tau \\ U \subset A}} U.
  \]
\end{definition}

\begin{theoreme}\label{thm:proprietes-interieur}
  Soit $(X,\tau)$ un espace topologique. Pour toutes parties $A, B \subset X$ :
  \begin{enumerate}[label=(\roman*)]
    \item\label{it:int-ouvert} $\Int{A}$ est ouvert (c'est le plus grand ouvert contenu dans $A$).
    \item\label{it:int-inclus} $\Int{A} \subset A$.
    \item\label{it:int-carac} $\Int{A} = A$ si et seulement si $A$ est ouvert.
    \item\label{it:int-idempotent} $\Int{(\Int{A})} = \Int{A}$ (idempotence).
    \item\label{it:int-monotone} Si $A \subset B$, alors $\Int{A} \subset \Int{B}$ (monotonie).
    \item\label{it:int-intersection} $\Int{(A \cap B)} = \Int{A} \cap \Int{B}$.
    \item\label{it:int-voisinage} $x \in \Int{A}$ si et seulement si $A \in \Vois{x}$.
  \end{enumerate}
\end{theoreme}

\begin{proof}
  \ref{it:int-ouvert} : $\Int{A}$ est une réunion d'ouverts, donc un ouvert par~\ref{ax:O2}. Par construction, tout ouvert $U \subset A$ vérifie $U \subset \Int{A}$, donc $\Int{A}$ est le plus grand ouvert contenu dans $A$.

  \ref{it:int-inclus} : Chaque ouvert $U$ dans la réunion satisfait $U \subset A$, donc $\Int{A} \subset A$.

  \ref{it:int-carac} : Si $A$ est ouvert, alors $A$ apparaît dans la réunion, donc $A \subset \Int{A}$, et avec \ref{it:int-inclus}, $\Int{A} = A$. Réciproquement, si $\Int{A} = A$, alors $A = \Int{A}$ est ouvert par~\ref{it:int-ouvert}.

  \ref{it:int-idempotent} : $\Int{A}$ est ouvert par \ref{it:int-ouvert}, donc $\Int{(\Int{A})} = \Int{A}$ par~\ref{it:int-carac}.

  \ref{it:int-monotone} : Si $A \subset B$, tout ouvert contenu dans $A$ est contenu dans $B$, donc $\Int{A} \subset \Int{B}$.

  \ref{it:int-intersection} : Comme $A \cap B \subset A$ et $A \cap B \subset B$, on a $\Int{(A \cap B)} \subset \Int{A}$ et $\Int{(A \cap B)} \subset \Int{B}$, d'où $\Int{(A \cap B)} \subset \Int{A} \cap \Int{B}$. Réciproquement, $\Int{A} \cap \Int{B}$ est ouvert (intersection finie d'ouverts) et $\Int{A} \cap \Int{B} \subset A \cap B$, donc $\Int{A} \cap \Int{B} \subset \Int{(A \cap B)}$.

  \ref{it:int-voisinage} : $x \in \Int{A}$ si et seulement s'il existe un ouvert $U$ avec $x \in U \subset A$, ce qui est exactement la définition de $A \in \Vois{x}$.
\end{proof}

\begin{remarque}\label{rem:int-reunion}
  En général, $\Int{(A \cup B)} \neq \Int{A} \cup \Int{B}$. On a seulement l'inclusion $\Int{A} \cup \Int{B} \subset \Int{(A \cup B)}$. Par exemple, dans $\R$ usuel, si $A = \Q$ et $B = \R \setminus \Q$, alors $\Int{A} = \Int{B} = \varnothing$, mais $\Int{(A \cup B)} = \Int{\R} = \R$.
\end{remarque}

% ============================================================
\section{Adhérence}
\label{sec:adherence}

\begin{definition}[Adhérence]\label{def:adherence}
  \index{adhérence}
  Soit $(X,\tau)$ un espace topologique et $A \subset X$. L'\textbf{adhérence} (ou \textbf{fermeture}) de $A$, notée $\Cl{A}$, est l'intersection de tous les fermés contenant $A$ :
  \[
    \Cl{A} = \bigcap_{\substack{F \text{ fermé} \\ A \subset F}} F.
  \]
\end{definition}

\begin{theoreme}\label{thm:proprietes-adherence}
  Soit $(X,\tau)$ un espace topologique. Pour toutes parties $A, B \subset X$ :
  \begin{enumerate}[label=(\roman*)]
    \item\label{it:cl-ferme} $\Cl{A}$ est fermé (c'est le plus petit fermé contenant $A$).
    \item\label{it:cl-contient} $A \subset \Cl{A}$.
    \item\label{it:cl-carac} $\Cl{A} = A$ si et seulement si $A$ est fermé.
    \item\label{it:cl-idempotent} $\Cl{\Cl{A}} = \Cl{A}$ (idempotence).
    \item\label{it:cl-monotone} Si $A \subset B$, alors $\Cl{A} \subset \Cl{B}$ (monotonie).
    \item\label{it:cl-reunion} $\Cl{A \cup B} = \Cl{A} \cup \Cl{B}$.
    \item\label{it:cl-voisinage} $x \in \Cl{A}$ si et seulement si tout voisinage de $x$ rencontre $A$ :
    \[
      x \in \Cl{A} \iff \forall V \in \Vois{x},\; V \cap A \neq \varnothing.
    \]
  \end{enumerate}
\end{theoreme}

\begin{proof}
  \ref{it:cl-ferme} : $\Cl{A}$ est une intersection de fermés, donc fermée par~\ref{ax:F2}. Tout fermé $F$ contenant $A$ vérifie $\Cl{A} \subset F$, donc $\Cl{A}$ est le plus petit fermé contenant $A$.

  \ref{it:cl-contient} : $A$ est contenu dans chaque fermé $F \supset A$, donc $A \subset \Cl{A}$.

  \ref{it:cl-carac} : Si $A$ est fermé, $A$ apparaît dans l'intersection, donc $\Cl{A} \subset A$, et avec \ref{it:cl-contient}, $\Cl{A} = A$. Réciproquement, si $\Cl{A} = A$, alors $A$ est fermé par~\ref{it:cl-ferme}.

  \ref{it:cl-idempotent} : $\Cl{A}$ est fermé, donc $\Cl{\Cl{A}} = \Cl{A}$ par~\ref{it:cl-carac}.

  \ref{it:cl-monotone} : Si $A \subset B$, tout fermé contenant $B$ contient $A$, donc $\Cl{A} \subset \Cl{B}$.

  \ref{it:cl-reunion} : Comme $A \subset A \cup B$ et $B \subset A \cup B$, on a $\Cl{A} \subset \Cl{A \cup B}$ et $\Cl{B} \subset \Cl{A \cup B}$, d'où $\Cl{A} \cup \Cl{B} \subset \Cl{A \cup B}$. Réciproquement, $\Cl{A} \cup \Cl{B}$ est fermé (réunion finie de fermés) et contient $A \cup B$, donc $\Cl{A \cup B} \subset \Cl{A} \cup \Cl{B}$.

  \ref{it:cl-voisinage} : Supposons $x \notin \Cl{A}$. Alors $U = X \setminus \Cl{A}$ est un ouvert contenant $x$, et $U \cap A = \varnothing$ car $A \subset \Cl{A}$. Donc $U \in \Vois{x}$ et $U \cap A = \varnothing$. Réciproquement, s'il existe $V \in \Vois{x}$ avec $V \cap A = \varnothing$, soit $U \in \tau$ avec $x \in U \subset V$. Alors $U \cap A = \varnothing$, donc $A \subset X \setminus U$. Comme $X \setminus U$ est fermé, $\Cl{A} \subset X \setminus U$, d'où $x \notin \Cl{A}$.
\end{proof}

\begin{proposition}[Dualité intérieur-adhérence]\label{prop:dualite-int-cl}
  \index{dualité intérieur-adhérence}
  Pour toute partie $A$ d'un espace topologique $(X,\tau)$ :
  \[
    X \setminus \Int{A} = \Cl{X \setminus A} \qquad \text{et} \qquad X \setminus \Cl{A} = \Int{X \setminus A}.
  \]
\end{proposition}

\begin{proof}
  Montrons la première égalité. On a :
  \begin{align*}
    X \setminus \Int{A}
    &= X \setminus \bigcup_{\substack{U \in \tau \\ U \subset A}} U
    = \bigcap_{\substack{U \in \tau \\ U \subset A}} (X \setminus U)
    = \bigcap_{\substack{F \text{ fermé} \\ X \setminus A \subset F}} F
    = \Cl{X \setminus A}.
  \end{align*}
  La deuxième égalité s'obtient en appliquant la première à $X \setminus A$ puis en prenant le complémentaire.
\end{proof}

\begin{remarque}[Caractérisation séquentielle]\label{rem:carac-sequentielle}
  \index{adhérence!caractérisation séquentielle}
  Dans un espace métrique $(X,d)$, on a $x \in \Cl{A}$ si et seulement s'il existe une suite $(a_n)_{n \geq 0}$ d'éléments de $A$ telle que $a_n \to x$. Cette caractérisation est \textbf{fausse en général} dans un espace topologique non premier-dénombrable. Par exemple, dans un espace non dénombrable $X$ muni de la topologie codénombrable $\tau_{\mathrm{cod}}$, un sous-ensemble dénombrable infini $A$ a pour adhérence $\Cl{A} = X$, mais les seules suites convergentes sont les suites stationnaires (voir exercice~\ref{exo:adherence-non-sequentielle}).
\end{remarque}

% ============================================================
\section{Frontière}
\label{sec:frontiere}

\begin{definition}[Frontière]\label{def:frontiere}
  \index{frontière}
  Soit $(X,\tau)$ un espace topologique et $A \subset X$. La \textbf{frontière} de $A$, notée $\Fr{A}$, est définie par
  \[
    \Fr{A} = \Cl{A} \setminus \Int{A} = \Cl{A} \cap \Cl{X \setminus A}.
  \]
\end{definition}

\begin{proposition}\label{prop:proprietes-frontiere}
  Pour toute partie $A$ d'un espace topologique $(X,\tau)$ :
  \begin{enumerate}[label=(\roman*)]
    \item $\Fr{A}$ est fermé.
    \item $\Fr{A} = \Fr{(X \setminus A)}$.
    \item $A$ est ouvert si et seulement si $A \cap \Fr{A} = \varnothing$.
    \item $A$ est fermé si et seulement si $\Fr{A} \subset A$.
    \item $A$ est ouvert et fermé si et seulement si $\Fr{A} = \varnothing$.
    \item $X = \Int{A} \sqcup \Fr{A} \sqcup \Int{(X \setminus A)}$ (partition).
  \end{enumerate}
\end{proposition}

\begin{proof}
  (i) $\Fr{A} = \Cl{A} \cap \Cl{X \setminus A}$ est l'intersection de deux fermés.

  (ii) $\Fr{(X \setminus A)} = \Cl{X \setminus A} \cap \Cl{X \setminus (X \setminus A)} = \Cl{X \setminus A} \cap \Cl{A} = \Fr{A}$.

  (iii) $A$ est ouvert $\Leftrightarrow$ $\Int{A} = A$ $\Leftrightarrow$ $\Fr{A} = \Cl{A} \setminus A \subset \Cl{A} \setminus A$. Si $A$ est ouvert, $\Int{A} = A$, donc $\Fr{A} = \Cl{A} \setminus A$ et $A \cap \Fr{A} = A \cap (\Cl{A} \setminus A) = \varnothing$. Réciproquement, si $A \cap \Fr{A} = \varnothing$, alors $A \cap (\Cl{A} \setminus \Int{A}) = \varnothing$. Comme $A \subset \Cl{A}$, cela donne $A \setminus \Int{A} = \varnothing$, donc $A \subset \Int{A}$, et ainsi $A = \Int{A}$ est ouvert.

  (iv) $A$ est fermé $\Leftrightarrow$ $\Cl{A} = A$ $\Leftrightarrow$ $\Fr{A} = A \setminus \Int{A} \subset A$.

  (v) Conséquence immédiate de (iii) et (iv).

  (vi) Les trois ensembles $\Int{A}$, $\Fr{A}$ et $\Int{(X \setminus A)}$ sont deux à deux disjoints. De plus, par la proposition~\ref{prop:dualite-int-cl}, $X \setminus \Int{A} = \Cl{X \setminus A}$ et $X \setminus \Int{(X \setminus A)} = \Cl{A}$, d'où
  \[
    \Int{A} \cup \Fr{A} \cup \Int{(X \setminus A)} = \Int{A} \cup (\Cl{A} \cap \Cl{X \setminus A}) \cup \Int{(X \setminus A)}.
  \]
  Tout point $x \in X$ est soit dans $\Int{A}$ (voisinage contenu dans $A$), soit dans $\Int{(X \setminus A)}$ (voisinage disjoint de $A$), soit dans aucun des deux, auquel cas tout voisinage de $x$ rencontre à la fois $A$ et $X \setminus A$, donc $x \in \Cl{A} \cap \Cl{X \setminus A} = \Fr{A}$.
\end{proof}

% ============================================================
\section{Ensemble dérivé, points isolés, ensembles denses}
\label{sec:derive-dense}

\begin{definition}[Point d'accumulation et ensemble dérivé]\label{def:derive}
  \index{point d'accumulation}\index{ensemble dérivé}
  Soit $(X,\tau)$ un espace topologique et $A \subset X$. Un point $x \in X$ est dit \textbf{point d'accumulation} (ou \textbf{point limite}) de $A$ si tout voisinage de $x$ contient au moins un point de $A$ distinct de $x$ :
  \[
    \forall V \in \Vois{x},\quad (V \setminus \{x\}) \cap A \neq \varnothing.
  \]
  L'\textbf{ensemble dérivé} de $A$, noté $A'$, est l'ensemble des points d'accumulation de $A$.
\end{definition}

\begin{proposition}\label{prop:adherence-derive}
  Pour toute partie $A$ d'un espace topologique :
  \[
    \Cl{A} = A \cup A'.
  \]
\end{proposition}

\begin{proof}
  Soit $x \in \Cl{A}$. Si $x \in A$, c'est fini. Si $x \notin A$, alors pour tout $V \in \Vois{x}$, $V \cap A \neq \varnothing$ (car $x \in \Cl{A}$), et comme $x \notin A$, on a $(V \setminus \{x\}) \cap A \neq \varnothing$, donc $x \in A'$.

  Réciproquement, si $x \in A$, alors $x \in \Cl{A}$ car $A \subset \Cl{A}$. Si $x \in A'$, alors tout voisinage de $x$ rencontre $A$ (même en un point différent de $x$), donc $x \in \Cl{A}$.
\end{proof}

\begin{definition}[Point isolé]\label{def:point-isole}
  \index{point isolé}
  Un point $x \in A$ est dit \textbf{isolé dans $A$} s'il existe un voisinage $V$ de $x$ tel que $V \cap A = \{x\}$, autrement dit si $x \in A \setminus A'$.
\end{definition}

\begin{definition}[Ensemble dense]\label{def:dense}
  \index{ensemble dense}\index{dense}
  Une partie $A$ d'un espace topologique $(X,\tau)$ est dite \textbf{dense} dans $X$ si $\Cl{A} = X$, c'est-à-dire si tout ouvert non vide de $X$ rencontre $A$ :
  \[
    \forall U \in \tau,\; U \neq \varnothing \implies U \cap A \neq \varnothing.
  \]
\end{definition}

\begin{exemple}\label{ex:Q-dense-R}
  \index{rationnels!denses dans $\R$}
  $\Q$ est dense dans $\R$ pour la topologie usuelle. En effet, pour tout ouvert non vide $U$ de $\R$, $U$ contient un intervalle ouvert $(a,b)$ avec $a < b$, et la densité de $\Q$ dans $\R$ (propriété archimédienne) assure l'existence d'un rationnel dans $(a,b)$.

  De même, $\R \setminus \Q$ (les irrationnels) est dense dans $\R$.
\end{exemple}

\begin{definition}[Ensemble nulle part dense]\label{def:nulle-part-dense}
  \index{nulle part dense}
  Une partie $A$ d'un espace topologique $(X,\tau)$ est dite \textbf{nulle part dense} (ou \textbf{rare}) si l'intérieur de son adhérence est vide :
  \[
    \Int{\Cl{A}} = \varnothing.
  \]
\end{definition}

\begin{exemple}\label{ex:Z-nulle-part-dense}
  $\Z$ est nulle part dense dans $\R$ pour la topologie usuelle. En effet, $\Cl{\Z} = \Z$ (car $\Z$ est fermé), et $\Int{\Z} = \varnothing$ (car $\Z$ ne contient aucun intervalle ouvert).

  De même, l'ensemble de Cantor est nulle part dense dans $[0,1]$.
\end{exemple}

% ============================================================
\section{Axiomes de fermeture de Kuratowski}
\label{sec:kuratowski-axiomes}

\index{Kuratowski!axiomes de fermeture}

Les propriétés de l'adhérence établies au théorème~\ref{thm:proprietes-adherence} permettent de caractériser axiomatiquement une topologie par son opérateur d'adhérence.

\begin{theoreme}[Axiomes de Kuratowski]\label{thm:kuratowski-axiomes}
  Soit $X$ un ensemble et $\varphi : \mathcal{P}(X) \to \mathcal{P}(X)$ une application satisfaisant les quatre axiomes suivants :
  \begin{enumerate}[label=\textbf{(K\arabic*)}]
    \item\label{ax:K1} $\varphi(\varnothing) = \varnothing$.
    \item\label{ax:K2} $A \subset \varphi(A)$ pour tout $A \subset X$.
    \item\label{ax:K3} $\varphi(A \cup B) = \varphi(A) \cup \varphi(B)$ pour tous $A, B \subset X$.
    \item\label{ax:K4} $\varphi(\varphi(A)) = \varphi(A)$ pour tout $A \subset X$ (idempotence).
  \end{enumerate}
  Alors il existe une unique topologie $\tau$ sur $X$ telle que $\varphi(A) = \Cl{A}$ pour tout $A \subset X$. Cette topologie est définie par
  \[
    \tau = \{U \subset X : \varphi(X \setminus U) = X \setminus U\}.
  \]
\end{theoreme}

\begin{proof}
  Posons $\mathcal{F} = \{F \subset X : \varphi(F) = F\}$. Montrons que $\mathcal{F}$ satisfait les axiomes des fermés.

  \ref{ax:F1} : $\varphi(\varnothing) = \varnothing$ par \ref{ax:K1}, donc $\varnothing \in \mathcal{F}$. Par \ref{ax:K2}, $X \subset \varphi(X) \subset X$, donc $\varphi(X) = X$ et $X \in \mathcal{F}$.

  \ref{ax:F2} : Soit $(F_i)_{i \in I}$ une famille d'éléments de $\mathcal{F}$. Posons $F = \bigcap_{i \in I} F_i$. Pour tout $i$, $F \subset F_i$, donc par \ref{ax:K2} et \ref{ax:K3}, et la monotonie (qui découle de \ref{ax:K2} et \ref{ax:K3} : si $A \subset B$, alors $\varphi(B) = \varphi(A \cup B) = \varphi(A) \cup \varphi(B) \supset \varphi(A)$), on a $\varphi(F) \subset \varphi(F_i) = F_i$ pour tout $i$. Donc $\varphi(F) \subset \bigcap_{i \in I} F_i = F$. Avec \ref{ax:K2}, $\varphi(F) = F$.

  \ref{ax:F3} : Si $F_1, F_2 \in \mathcal{F}$, alors $\varphi(F_1 \cup F_2) = \varphi(F_1) \cup \varphi(F_2) = F_1 \cup F_2$ par \ref{ax:K3}, donc $F_1 \cup F_2 \in \mathcal{F}$. Par récurrence, toute réunion finie d'éléments de $\mathcal{F}$ est dans $\mathcal{F}$.

  Ainsi $\tau = \{X \setminus F : F \in \mathcal{F}\}$ est une topologie. Montrons que $\varphi(A) = \Cl{A}$ pour tout $A$. On a $\varphi(A) \in \mathcal{F}$ par \ref{ax:K4}, et $A \subset \varphi(A)$ par \ref{ax:K2}, donc $\varphi(A)$ est un fermé contenant $A$. Si $F \in \mathcal{F}$ avec $A \subset F$, alors $\varphi(A) \subset \varphi(F) = F$ par monotonie. Donc $\varphi(A)$ est le plus petit fermé contenant $A$, c'est-à-dire $\Cl{A}$.

  L'unicité découle du fait que la topologie est déterminée par ses fermés, qui sont eux-mêmes déterminés par $\varphi$.
\end{proof}

% ============================================================
\section{Théorème des 14 ensembles de Kuratowski}
\label{sec:kuratowski-14}

\index{Kuratowski!théorème des 14 ensembles}

\begin{theoreme}[Kuratowski, 1922]\label{thm:kuratowski-14}
  Soit $(X,\tau)$ un espace topologique et $A \subset X$. En appliquant itérativement les opérations de fermeture $A \mapsto \Cl{A}$ et de complémentation $A \mapsto X \setminus A$ à partir de $A$, on obtient au plus $14$ ensembles distincts. De plus, cette borne est atteinte.
\end{theoreme}

\begin{proof}[Démonstration (esquisse)]
  Notons $c(A) = X \setminus A$ la complémentation et $k(A) = \Cl{A}$ la fermeture. On remarque que :
  \begin{itemize}
    \item $c$ est une involution : $c \circ c = \mathrm{id}$.
    \item $k$ est idempotente : $k \circ k = k$.
    \item L'intérieur s'écrit $i = c \circ k \circ c$ (par la proposition~\ref{prop:dualite-int-cl}).
  \end{itemize}

  Les $14$ ensembles possibles à partir de $A$ sont :
  \begin{align*}
    &A, \quad k(A), \quad ck(A), \quad kck(A), \quad ckck(A), \quad kckck(A), \quad ckckck(A), \\
    &c(A), \quad kc(A), \quad ckc(A), \quad kckcA), \quad ckckcA), \quad kckckcA), \quad ckckckc(A).
  \end{align*}
  Réécrivons plus lisiblement. À partir de $A$, on peut appliquer :
  \[
    A,\; kA,\; ckA,\; kckA,\; ckckA,\; kckckA,\; ckckckA
  \]
  et, en commençant par $c$ :
  \[
    cA,\; kcA,\; ckcA,\; kckcA,\; ckckcA,\; kckckcA,\; ckckckcA.
  \]

  La clé est de montrer que $kckckck = kck$, c'est-à-dire que toute chaîne d'opérations se stabilise après au plus $7$ étapes. Cela repose sur l'identité fondamentale :
  \[
    k \circ c \circ k \circ c \circ k \circ c \circ k = k \circ c \circ k.
  \]

  Pour établir cette identité, posons $i = ckc$ (l'opérateur intérieur). On doit montrer $kik = k$. Comme $ikA \subset kA$ (l'intérieur est contenu dans l'ensemble) et $kA$ est fermé, on a $kikA \subset kA$. Par ailleurs, $ikA \subset kA$ donne (par monotonie de $k$ et idempotence) $kikA \supset kA$ --- plus précisément, soit $U = ikA = \Int{\Cl{A}}$. Alors $\Cl{U} \subset \Cl{\Cl{A}} = \Cl{A}$, d'où $k(ikA) \subset kA$. Pour l'autre sens, $A \subset kA$ donne $ikA \supset iA$, puis $kikA \supset kiA \supset kA$ n'est pas vrai en général. Il faut un argument plus fin.

  Montrons directement que la chaîne $k, ck, kck, ckck, kckck, ckckck$ est la plus longue non stabilisée. L'identité clé est : pour toute partie $B$ fermée ($kB = B$), on a $kckck(B) = kck(B)$, car $ck(B) \supset ckck(B)$ (puisque $kck(B) \supset k(B) = B$, d'où $ck(B) \supset ckck(B)$ est faux dans ce sens). Prouvons autrement :

  Puisque $\Cl{A} \supset A$, on a $\Int{\Cl{A}} \supset \Int{A}$, d'où $\Cl{\Int{\Cl{A}}} \supset \Cl{\Int{A}}$, puis $\Int{\Cl{\Int{\Cl{A}}}} \supset \Int{\Cl{\Int{A}}}$. D'autre part, posons $B = \Cl{A}$. Alors $\Int{B}$ est ouvert, donc $\Cl{\Int{B}}$ est fermé, donc $\Int{\Cl{\Int{B}}}$ est ouvert. On a $\Int{B} \subset \Cl{\Int{B}}$, d'où $\Int{B} \subset \Int{\Cl{\Int{B}}}$. Mais aussi $\Int{B} \subset B$, donc $\Cl{\Int{B}} \subset \Cl{B} = B$ (car $B$ est fermé), donc $\Int{\Cl{\Int{B}}} \subset \Int{B}$. Finalement $\Int{\Cl{\Int{B}}} = \Int{B}$, c'est-à-dire $ickA = ikA$, soit $kckckA = kckA$ après application de $k$ et $c$ de manière appropriée.

  Plus précisément, l'identité $\Int{\Cl{\Int{\Cl{A}}}} = \Int{\Cl{A}}$ pour tout $A$ se démontre ainsi :
  \begin{itemize}
    \item $\Int{\Cl{A}}$ est ouvert, donc $\Cl{\Int{\Cl{A}}} \supset \Int{\Cl{A}}$, d'où $\Int{\Cl{\Int{\Cl{A}}}} \supset \Int{\Cl{A}}$ (car $\Int{\Cl{A}}$ est ouvert et contenu dans $\Cl{\Int{\Cl{A}}}$).
    \item $\Int{\Cl{A}} \subset \Cl{A}$, donc $\Cl{\Int{\Cl{A}}} \subset \Cl{\Cl{A}} = \Cl{A}$, d'où $\Int{\Cl{\Int{\Cl{A}}}} \subset \Int{\Cl{A}}$.
  \end{itemize}
  Ainsi $\Int{\Cl{\Int{\Cl{A}}}} = \Int{\Cl{A}}$, ce qui donne $kckckck = kck$ après complémentation, et la chaîne se stabilise.

  La borne $14$ est atteinte dans $\R$ avec la topologie usuelle pour l'ensemble
  \[
    A = (0,1) \cup (1,2) \cup \{3\} \cup ([4,5] \cap \Q).
  \]
\end{proof}

% ============================================================
\section{Illustrations}
\label{sec:illustrations-chap2}

\begin{center}
\begin{tikzpicture}[scale=1.2]
  % Interior, closure, boundary visualization
  % The set A
  \fill[blue!10] plot[smooth cycle,tension=0.7] coordinates {(-2,0) (-1.5,1.5) (0.5,1.8) (2,1) (2.2,-0.5) (1,-1.5) (-1.5,-1.2)};

  % Interior (slightly smaller)
  \fill[green!25] plot[smooth cycle,tension=0.7] coordinates {(-1.5,0.1) (-1.1,1.1) (0.3,1.3) (1.5,0.7) (1.7,-0.3) (0.7,-1) (-1.2,-0.8)};

  % Boundary of A
  \draw[red,ultra thick] plot[smooth cycle,tension=0.7] coordinates {(-2,0) (-1.5,1.5) (0.5,1.8) (2,1) (2.2,-0.5) (1,-1.5) (-1.5,-1.2)};

  % Labels
  \node at (0,0) {$\Int{A}$};
  \node at (-1.8,1.2) {$A$};
  \node[red] at (2.5,0.3) {$\Fr{A}$};

  % Closure extends slightly beyond boundary
  \draw[blue,thick,dashed] plot[smooth cycle,tension=0.7] coordinates {(-2.2,0.1) (-1.7,1.7) (0.6,2.0) (2.2,1.2) (2.4,-0.6) (1.1,-1.7) (-1.7,-1.4)};
  \node[blue] at (-2.5,1.5) {$\Cl{A}$};

  \node[below] at (0,-2.2) {Intérieur, adhérence et frontière d'un ensemble $A$};
\end{tikzpicture}
\end{center}

\begin{center}
\begin{tikzpicture}[scale=1.0]
  % Number line illustration for Q dense in R
  \draw[thick,->] (-4,0) -- (4,0) node[right] {$\R$};

  % Some rational points
  \foreach \x/\lab in {-3/{-3}, -2/{-2}, -1/{-1}, 0/{0}, 1/{1}, 2/{2}, 3/{3}} {
    \fill[blue] (\x,0) circle (1.5pt);
    \node[below,font=\small] at (\x,-0.15) {$\lab$};
  }
  % Some more rationals
  \foreach \x in {-2.5,-1.5,-0.5,0.5,1.5,2.5,0.333,-0.333,0.667,-0.667,1.333,-1.333} {
    \fill[blue] (\x,0) circle (1pt);
  }

  % Indicate density with dots
  \node[above,blue,font=\small] at (0,0.3) {$\Q$ dense dans $\R$};
  \node[below,font=\small] at (0,-0.6) {Tout intervalle ouvert contient un rationnel};
\end{tikzpicture}
\end{center}

% ============================================================
\section{Contre-exemples}
\label{sec:contre-exemples-chap2}

\begin{exemple}[Adhérence et suites : le cas non premier-dénombrable]\label{ex:adherence-non-sequentielle}
  \index{adhérence!non séquentielle}
  Soit $X$ un ensemble non dénombrable muni de la topologie codénombrable $\tau_{\mathrm{cod}}$. Soit $A \subset X$ un sous-ensemble dénombrable infini.

  \emph{Adhérence de $A$.} Soit $x \in X$ et $U \in \tau_{\mathrm{cod}}$ avec $x \in U$. Alors $X \setminus U$ est dénombrable, donc $U$ est non dénombrable (car $X$ l'est). En particulier, $A \not\subset X \setminus U$ (car $A$ est infini et $U$ manque seulement un ensemble dénombrable, mais il faut vérifier que $U$ rencontre $A$). Plus précisément, $A \cap U = A \setminus (X \setminus U)$ : comme $A$ est infini et $X \setminus U$ est dénombrable, si $A \subset X \setminus U$ cela est possible. Prenons plutôt $A$ infini avec $\Cl{A} = X$ directement : le seul fermé contenant $A$ et distinct de $X$ aurait son complémentaire ouvert non vide, donc de complémentaire dénombrable, mais ce complémentaire contiendrait $A$. Ainsi, si $F$ est fermé avec $A \subset F$ et $F \neq X$, alors $X \setminus F \neq \varnothing$ est ouvert, donc $X \setminus (X \setminus F) = F$ est dénombrable. Comme $A \subset F$ et $A$ est infini, cela est possible. Corrigeons : les fermés de $\tau_{\mathrm{cod}}$ sont $X$ et les ensembles dénombrables. Comme $A$ est dénombrable, $A$ lui-même est fermé! Donc $\Cl{A} = A \neq X$.

  Prenons plutôt $A$ \emph{non dénombrable} avec $X \setminus A$ non dénombrable. Alors $A$ n'est ni ouvert (car $X \setminus A$ n'est pas dénombrable) ni fermé (car $A$ n'est pas dénombrable et $A \neq X$). L'adhérence de $A$ est $\Cl{A} = X$ (le seul fermé non dénombrable est $X$). Or, les suites convergentes pour $\tau_{\mathrm{cod}}$ sont les suites stationnaires : si $(x_n)$ converge vers $\ell$ et n'est pas stationnaire, le complémentaire $\{x_n : x_n \neq \ell\}$ est dénombrable, donc son complémentaire est cocountable, donc ouvert, et contient $\ell$ mais pas les $x_n \neq \ell$, contradiction. Ainsi aucun point de $X \setminus A$ n'est limite d'une suite d'éléments de $A$.
\end{exemple}

\begin{exemple}[La droite de Sorgenfrey]\label{ex:sorgenfrey}
  \index{droite de Sorgenfrey}\index{Sorgenfrey}
  La \textbf{droite de Sorgenfrey} est l'ensemble $\R$ muni de la topologie $\tau_S$ engendrée par les intervalles semi-ouverts à droite $[a,b)$ avec $a < b$. Vérifions que ces intervalles forment une base de topologie (ce fait sera formalisé au chapitre suivant, mais l'idée est claire) :
  \begin{itemize}
    \item $\R = \bigcup_{n \in \Z} [n, n+1)$.
    \item $[a,b) \cap [c,d) = [\max(a,c), \min(b,d))$ (qui est vide si $\max(a,c) \geq \min(b,d)$, ou un intervalle semi-ouvert sinon).
  \end{itemize}

  \emph{Propriétés.}
  \begin{enumerate}[label=(\alph*)]
    \item $\tau_S$ est strictement plus fine que la topologie usuelle $\tau_{\mathrm{us}}$. En effet, tout ouvert $(a,b)$ de $\tau_{\mathrm{us}}$ s'écrit $\bigcup_{n \geq 1} [a + 1/n, b)$, donc $(a,b) \in \tau_S$. Mais $[0,1) \in \tau_S \setminus \tau_{\mathrm{us}}$.
    \item L'ensemble $\Q$ est dense dans $(\R, \tau_S)$. En effet, tout intervalle $[a,b)$ non vide contient un rationnel (par densité de $\Q$ dans $\R$ usuel).
    \item Les fermés de $\tau_S$ incluent les intervalles de la forme $(-\infty, a)$ et $(b, +\infty]$ : le complémentaire de $(-\infty, a)$ est $[a, +\infty) = \bigcup_{n \geq 0} [a+n, a+n+1)$, qui est ouvert dans $\tau_S$.
    \item L'adhérence de $(0,1)$ dans $\tau_S$ est $[0,1)$. En effet, $[0,1)$ est fermé dans $\tau_S$ (car $\R \setminus [0,1) = (-\infty, 0) \cup [1, +\infty)$ est ouvert). De plus, pour tout $\varepsilon > 0$, $[0, \varepsilon) \cap (0,1) \neq \varnothing$, donc $0 \in \Cl{(0,1)}$.
  \end{enumerate}
\end{exemple}

% ============================================================
\section{Exemples travaillés}
\label{sec:exemples-travailles-chap2}

\begin{exemple}[Intérieur et adhérence de $\Q$ dans $\R$]\label{ex:int-cl-Q}
  Dans $(\R, \tau_{\mathrm{us}})$ :
  \begin{itemize}
    \item $\Int{\Q} = \varnothing$. En effet, tout intervalle ouvert $(a,b)$ contient un irrationnel, donc $(a,b) \not\subset \Q$. Aucun ouvert non vide n'est contenu dans $\Q$.
    \item $\Cl{\Q} = \R$. Tout intervalle ouvert $(a,b)$ non vide contient un rationnel, donc tout ouvert non vide rencontre $\Q$.
    \item $\Fr{\Q} = \Cl{\Q} \setminus \Int{\Q} = \R \setminus \varnothing = \R$.
  \end{itemize}
  De même, $\Int{(\R \setminus \Q)} = \varnothing$, $\Cl{\R \setminus \Q} = \R$, et $\Fr{(\R \setminus \Q)} = \R$.
\end{exemple}

\begin{exemple}[Intérieur et adhérence dans la topologie cofinie]\label{ex:int-cl-cofinie}
  Soit $X = \R$ muni de la topologie cofinie $\tau_{\mathrm{cof}}$. Soit $A = [0,1]$.
  \begin{itemize}
    \item $\Int{A}$ : un ouvert non vide de $\tau_{\mathrm{cof}}$ a un complémentaire fini, donc contient tous les réels sauf un nombre fini. Ainsi, aucun ouvert non vide n'est contenu dans $[0,1]$ (car $[0,1]^c$ est infini). Donc $\Int{A} = \varnothing$.
    \item $\Cl{A}$ : les fermés de $\tau_{\mathrm{cof}}$ sont $\R$ et les ensembles finis. Comme $A = [0,1]$ est infini, le seul fermé le contenant est $\R$. Donc $\Cl{A} = \R$.
    \item $\Fr{A} = \R \setminus \varnothing = \R$.
  \end{itemize}
  Plus généralement, dans $(\R, \tau_{\mathrm{cof}})$, toute partie infinie a un intérieur vide et une adhérence égale à $\R$.
\end{exemple}

\begin{exemple}[Les 14 ensembles de Kuratowski]\label{ex:kuratowski-14-explicite}
  \index{Kuratowski!exemple des 14 ensembles}
  Prenons $X = \R$ avec la topologie usuelle et
  \[
    A = (0,1) \cup (1,2) \cup \{3\} \cup ([4,5] \cap \Q).
  \]
  Calculons les $14$ ensembles. Notons $k = \Cl{\cdot}$ et $c = X \setminus (\cdot)$.

  \begin{enumerate}
    \item $A = (0,1) \cup (1,2) \cup \{3\} \cup ([4,5] \cap \Q)$.
    \item $kA = [0,2] \cup \{3\} \cup [4,5]$.
    \item $ckA = (-\infty,0) \cup (2,3) \cup (3,4) \cup (5,+\infty)$.
    \item $kckA = (-\infty,0] \cup [2,4] \cup [5,+\infty)$.
    \item $ckckA = (0,2) \cup (4,5)$.
    \item $kckckA = [0,2] \cup [4,5]$.
    \item $ckckckA = (-\infty,0) \cup (2,4) \cup (5,+\infty)$.
    \item $cA = (-\infty,0] \cup \{1\} \cup [2,3) \cup (3,4) \cup ([4,5] \setminus \Q) \cup (5,+\infty)$.
    \item $kcA = (-\infty,0] \cup \{1\} \cup [2,+\infty)$.
    \item $ckcA = (0,1) \cup (1,2)$.
    \item $kckcA = [0,2]$.
    \item $ckckcA = (-\infty,0) \cup (2,+\infty)$.
    \item $kckckcA = (-\infty,0] \cup [2,+\infty)$.
    \item $ckckckcA = (0,2)$.
  \end{enumerate}
  On vérifie que ces $14$ ensembles sont tous distincts, et que toute application supplémentaire de $k$ ou $c$ redonne un ensemble déjà obtenu.
\end{exemple}

% ============================================================
\section{Exercices}
\label{sec:exercices-chap2}

\begin{exercice}[$\star$]\label{exo:fermes-usuelle}
  Déterminer si les ensembles suivants sont ouverts, fermés, les deux, ou ni l'un ni l'autre, dans $\R$ muni de la topologie usuelle :
  \begin{enumerate}[label=(\alph*)]
    \item $[0,1)$
    \item $\Z$
    \item $\{1/n : n \geq 1\}$
    \item $\{1/n : n \geq 1\} \cup \{0\}$
    \item $(0,1) \cup (2,3)$
  \end{enumerate}
\end{exercice}

\begin{exercice}[$\star$]\label{exo:int-cl-calculs}
  Calculer l'intérieur, l'adhérence et la frontière des ensembles suivants dans $(\R, \tau_{\mathrm{us}})$ :
  \begin{enumerate}[label=(\alph*)]
    \item $A = (0,1]$
    \item $B = \Z$
    \item $C = \{1/n : n \geq 1\}$
    \item $D = [0,1] \cap \Q$
  \end{enumerate}
\end{exercice}

\begin{exercice}[$\star$]\label{exo:dense-R}
  Montrer que $\R \setminus \Q$ est dense dans $(\R, \tau_{\mathrm{us}})$.
\end{exercice}

\begin{exercice}[$\star\star$]\label{exo:frontiere-frontiere}
  Soit $A$ une partie d'un espace topologique $(X,\tau)$. Montrer que $\Fr{(\Fr{A})} \subset \Fr{A}$. Donner un exemple où l'inclusion est stricte.
\end{exercice}

\begin{exercice}[$\star\star$]\label{exo:int-cl-formules}
  Soit $(X,\tau)$ un espace topologique et $A, B \subset X$. Montrer :
  \begin{enumerate}[label=(\alph*)]
    \item $\Cl{A \cap B} \subset \Cl{A} \cap \Cl{B}$ (l'égalité est-elle vraie en général ?)
    \item $\Int{A} \cup \Int{B} \subset \Int{(A \cup B)}$ (l'égalité est-elle vraie en général ?)
  \end{enumerate}
  Donner des contre-exemples pour les égalités.
\end{exercice}

\begin{exercice}[$\star\star$]\label{exo:derive-proprietes}
  Soit $(X,\tau)$ un espace topologique et $A, B \subset X$. Montrer :
  \begin{enumerate}[label=(\alph*)]
    \item $(A \cup B)' = A' \cup B'$.
    \item $A' \subset \Cl{A}$.
    \item Si $A \subset B$, alors $A' \subset B'$.
  \end{enumerate}
\end{exercice}

\begin{exercice}[$\star\star$]\label{exo:adherence-non-sequentielle}
  \index{topologie codénombrable!suites convergentes}
  Soit $X$ un ensemble non dénombrable muni de la topologie codénombrable $\tau_{\mathrm{cod}}$.
  \begin{enumerate}[label=(\alph*)]
    \item Montrer que les seules suites convergentes dans $(X, \tau_{\mathrm{cod}})$ sont les suites stationnaires à partir d'un certain rang.
    \item En déduire qu'il existe une partie $A$ de $X$ et un point $x \in \Cl{A}$ tels qu'aucune suite d'éléments de $A$ ne converge vers $x$.
  \end{enumerate}
\end{exercice}

\begin{exercice}[$\star\star$]\label{exo:sorgenfrey-calculs}
  Dans la droite de Sorgenfrey $(\R, \tau_S)$ :
  \begin{enumerate}[label=(\alph*)]
    \item Calculer $\Int{(0,1]}$, $\Cl{(0,1)}$, $\Fr{[0,1)}$.
    \item Montrer que $\Q$ est dense.
    \item Montrer que $[a,b)$ est à la fois ouvert et fermé.
  \end{enumerate}
\end{exercice}

\begin{exercice}[$\star\star\star$]\label{exo:nulle-part-dense}
  \begin{enumerate}[label=(\alph*)]
    \item Montrer que $A$ est nulle part dense si et seulement si pour tout ouvert non vide $U$, il existe un ouvert non vide $V \subset U$ tel que $V \cap A = \varnothing$.
    \item Montrer qu'une réunion finie d'ensembles nulle part denses est nulle part dense.
    \item Montrer qu'une réunion dénombrable d'ensembles nulle part denses n'est pas nécessairement nulle part dense. \emph{Indication : considérer $\Q = \bigcup_{q \in \Q} \{q\}$.}
  \end{enumerate}
\end{exercice}

\begin{exercice}[$\star\star\star$]\label{exo:kuratowski-verification}
  Vérifier que l'ensemble $A = (0,1) \cup (1,2) \cup \{3\} \cup ([4,5] \cap \Q)$ de $\R$ engendre bien $14$ ensembles distincts par les opérations de fermeture et de complémentation.
\end{exercice}

% ============================================================
\section*{Résumé du chapitre}
\addcontentsline{toc}{section}{Résumé du chapitre}

\begin{itemize}
  \item Un ensemble est \textbf{fermé} si son complémentaire est ouvert. Les fermés sont stables par intersection quelconque et réunion finie.
  \item L'\textbf{intérieur} $\Int{A}$ est le plus grand ouvert contenu dans $A$ ; l'\textbf{adhérence} $\Cl{A}$ est le plus petit fermé contenant $A$. Ces opérateurs sont liés par dualité : $\Cl{A} = X \setminus \Int{(X \setminus A)}$.
  \item La \textbf{frontière} $\Fr{A} = \Cl{A} \setminus \Int{A}$ est un fermé. Elle est vide si et seulement si $A$ est clopen.
  \item L'\textbf{ensemble dérivé} $A'$ est l'ensemble des points d'accumulation. On a $\Cl{A} = A \cup A'$.
  \item $A$ est \textbf{dense} si $\Cl{A} = X$ ; $A$ est \textbf{nulle part dense} si $\Int{\Cl{A}} = \varnothing$.
  \item Les \textbf{axiomes de Kuratowski} \ref{ax:K1}--\ref{ax:K4} donnent une axiomatisation alternative de la topologie via l'opérateur d'adhérence.
  \item Le \textbf{théorème des 14 ensembles} de Kuratowski affirme que les opérations de fermeture et complémentation engendrent au plus $14$ ensembles distincts à partir d'une partie donnée.
  \item La caractérisation séquentielle de l'adhérence, valable dans les espaces métriques, \textbf{échoue} dans les espaces non premier-dénombrables (par exemple, la topologie codénombrable).
  \item La \textbf{droite de Sorgenfrey} fournit un riche réservoir de contre-exemples ; elle est plus fine que la topologie usuelle.
\end{itemize}
% chapters3_4.tex — Chapitres 3 et 4 : Bases et sous-bases ; Continuité et homéomorphismes
% Cours de Topologie Générale — Partie 2
% Convention : les environnements definition, theoreme, proposition, lemme,
% corollaire, remarque, exemple, exercice sont définis dans le préambule.

% ============================================================================
%
%                         CHAPITRE 3
%
% ============================================================================

\chapter{Bases et sous-bases, topologies engendrées}\label{ch:bases-sousbases}

\section*{Introduction}
\addcontentsline{toc}{section}{Introduction}

Dans un espace métrique $(X,d)$, la topologie est entièrement déterminée par
les boules ouvertes~: un sous-ensemble $U\subseteq X$ est ouvert si et
seulement si, pour tout $x\in U$, il existe $r>0$ tel que $B(x,r)\subseteq U$.
Autrement dit, la collection des boules ouvertes \og{}engendre\fg{} la
topologie. Cette observation conduit à une question naturelle~: dans un espace
topologique quelconque, peut-on trouver une collection plus maniable d'ouverts
à partir de laquelle on reconstruit toute la topologie~?

\medskip

La réponse est affirmative et fait l'objet de ce chapitre. Nous introduisons
les notions de \emph{base} et de \emph{sous-base} d'une topologie, puis nous
étudions les topologies \emph{engendrées} par une collection arbitraire de
sous-ensembles. Ces outils sont fondamentaux~: ils permettent de construire
efficacement de nouvelles topologies et de vérifier des propriétés en ne
travaillant que sur une collection restreinte d'ouverts. Nous examinons ensuite
les axiomes de dénombrabilité --- espaces à base dénombrable, premier
dénombrable, séparables --- et leurs relations.

\begin{center}
\begin{tikzpicture}[scale=1.1]
  % Big open set
  \draw[thick, blue!60!black, fill=blue!5, rounded corners=18pt]
    (0,0) -- (5,0.3) -- (5.5,2) -- (4.8,4) -- (2,4.5) -- (-0.3,3.5) -- cycle;
  \node[blue!60!black] at (2.6,4.8) {$U$ (ouvert)};
  % Base elements inside
  \draw[thick, red!70!black, fill=red!8] (1.2,1.5) ellipse (0.8 and 0.6);
  \node[red!70!black, font=\small] at (1.2,1.5) {$B_1$};
  \draw[thick, red!70!black, fill=red!8] (3.0,2.8) ellipse (1.0 and 0.7);
  \node[red!70!black, font=\small] at (3.0,2.8) {$B_2$};
  \draw[thick, red!70!black, fill=red!8] (4.0,1.2) ellipse (0.7 and 0.55);
  \node[red!70!black, font=\small] at (4.0,1.2) {$B_3$};
  \draw[thick, red!70!black, fill=red!8] (2.0,3.5) ellipse (0.6 and 0.45);
  \node[red!70!black, font=\small] at (2.0,3.5) {$B_4$};
  \draw[thick, red!70!black, fill=red!8] (3.8,3.2) ellipse (0.55 and 0.4);
  \node[red!70!black, font=\small] at (3.8,3.2) {$B_5$};
  % Points
  \foreach \x/\y in {0.8/2.2, 1.5/0.9, 2.5/1.8, 3.5/2.0, 4.3/1.6, 1.8/2.9, 4.2/3.5} {
    \fill[black] (\x,\y) circle (1pt);
  }
  % Caption
  \node[below, font=\small\itshape] at (2.7,-0.5)
    {Tout ouvert $U$ est réunion d'éléments de base $B_i$.};
\end{tikzpicture}
\captionof{figure}{Éléments de base engendrant un ouvert.}
\label{fig:base-elements}
\end{center}

% ============================================================
\section{Bases d'une topologie}\label{sec:bases}
% ============================================================

\subsection{Définition et premières propriétés}\label{subsec:base-def}

\begin{definition}[Base d'une topologie]\label{def:base}
  Soit $(X,\T)$ un espace topologique. Une collection
  $\B\subseteq\T$\index{base d'une topologie} est une \emph{base} de la
  topologie $\T$ si tout ouvert non vide $U\in\T$ est réunion d'éléments
  de~$\B$~:
  \[
    \forall\, U\in\T,\quad
    U = \bigcup_{\substack{B\in\B \\ B\subseteq U}} B.
  \]
  De manière équivalente, $\B$ est une base de $\T$ si, pour tout $U\in\T$ et
  tout $x\in U$, il existe $B\in\B$ tel que $x\in B\subseteq U$.
\end{definition}

\begin{exemple}
  Dans l'espace métrique $(\R^n,d_2)$, la collection des boules ouvertes
  $\{B(x,r) : x\in\R^n,\, r>0\}$ est une base de la topologie usuelle. On
  peut même se restreindre aux boules de centres rationnels et de rayons
  rationnels~: $\{B(q,r) : q\in\Q^n,\, r\in\Q_{>0}\}$, qui forme une base
  \emph{dénombrable}.
\end{exemple}

\begin{proposition}[Caractérisation d'une base]\label{prop:carac-base}
  Soit $(X,\T)$ un espace topologique et $\B\subseteq\T$. Les conditions
  suivantes sont équivalentes~:
  \begin{enumerate}[label=\textup{(\roman*)}]
    \item $\B$ est une base de $\T$.
    \item Pour tout $U\in\T$ et tout $x\in U$, il existe $B\in\B$ tel que
      $x\in B\subseteq U$.
    \item $\T = \bigl\{\bigcup_{i\in I} B_i : B_i\in\B,\, I \text{ ensemble
      d'indices}\bigr\} \cup \{\varnothing\}$.
  \end{enumerate}
\end{proposition}

\begin{proof}
  $\mathrm{(i)\Rightarrow(ii)}$~: Si $U=\bigcup_{j\in J} B_j$ avec
  $B_j\in\B$, alors pour tout $x\in U$, il existe $j\in J$ avec $x\in B_j$ et
  $B_j\subseteq U$.

  $\mathrm{(ii)\Rightarrow(iii)}$~: Pour tout $U\in\T$, on a
  $U = \bigcup_{x\in U} B_x$ où $B_x\in\B$ satisfait $x\in B_x\subseteq U$.
  Donc $U$ est réunion d'éléments de $\B$.

  $\mathrm{(iii)\Rightarrow(i)}$~: Immédiat par définition.
\end{proof}

\subsection{Conditions nécessaires et suffisantes}\label{subsec:base-cns}

Le théorème suivant caractérise les collections de sous-ensembles qui sont des
bases d'une topologie, sans référence préalable à une topologie.

\begin{theoreme}[Critère de base]\label{thm:critere-base}
  Soit $X$ un ensemble et $\B$ une collection de sous-ensembles de $X$. Alors
  $\B$ est la base d'une topologie sur $X$ si et seulement si les deux
  conditions suivantes sont satisfaites~:
  \begin{enumerate}[label=\textup{(B\arabic*)}]
    \item\label{B1} $\displaystyle X = \bigcup_{B\in\B} B$.
    \item\label{B2} Pour tous $B_1,B_2\in\B$ et tout $x\in B_1\cap B_2$, il
      existe $B_3\in\B$ tel que $x\in B_3\subseteq B_1\cap B_2$.
  \end{enumerate}
  Dans ce cas, la topologie engendrée est
  \[
    \T_{\B} = \Bigl\{\bigcup_{i\in I} B_i : B_i\in\B,\;
    I\text{ ensemble d'indices}\Bigr\}.
  \]
\end{theoreme}

\begin{proof}
  \textbf{Nécessité.} Supposons que $\B$ est une base de $\T$. Puisque
  $X\in\T$, on a $X=\bigcup_{B\in\B,\, B\subseteq X} B$, d'où \ref{B1}. Pour
  \ref{B2}, soient $B_1,B_2\in\B\subseteq\T$ et $x\in B_1\cap B_2$. Comme
  $B_1\cap B_2\in\T$ (intersection finie d'ouverts), il existe $B_3\in\B$ tel
  que $x\in B_3\subseteq B_1\cap B_2$.

  \textbf{Suffisance.} Posons
  $\T_{\B}=\bigl\{\bigcup_{i\in I} B_i : B_i\in\B\bigr\}\cup\{\varnothing\}$.
  Vérifions que $\T_{\B}$ est une topologie.
  \begin{itemize}
    \item $\varnothing\in\T_{\B}$ par convention (réunion vide) et
      $X\in\T_{\B}$ par \ref{B1}.
    \item $\T_{\B}$ est stable par réunion quelconque~: une réunion de réunions
      d'éléments de $\B$ est encore une réunion d'éléments de $\B$.
    \item Montrons la stabilité par intersection finie. Il suffit de traiter le
      cas de deux éléments. Soient
      $U=\bigcup_{i\in I} B_i$ et $V=\bigcup_{j\in J} B_j$. Alors
      $U\cap V = \bigcup_{(i,j)\in I\times J} (B_i\cap B_j)$. Pour tout
      $x\in B_i\cap B_j$, la condition \ref{B2} fournit $B_x\in\B$ avec
      $x\in B_x\subseteq B_i\cap B_j$. Donc
      $B_i\cap B_j = \bigcup_{x\in B_i\cap B_j} B_x$, et
      $U\cap V\in\T_{\B}$.
  \end{itemize}
  Par construction, $\B$ est une base de $\T_{\B}$.
\end{proof}

\begin{remarque}
  La condition \ref{B2} exprime que l'intersection de deux éléments de base,
  bien qu'elle ne soit pas nécessairement dans $\B$, peut être \og{}raffinée\fg{}
  par des éléments de $\B$. C'est l'analogue topologique de la propriété de
  filtre-base.
\end{remarque}

% ============================================================
\section{Sous-bases et topologies engendrées}\label{sec:sousbases}
% ============================================================

\subsection{Sous-base}\label{subsec:sousbase-def}

\begin{definition}[Sous-base d'une topologie]\label{def:sousbase}
  Soit $(X,\T)$ un espace topologique. Une collection
  $\S\subseteq\T$\index{sous-base} est une \emph{sous-base} de~$\T$ si la
  collection des intersections finies d'éléments de~$\S$~:
  \[
    \B_{\S} = \Bigl\{ S_1\cap S_2\cap\cdots\cap S_n :
    n\in\N^*,\; S_1,\ldots,S_n\in\S \Bigr\} \cup \{X\},
  \]
  est une base de~$\T$.
\end{definition}

\begin{theoreme}[Topologie engendrée par une sous-base]\label{thm:topo-sousbase}
  Soit $X$ un ensemble et $\S$ une collection de sous-ensembles de $X$ telle
  que $\bigcup_{S\in\S} S = X$. Alors la collection $\B_{\S}$ des
  intersections finies d'éléments de $\S$ satisfait les conditions
  \ref{B1}--\ref{B2} du 
ef{thm:critere-base} et engendre donc une
  topologie $\T(\S)$ sur $X$, appelée \emph{topologie engendrée par
  $\S$}\index{topologie engendrée}.
\end{theoreme}

\begin{proof}
  \ref{B1} est satisfaite par hypothèse. Pour \ref{B2}, soient
  $B_1 = \bigcap_{i=1}^{m} S_i$ et $B_2 = \bigcap_{j=1}^{n} T_j$ deux
  éléments de $\B_{\S}$. Alors
  $B_1\cap B_2 = \bigl(\bigcap_{i=1}^{m} S_i\bigr) \cap
  \bigl(\bigcap_{j=1}^{n} T_j\bigr)$ est une intersection finie d'éléments
  de $\S$, donc $B_1\cap B_2\in\B_{\S}$. Il suffit alors de prendre
  $B_3 = B_1\cap B_2$.
\end{proof}

\subsection{Topologie engendrée comme intersection}\label{subsec:topo-engendree}

\begin{theoreme}[Topologie engendrée par une collection]\label{thm:topo-engendree}
  Soit $X$ un ensemble et $\mathcal{C}$ une collection de sous-ensembles de
  $X$. L'intersection
  \[
    \T(\mathcal{C}) = \bigcap \bigl\{ \T : \T \text{ est une topologie sur }
    X \text{ et } \mathcal{C}\subseteq\T \bigr\}
  \]
  est une topologie sur $X$, appelée \emph{topologie engendrée par
  $\mathcal{C}$}. C'est la plus petite topologie (au sens de l'inclusion)
  contenant $\mathcal{C}$.
\end{theoreme}

\begin{proof}
  L'intersection est bien définie car la topologie discrète
  $\mathcal{P}(X)$ contient $\mathcal{C}$. Vérifions les axiomes.
  \begin{itemize}
    \item $\varnothing$ et $X$ appartiennent à toute topologie contenant
      $\mathcal{C}$, donc à l'intersection.
    \item Si $(U_i)_{i\in I}$ est une famille d'éléments de
      $\T(\mathcal{C})$, alors pour toute topologie $\T\supseteq\mathcal{C}$,
      chaque $U_i\in\T$, donc $\bigcup_{i\in I} U_i\in\T$. Ceci étant vrai
      pour toute telle $\T$, on a $\bigcup_{i\in I} U_i\in\T(\mathcal{C})$.
    \item Le raisonnement est identique pour les intersections finies.
  \end{itemize}
  Si $\T'$ est une topologie contenant $\mathcal{C}$, alors $\T'$ fait
  partie de la famille dont on prend l'intersection, donc
  $\T(\mathcal{C})\subseteq\T'$.
\end{proof}

\begin{proposition}[Lien entre les deux constructions]\label{prop:lien-sousbase-engendree}
  Si $\S$ est une collection de sous-ensembles de $X$ avec
  $\bigcup_{S\in\S}S=X$, alors la topologie engendrée par $\S$ au sens du
  
ef{thm:topo-sousbase} coïncide avec la topologie engendrée par $\S$ au
  sens du 
ef{thm:topo-engendree}.
\end{proposition}

\begin{proof}
  Notons $\T_1$ la topologie du 
ef{thm:topo-sousbase} et $\T_2$ celle du
  
ef{thm:topo-engendree}. Puisque $\S\subseteq\T_1$ et que $\T_2$ est la
  plus petite topologie contenant $\S$, on a $\T_2\subseteq\T_1$. Réciproquement,
  tout élément de $\B_{\S}$ (intersection finie d'éléments de $\S$) appartient
  à $\T_2$, et toute réunion d'éléments de $\B_{\S}$ appartient donc à $\T_2$.
  Ainsi $\T_1\subseteq\T_2$.
\end{proof}

% ============================================================
\section{Exemples de bases}\label{sec:exemples-bases}
% ============================================================

\begin{exemple}\label{ex:base-R-usuelle}
  \textbf{Base de la topologie usuelle de $\R$.}
  La collection $\B=\{(a,b) : a,b\in\R,\; a<b\}$ des intervalles ouverts
  bornés est une base de la topologie usuelle. On peut la réduire à
  $\B_{\Q}=\{(a,b) : a,b\in\Q,\; a<b\}$, qui est dénombrable.
\end{exemple}

\begin{exemple}\label{ex:base-sorgenfrey}
  \textbf{Droite de Sorgenfrey\index{droite de Sorgenfrey}.}
  Sur $\R$, considérons la collection
  $\B_\ell = \{[a,b) : a,b\in\R,\; a<b\}$.
  Vérifions \ref{B1}--\ref{B2} du 
ef{thm:critere-base}~:
  \begin{itemize}
    \item \ref{B1}~: Pour tout $x\in\R$, on a $x\in[x,x+1)\in\B_\ell$.
    \item \ref{B2}~: Soient $[a_1,b_1)$ et $[a_2,b_2)$ dans $\B_\ell$.
      Alors $[a_1,b_1)\cap[a_2,b_2) = [\max(a_1,a_2),\min(b_1,b_2))$, qui est
      soit vide, soit un élément de $\B_\ell$.
  \end{itemize}
  La topologie engendrée $\T_\ell$ est la \emph{topologie de Sorgenfrey}. Elle
  est strictement plus fine que la topologie usuelle~: tout intervalle ouvert
  $(a,b)$ est réunion $\bigcup_{n\geq 1}[a+1/n,b)$, mais $[0,1)$ n'est pas
  ouvert pour la topologie usuelle.
\end{exemple}

\begin{exemple}\label{ex:base-ordre}
  \textbf{Topologie de l'ordre\index{topologie de l'ordre}.}
  Soit $(X,\leq)$ un ensemble totalement ordonné. La collection des
  intervalles ouverts $\{(a,b) : a<b\}$, des rayons $\{[a_0,b) : b\in X\}$
  (si $X$ admet un plus petit élément $a_0$) et $\{(a,b_0] : a\in X\}$ (si
  $X$ admet un plus grand élément $b_0$) forme une base de la
  \emph{topologie de l'ordre} sur $X$.
\end{exemple}

% ============================================================
\section{Axiomes de dénombrabilité}\label{sec:axiomes-denombrabilite}
% ============================================================

\subsection{Espaces premier et second dénombrable}\label{subsec:premier-second}

\begin{definition}[Espace second dénombrable]\label{def:second-denombrable}
  Un espace topologique $(X,\T)$ est dit \emph{à base dénombrable} ou
  \emph{second dénombrable}\index{second dénombrable} (on dit aussi qu'il
  satisfait le \emph{deuxième axiome de dénombrabilité}) s'il existe une base
  $\B$ de $\T$ qui est dénombrable.
\end{definition}

\begin{definition}[Espace premier dénombrable]\label{def:premier-denombrable}
  Un espace topologique $(X,\T)$ est dit \emph{premier
  dénombrable}\index{premier dénombrable} (ou satisfait le \emph{premier axiome
  de dénombrabilité}) si tout point $x\in X$ possède une base de voisinages
  dénombrable~: il existe une suite $(V_n)_{n\in\N}$ de voisinages de $x$
  telle que, pour tout voisinage $V$ de $x$, il existe $n\in\N$ avec
  $V_n\subseteq V$.
\end{definition}

\begin{proposition}[Second dénombrable implique premier dénombrable]\label{prop:A2-implique-A1}
  Tout espace second dénombrable est premier dénombrable.
\end{proposition}

\begin{proof}
  Soit $\B=\{B_n : n\in\N\}$ une base dénombrable de $\T$. Pour $x\in X$,
  posons $\B_x = \{B_n\in\B : x\in B_n\}$. C'est une partie dénombrable de
  $\B$. Pour tout voisinage $V$ de $x$, il existe un ouvert $U$ avec
  $x\in U\subseteq V$, puis $B_n\in\B$ avec $x\in B_n\subseteq U\subseteq V$.
  Donc $B_n\in\B_x$ et $\B_x$ est une base de voisinages dénombrable de $x$.
\end{proof}

\subsection{Espaces séparables}\label{subsec:separables}

\begin{definition}[Espace séparable]\label{def:separable}
  Un espace topologique $(X,\T)$ est dit
  \emph{séparable}\index{espace séparable} s'il contient un sous-ensemble
  dense dénombrable~: il existe $D\subseteq X$ dénombrable tel que
  $\overline{D}=X$.
\end{definition}

\begin{proposition}[Second dénombrable implique séparable]\label{prop:A2-implique-separable}
  Tout espace second dénombrable est séparable.
\end{proposition}

\begin{proof}
  Soit $\B=\{B_n : n\in\N\}$ une base dénombrable de $\T$. Pour chaque
  $B_n\neq\varnothing$, choisissons $x_n\in B_n$ (axiome du choix dénombrable).
  L'ensemble $D=\{x_n : n\in\N,\, B_n\neq\varnothing\}$ est dénombrable.
  Montrons que $\overline{D}=X$. Soit $x\in X$ et $U$ un ouvert contenant $x$.
  Il existe $B_n\in\B$ avec $x\in B_n\subseteq U$. Alors $x_n\in B_n\subseteq U$,
  donc $U\cap D\neq\varnothing$. Ainsi $x\in\overline{D}$.
\end{proof}

\subsection{Théorème de Lindelöf}\label{subsec:lindelof}

\begin{definition}[Espace de Lindelöf]\label{def:lindelof}
  Un espace topologique $(X,\T)$ est un \emph{espace de
  Lindelöf}\index{espace de Lindelöf} si tout recouvrement ouvert de $X$
  admet un sous-recouvrement dénombrable.
\end{definition}

\begin{theoreme}[Théorème de Lindelöf]\label{thm:lindelof}
  Tout espace topologique second dénombrable est un espace de Lindelöf.
\end{theoreme}

\begin{proof}
  Soit $\B = \{B_n : n\in\N\}$ une base dénombrable et soit
  $\{U_i\}_{i\in I}$ un recouvrement ouvert de $X$. Pour tout $x\in X$, il
  existe $i_x\in I$ avec $x\in U_{i_x}$, puis $n_x\in\N$ avec
  $x\in B_{n_x}\subseteq U_{i_x}$.

  Posons $J = \{n_x : x\in X\}\subseteq\N$, qui est dénombrable. Pour chaque
  $n\in J$, choisissons $i(n)\in I$ tel que $B_n\subseteq U_{i(n)}$. La
  famille $\{U_{i(n)}\}_{n\in J}$ est dénombrable. Elle recouvre $X$~: pour
  tout $x\in X$, on a $x\in B_{n_x}\subseteq U_{i(n_x)}$.
\end{proof}

\begin{corollaire}\label{cor:cor:Rn-lindelof}
  L'espace $\R^n$ muni de la topologie usuelle est un espace de Lindelöf.
\end{corollaire}

\begin{proof}
  $\R^n$ est second dénombrable (base de boules ouvertes à centres et rayons
  rationnels), donc de Lindelöf par le 
ef{thm:lindelof}.
\end{proof}

% ============================================================
\section{Espaces métriques et dénombrabilité}\label{sec:metriques-denombrabilite}
% ============================================================

\begin{proposition}[Espaces métriques et premier dénombrabilité]\label{prop:metrique-A1}
  Tout espace métrique $(X,d)$ est premier dénombrable.
\end{proposition}

\begin{proof}
  Pour tout $x\in X$, la famille $\{B(x,1/n)\}_{n\geq 1}$ est une base de
  voisinages dénombrable de $x$. En effet, pour tout voisinage $V$ de $x$, il
  existe $\varepsilon>0$ avec $B(x,\varepsilon)\subseteq V$. Choisissons
  $n\in\N^*$ tel que $1/n<\varepsilon$~; alors
  $B(x,1/n)\subseteq B(x,\varepsilon)\subseteq V$.
\end{proof}

\begin{proposition}[Caractérisation de la séparabilité dans les espaces métriques]\label{prop:metrique-separable}
  Pour un espace métrique $(X,d)$, les conditions suivantes sont
  équivalentes~:
  \begin{enumerate}[label=\textup{(\roman*)}]
    \item $X$ est second dénombrable.
    \item $X$ est séparable.
    \item $X$ est de Lindelöf.
  \end{enumerate}
\end{proposition}

\begin{proof}
  Les implications $\mathrm{(i)\Rightarrow(ii)}$ et
  $\mathrm{(i)\Rightarrow(iii)}$ sont déjà établies. Montrons
  $\mathrm{(ii)\Rightarrow(i)}$. Soit $D=\{d_n : n\in\N\}$ un sous-ensemble
  dense dénombrable. La collection
  $\B = \{B(d_n,1/k) : n\in\N,\, k\in\N^*\}$ est dénombrable. Montrons que
  c'est une base. Soit $U$ ouvert et $x\in U$. Il existe $\varepsilon>0$ avec
  $B(x,\varepsilon)\subseteq U$. Puisque $D$ est dense, il existe $d_n\in D$
  avec $d(x,d_n)<\varepsilon/2$. Choisissons $k\in\N^*$ avec
  $1/k<\varepsilon/2$. Alors $x\in B(d_n,1/k)$. De plus, pour tout
  $y\in B(d_n,1/k)$,
  $d(x,y)\leq d(x,d_n)+d(d_n,y) < \varepsilon/2+\varepsilon/2 = \varepsilon$,
  donc $B(d_n,1/k)\subseteq B(x,\varepsilon)\subseteq U$.

  Montrons $\mathrm{(iii)\Rightarrow(ii)}$. Pour chaque $n\geq 1$, le
  recouvrement $\{B(x,1/n) : x\in X\}$ admet un sous-recouvrement
  dénombrable~: il existe $D_n\subseteq X$ dénombrable tel que
  $X=\bigcup_{x\in D_n} B(x,1/n)$. Posons $D=\bigcup_{n\geq 1} D_n$,
  qui est dénombrable. Pour tout $x\in X$ et $\varepsilon>0$, choisissons
  $n$ tel que $1/n<\varepsilon$. Il existe $y\in D_n$ avec $x\in B(y,1/n)$,
  donc $d(x,y)<1/n<\varepsilon$ et $y\in D\cap B(x,\varepsilon)$. Ainsi $D$
  est dense.
\end{proof}

% ============================================================
\section{Contre-exemples}\label{sec:contre-exemples-ch3}
% ============================================================

\begin{proposition}[Propriétés de la droite de Sorgenfrey]\label{prop:sorgenfrey-proprietes}
  La droite de Sorgenfrey $(\R,\T_\ell)$ est~:
  \begin{enumerate}[label=\textup{(\roman*)}]
    \item premier dénombrable,
    \item séparable,
    \item \emph{non} second dénombrable.
  \end{enumerate}
\end{proposition}

\begin{proof}
  \textup{(i)}~: Pour tout $x\in\R$, la famille
  $\{[x,x+1/n) : n\geq 1\}$ est une base de voisinages dénombrable.

  \textup{(ii)}~: Montrons que $\Q$ est dense dans $(\R,\T_\ell)$. Soit
  $[a,b)$ un ouvert de base. L'intervalle $[a,b)$ contient un rationnel
  (densité de $\Q$ dans $\R$ pour la topologie usuelle), donc
  $[a,b)\cap\Q\neq\varnothing$.

  \textup{(iii)}~: Supposons par l'absurde que $\B=\{B_n : n\in\N\}$ soit une
  base dénombrable de $\T_\ell$. Pour tout $x\in\R$, l'ouvert $[x,x+1)$
  contient un $B_{n_x}$ avec $x\in B_{n_x}$. Puisque $B_{n_x}$ est ouvert
  dans $\T_\ell$ et $x\in B_{n_x}$, on a $\inf B_{n_x}=x$ (car $B_{n_x}$
  contient $x$ et est contenu dans $[x,x+1)$, donc ne contient aucun point
  strictement inférieur à $x$). L'application $x\mapsto n_x$ est donc
  injective (si $n_x=n_y$, alors $x=\inf B_{n_x}=\inf B_{n_y}=y$). Mais
  $\R$ est non dénombrable et $\N$ est dénombrable~: contradiction.
\end{proof}

\begin{remarque}
  La droite de Sorgenfrey montre que \og{}séparable\fg{} n'implique pas
  \og{}second dénombrable\fg{} en dehors du cadre métrique. La
  
ef{prop:metrique-separable} ne se généralise donc pas aux espaces
  topologiques arbitraires.
\end{remarque}

% ============================================================
\section{Exemples résolus}\label{sec:exemples-resolus-ch3}
% ============================================================

\begin{exemple}\label{ex:base-produit}
  \textbf{Base de la topologie produit.}
  Soient $(X,\T_X)$ et $(Y,\T_Y)$ deux espaces topologiques avec des bases
  $\B_X$ et $\B_Y$ respectivement. Montrons que
  $\B = \{B_X\times B_Y : B_X\in\B_X,\, B_Y\in\B_Y\}$ est une base de la
  topologie produit sur $X\times Y$.

  \textbf{Solution.} Vérifions les conditions du 
ef{thm:critere-base}.
  \ref{B1}~: Pour $(x,y)\in X\times Y$, il existe $B_X\in\B_X$ avec
  $x\in B_X$ et $B_Y\in\B_Y$ avec $y\in B_Y$, donc
  $(x,y)\in B_X\times B_Y$. \ref{B2}~: Soient
  $B_1\times C_1,\, B_2\times C_2\in\B$ et
  $(x,y)\in (B_1\times C_1)\cap(B_2\times C_2) = (B_1\cap B_2)\times(C_1\cap C_2)$.
  Alors $x\in B_1\cap B_2$, donc il existe $B_3\in\B_X$ avec
  $x\in B_3\subseteq B_1\cap B_2$. De même, il existe $C_3\in\B_Y$ avec
  $y\in C_3\subseteq C_1\cap C_2$. Alors
  $(x,y)\in B_3\times C_3\subseteq (B_1\cap B_2)\times(C_1\cap C_2)$.

  La topologie engendrée par $\B$ coïncide avec la topologie produit
  $\T_X\times\T_Y$.
\end{exemple}

\begin{exemple}\label{ex:topo-engendree-R}
  \textbf{Topologie engendrée par les demi-droites.}
  Soit $\S = \{(-\infty,a) : a\in\R\}\cup\{(b,+\infty) : b\in\R\}$. Montrons
  que $\T(\S)$ est la topologie usuelle de $\R$.

  \textbf{Solution.} Les intersections finies d'éléments de $\S$ sont de la
  forme $(-\infty,a)\cap(b,+\infty)=(b,a)$ lorsque $b<a$, ou $\R$, ou
  $\varnothing$ (si $b\geq a$), ou les demi-droites elles-mêmes. Les
  intervalles ouverts $(b,a)$ forment une base de la topologie usuelle.
  Réciproquement, toute demi-droite est ouverte pour la topologie usuelle.
  Donc $\T(\S)$ est exactement la topologie usuelle.
\end{exemple}

\begin{exemple}\label{ex:sousbase-cofinite}
  \textbf{Sous-base de la topologie cofinie.}
  Soit $X$ un ensemble infini. Pour chaque $x\in X$, posons
  $S_x = X\setminus\{x\}$. La collection $\S=\{S_x : x\in X\}$ est une
  sous-base de la topologie cofinie\index{topologie cofinie} sur $X$.

  \textbf{Solution.} Une intersection finie $S_{x_1}\cap\cdots\cap S_{x_n}
  = X\setminus\{x_1,\ldots,x_n\}$ est le complémentaire d'un ensemble fini,
  donc un ouvert cofini. Les réunions de tels ensembles donnent tous les
  ouverts cofinis (et $\varnothing$).
\end{exemple}

% ============================================================
\section{Exercices}\label{sec:exercices-ch3}
% ============================================================

\begin{exercice}[$\bigstar$]
  Soit $(X,d)$ un espace métrique et $E\subseteq X$ un sous-ensemble dense.
  Montrer que $\{B(e,r) : e\in E,\, r\in\Q_{>0}\}$ est une base de la
  topologie induite par~$d$.
\end{exercice}

\begin{exercice}[$\bigstar$]
  Soient $\B_1$ et $\B_2$ deux bases de topologies $\T_1$ et $\T_2$ sur un
  même ensemble $X$. Montrer que $\T_1\subseteq\T_2$ si et seulement si, pour
  tout $B_1\in\B_1$ et tout $x\in B_1$, il existe $B_2\in\B_2$ tel que
  $x\in B_2\subseteq B_1$.
\end{exercice}

\begin{exercice}[$\bigstar$]
  Montrer que la topologie usuelle de $\R$ admet
  $\{(a,b) : a,b\in\Q,\, a<b\}$ comme base dénombrable. En déduire que $\R$
  est second dénombrable.
\end{exercice}

\begin{exercice}[$\bigstar\bigstar$]
  Soit $\S = \{(a,+\infty) : a\in\R\}$. Décrire la topologie $\T(\S)$
  engendrée par $\S$ sur $\R$. Cette topologie est-elle séparée~? Est-elle
  première dénombrable~?
\end{exercice}

\begin{exercice}[$\bigstar\bigstar$]\label{exo:sorgenfrey-produit}
  Soit $\R_\ell$ la droite de Sorgenfrey. Montrer que $\R_\ell\times\R_\ell$
  n'est pas un espace de Lindelöf.
  \emph{Indication~:} considérer l'anti-diagonale
  $\{(x,-x) : x\in\R\}$ et le recouvrement par des ouverts de la forme
  $[x,x+\varepsilon)\times[-x,-x+\varepsilon)$.
\end{exercice}

\begin{exercice}[$\bigstar\bigstar$]
  Soit $(X,\T)$ un espace topologique séparable. Montrer que toute collection
  d'ouverts deux à deux disjoints est dénombrable.
\end{exercice}

\begin{exercice}[$\bigstar\bigstar$]
  Soit $X$ un ensemble non dénombrable et $\T_{\mathrm{coc}}$ la topologie
  co-dénombrable (les ouverts sont $\varnothing$ et les complémentaires
  d'ensembles dénombrables). Montrer que $(X,\T_{\mathrm{coc}})$ n'est pas
  premier dénombrable.
\end{exercice}

\begin{exercice}[$\bigstar\bigstar\bigstar$]
  Soit $(X,\T)$ un espace de Lindelöf et $f\colon X\to Y$ une application
  continue et surjective. Montrer que $Y$ est un espace de Lindelöf.
\end{exercice}

\begin{exercice}[$\bigstar\bigstar\bigstar$]
  Montrer que tout sous-espace fermé d'un espace de Lindelöf est de Lindelöf.
  Donner un exemple de sous-espace ouvert d'un espace de Lindelöf qui n'est
  pas de Lindelöf.
  \emph{Indication~:} pour le contre-exemple, utiliser le plan de Sorgenfrey
  $\R_\ell\times\R_\ell$.
\end{exercice}

% ============================================================
\section*{Résumé du chapitre}
\addcontentsline{toc}{section}{Résumé du chapitre}
% ============================================================

\begin{itemize}
  \item Une \textbf{base} $\B$ d'une topologie $\T$ est une sous-collection de
    $\T$ telle que tout ouvert est réunion d'éléments de $\B$. Les conditions
    \ref{B1}--\ref{B2} caractérisent les collections qui sont des bases.
  \item Une \textbf{sous-base} $\S$ engendre une topologie dont une base est
    formée des intersections finies d'éléments de $\S$.
  \item La \textbf{topologie engendrée} par une collection $\mathcal{C}$ est la
    plus petite topologie contenant $\mathcal{C}$.
  \item Un espace est \textbf{second dénombrable} s'il admet une base
    dénombrable~; \textbf{premier dénombrable} si chaque point admet une base
    de voisinages dénombrable~; \textbf{séparable} s'il admet un sous-ensemble
    dense dénombrable.
  \item On a les implications~: second dénombrable
    $\Rightarrow$ premier dénombrable, et second dénombrable $\Rightarrow$
    séparable $+$ Lindelöf. Dans les espaces métriques, les trois propriétés
    (second dénombrable, séparable, Lindelöf) sont équivalentes.
  \item La \textbf{droite de Sorgenfrey} est un contre-exemple fondamental~:
    séparable et premier dénombrable, mais non second dénombrable.
\end{itemize}


% ============================================================================
%
%                         CHAPITRE 4
%
% ============================================================================

\chapter{Continuité, homéomorphismes, invariants topologiques}%
\label{ch:continuite-homeomorphismes}

\section*{Introduction}
\addcontentsline{toc}{section}{Introduction}

En analyse réelle, la continuité d'une fonction $f\colon\R\to\R$ en un point
$a$ est définie par la condition classique en $\varepsilon$-$\delta$~:
\[
  \forall\,\varepsilon>0,\quad \exists\,\delta>0,\quad
  |x-a|<\delta \;\Longrightarrow\; |f(x)-f(a)|<\varepsilon.
\]
Cette formulation fait appel à la métrique. Comment généraliser cette notion à
des espaces topologiques quelconques, où l'on ne dispose pas nécessairement
d'une distance~?

\medskip

La réponse réside dans la reformulation suivante~: $f$ est continue si et
seulement si l'\emph{image réciproque} de tout ouvert est un ouvert. Cette
caractérisation, purement topologique, ne fait intervenir ni $\varepsilon$ ni
$\delta$, mais seulement la structure des ouverts. Elle constitue le point de
départ de ce chapitre.

\medskip

Nous étudions ensuite les \emph{homéomorphismes}, qui sont les isomorphismes
de la catégorie des espaces topologiques~: des bijections continues dont la
réciproque est aussi continue. Deux espaces homéomorphes sont
\emph{topologiquement indistinguables}~: toute propriété exprimable en termes
d'ouverts est préservée. Ces propriétés sont les \emph{invariants
topologiques}.

\begin{center}
\begin{tikzpicture}[scale=1.0]
  % Left: interval (0,1)
  \draw[thick, blue!70!black] (0,0) -- (3,0);
  \draw[blue!70!black, fill=white] (0,0) circle (2pt);
  \draw[blue!70!black, fill=white] (3,0) circle (2pt);
  \node[below, blue!70!black] at (0,0) {$0$};
  \node[below, blue!70!black] at (3,0) {$1$};
  \node[above, blue!70!black] at (1.5,0.1) {$(0,1)$};
  % Arrow
  \draw[-{Stealth}, thick, red!60!black] (3.8,0.3) to[bend left=20]
    node[above,font=\small] {$f(x)=\tan\bigl(\pi(x-\tfrac{1}{2})\bigr)$}
    (6.2,0.3);
  \draw[-{Stealth}, thick, red!60!black] (6.2,-0.3) to[bend left=20]
    node[below,font=\small] {$f^{-1}$} (3.8,-0.3);
  % Right: real line
  \draw[thick, green!50!black, -{Stealth}] (6,0) -- (10.5,0);
  \draw[thick, green!50!black, {Stealth}-] (5.5,0) -- (6,0);
  \node[below, green!50!black] at (8,0) {$\R$};
  % Label
  \node[below, font=\small\itshape, text width=10cm, align=center] at (5,-1.3)
    {Homéomorphisme entre $(0,1)$ et $\R$ via la tangente.};
\end{tikzpicture}
\captionof{figure}{L'intervalle $(0,1)$ est homéomorphe à $\R$.}
\label{fig:homeo-01-R}
\end{center}

% ============================================================
\section{Applications continues}\label{sec:applications-continues}
% ============================================================

\subsection{Définition par image réciproque d'ouverts}\label{subsec:continuite-def}

\begin{definition}[Application continue]\label{def:application-continue}
  Soient $(X,\T_X)$ et $(Y,\T_Y)$ deux espaces topologiques. Une application
  $f\colon X\to Y$ est dite \emph{continue}\index{application continue} si
  l'image réciproque de tout ouvert de $Y$ est un ouvert de $X$~:
  \[
    \forall\, V\in\T_Y,\quad f^{-1}(V)\in\T_X.
  \]
\end{definition}

\begin{remarque}
  Si $(X,d_X)$ et $(Y,d_Y)$ sont des espaces métriques, on retrouve la
  continuité usuelle~: $f$ est continue au sens de la définition ci-dessus si
  et seulement si elle est continue en tout point au sens
  $\varepsilon$-$\delta$.
\end{remarque}

\subsection{Caractérisations équivalentes}\label{subsec:continuite-equiv}

\begin{theoreme}[Caractérisations de la continuité]\label{thm:carac-continuite}
  Soient $(X,\T_X)$ et $(Y,\T_Y)$ deux espaces topologiques et
  $f\colon X\to Y$ une application. Les conditions suivantes sont
  équivalentes~:
  \begin{enumerate}[label=\textup{(\roman*)}]
    \item\label{cont:ouverts} $f$ est continue (l'image réciproque de tout
      ouvert est un ouvert).
    \item\label{cont:fermes} L'image réciproque de tout fermé de $Y$ est un
      fermé de $X$.
    \item\label{cont:adherence} Pour tout $A\subseteq X$,
      $f(\overline{A})\subseteq\overline{f(A)}$.
    \item\label{cont:interieur} Pour tout $B\subseteq Y$,
      $f^{-1}(\mathring{B})\subseteq\bigl(f^{-1}(B)\bigr)^{\circ}$.
    \item\label{cont:voisinages} Pour tout $x\in X$ et tout voisinage $V$ de
      $f(x)$ dans $Y$, l'ensemble $f^{-1}(V)$ est un voisinage de $x$
      dans $X$.
    \item\label{cont:base} Il existe une base $\B$ de $\T_Y$ telle que
      $f^{-1}(B)\in\T_X$ pour tout $B\in\B$.
    \item\label{cont:sousbase} Il existe une sous-base $\S$ de $\T_Y$ telle
      que $f^{-1}(S)\in\T_X$ pour tout $S\in\S$.
  \end{enumerate}
\end{theoreme}

\begin{proof}
  $\ref{cont:ouverts}\Leftrightarrow\ref{cont:fermes}$~: Pour tout fermé $F$
  de $Y$, on a $f^{-1}(Y\setminus F) = X\setminus f^{-1}(F)$. Donc
  $f^{-1}(F)$ est fermé si et seulement si $f^{-1}(Y\setminus F)$ est
  ouvert. Comme les fermés de $Y$ sont exactement les complémentaires des
  ouverts, l'équivalence suit.

  $\ref{cont:ouverts}\Rightarrow\ref{cont:adherence}$~: Soit $A\subseteq X$.
  L'ensemble $\overline{f(A)}$ est fermé dans $Y$, donc
  $f^{-1}\bigl(\overline{f(A)}\bigr)$ est fermé dans $X$ par
  \ref{cont:fermes}. Puisque $A\subseteq f^{-1}(f(A))\subseteq
  f^{-1}\bigl(\overline{f(A)}\bigr)$, on a
  $\overline{A}\subseteq f^{-1}\bigl(\overline{f(A)}\bigr)$, d'où
  $f(\overline{A})\subseteq\overline{f(A)}$.

  $\ref{cont:adherence}\Rightarrow\ref{cont:fermes}$~: Soit $F$ fermé dans
  $Y$ et $A=f^{-1}(F)$. Alors
  $f(\overline{A})\subseteq\overline{f(A)}\subseteq\overline{F}=F$, donc
  $\overline{A}\subseteq f^{-1}(F)=A$. Comme $A\subseteq\overline{A}$
  toujours, on a $\overline{A}=A$, c'est-à-dire $f^{-1}(F)$ est fermé.

  $\ref{cont:ouverts}\Rightarrow\ref{cont:interieur}$~: Soit $B\subseteq Y$.
  L'ensemble $\mathring{B}$ est ouvert dans $Y$, donc
  $f^{-1}(\mathring{B})$ est ouvert dans $X$. Puisque
  $f^{-1}(\mathring{B})\subseteq f^{-1}(B)$ et $f^{-1}(\mathring{B})$ est
  ouvert, on a
  $f^{-1}(\mathring{B})\subseteq\bigl(f^{-1}(B)\bigr)^{\circ}$.

  $\ref{cont:interieur}\Rightarrow\ref{cont:ouverts}$~: Soit $V\in\T_Y$.
  Alors $V=\mathring{V}$, donc
  $f^{-1}(V) = f^{-1}(\mathring{V})\subseteq\bigl(f^{-1}(V)\bigr)^{\circ}
  \subseteq f^{-1}(V)$. Ainsi
  $f^{-1}(V)=\bigl(f^{-1}(V)\bigr)^{\circ}$, c'est-à-dire $f^{-1}(V)$ est
  ouvert.

  $\ref{cont:ouverts}\Rightarrow\ref{cont:voisinages}$~: Soit $V$ un
  voisinage de $f(x)$. Il existe $U\in\T_Y$ avec
  $f(x)\in U\subseteq V$. Alors $x\in f^{-1}(U)\subseteq f^{-1}(V)$ et
  $f^{-1}(U)\in\T_X$, donc $f^{-1}(V)$ est voisinage de $x$.

  $\ref{cont:voisinages}\Rightarrow\ref{cont:ouverts}$~: Soit $V\in\T_Y$.
  Pour tout $x\in f^{-1}(V)$, l'ensemble $V$ est un voisinage ouvert de
  $f(x)$, donc $f^{-1}(V)$ est un voisinage de $x$. Ainsi $f^{-1}(V)$ est
  voisinage de chacun de ses points, donc est ouvert.

  $\ref{cont:ouverts}\Rightarrow\ref{cont:base}$~: Trivial (prendre
  $\B=\T_Y$).

  $\ref{cont:base}\Rightarrow\ref{cont:ouverts}$~: Soit $V\in\T_Y$. Alors
  $V=\bigcup_{i\in I} B_i$ avec $B_i\in\B$. Donc
  $f^{-1}(V) = \bigcup_{i\in I} f^{-1}(B_i)$, qui est réunion d'ouverts de
  $\T_X$, donc est ouvert.

  $\ref{cont:base}\Rightarrow\ref{cont:sousbase}$~: Trivial (toute base est
  une sous-base).

  $\ref{cont:sousbase}\Rightarrow\ref{cont:base}$~: Soit $\S$ la sous-base
  donnée et $\B_{\S}$ la base associée (intersections finies d'éléments de
  $\S$). Pour $B=S_1\cap\cdots\cap S_n\in\B_{\S}$, on a
  $f^{-1}(B) = f^{-1}(S_1)\cap\cdots\cap f^{-1}(S_n)$, qui est une
  intersection finie d'ouverts, donc un ouvert. On conclut par
  $\ref{cont:base}\Rightarrow\ref{cont:ouverts}$.
\end{proof}

\subsection{Composition}\label{subsec:composition-continuite}

\begin{theoreme}[Composition d'applications continues]\label{thm:composition-continue}
  Soient $(X,\T_X)$, $(Y,\T_Y)$, $(Z,\T_Z)$ trois espaces topologiques. Si
  $f\colon X\to Y$ et $g\colon Y\to Z$ sont continues, alors
  $g\circ f\colon X\to Z$\index{composition!d'applications continues} est
  continue.
\end{theoreme}

\begin{proof}
  Soit $W\in\T_Z$. Alors $g^{-1}(W)\in\T_Y$ par continuité de $g$, puis
  $f^{-1}(g^{-1}(W))\in\T_X$ par continuité de $f$. Puisque
  $(g\circ f)^{-1}(W) = f^{-1}(g^{-1}(W))$, le résultat suit.
\end{proof}

\begin{corollaire}\label{cor:cor:identite-continue}
  L'identité $\mathrm{id}_X\colon(X,\T)\to(X,\T)$ est continue. L'identité
  $\mathrm{id}_X\colon(X,\T_1)\to(X,\T_2)$ est continue si et seulement si
  $\T_2\subseteq\T_1$ (i.e.\ $\T_1$ est plus fine que $\T_2$).
\end{corollaire}

\begin{proof}
  Pour la première assertion, $\mathrm{id}_X^{-1}(U)=U\in\T$ pour tout
  $U\in\T$. Pour la seconde, $\mathrm{id}_X^{-1}(U)=U$ est dans $\T_1$
  pour tout $U\in\T_2$ si et seulement si $\T_2\subseteq\T_1$.
\end{proof}

% ============================================================
\section{Homéomorphismes}\label{sec:homeomorphismes}
% ============================================================

\subsection{Définition}\label{subsec:homeo-def}

\begin{definition}[Homéomorphisme]\label{def:homeomorphisme}
  Soient $(X,\T_X)$ et $(Y,\T_Y)$ deux espaces topologiques. Un
  \emph{homéomorphisme}\index{homéomorphisme} de $X$ sur $Y$ est une
  bijection $f\colon X\to Y$ telle que $f$ et $f^{-1}$ soient toutes deux
  continues. On dit alors que $X$ et $Y$ sont \emph{homéomorphes} et on
  note $X\cong Y$.
\end{definition}

\begin{remarque}
  La relation \og{}être homéomorphe\fg{} est une relation d'équivalence sur la
  classe des espaces topologiques~: réflexivité (via $\mathrm{id}$), symétrie
  (si $f$ est un homéomorphisme, $f^{-1}$ aussi), transitivité (composition de
  deux homéomorphismes, par le 
ef{thm:composition-continue}).
\end{remarque}

\begin{definition}[Invariant topologique]\label{def:invariant-topologique}
  Une propriété $\mathcal{P}$ des espaces topologiques est un \emph{invariant
  topologique}\index{invariant topologique} si, chaque fois qu'un espace
  $(X,\T_X)$ satisfait $\mathcal{P}$ et qu'un espace $(Y,\T_Y)$ est
  homéomorphe à $X$, l'espace $Y$ satisfait aussi $\mathcal{P}$.
\end{definition}

\begin{exemple}
  La séparabilité, la compacité, la connexité, les axiomes de séparation, les
  axiomes de dénombrabilité sont des invariants topologiques. Le diamètre
  d'un espace métrique ou le caractère borné ne sont \emph{pas} des invariants
  topologiques ($(0,1)\cong\R$, mais $(0,1)$ est borné et $\R$ ne l'est pas).
\end{exemple}

\subsection{Applications ouvertes et fermées}\label{subsec:ouvertes-fermees}

\begin{definition}[Application ouverte, application fermée]\label{def:ouverte-fermee}
  Soient $(X,\T_X)$ et $(Y,\T_Y)$ deux espaces topologiques et
  $f\colon X\to Y$ une application.
  \begin{itemize}
    \item $f$ est \emph{ouverte}\index{application ouverte} si l'image de tout
      ouvert de $X$ est un ouvert de $Y$~: $\forall\, U\in\T_X,\;
      f(U)\in\T_Y$.
    \item $f$ est \emph{fermée}\index{application fermée} si l'image de tout
      fermé de $X$ est un fermé de $Y$.
  \end{itemize}
\end{definition}

\begin{remarque}
  Une application continue n'est en général ni ouverte ni fermée, et
  réciproquement. Par exemple, l'application $f\colon\R\to\R$,
  $f(x)=x^2$ est continue mais pas ouverte ($f(\R)=[0,+\infty)$ qui est
  bien ouvert... mais $f$ n'est pas injective, ce qui crée des
  subtilités). L'inclusion $\iota\colon(0,1)\hookrightarrow\R$ est
  continue et ouverte (dans la topologie induite), mais l'image de
  l'ensemble fermé $(0,1)$ (fermé dans $(0,1)$) est $(0,1)$, qui n'est
  pas fermé dans $\R$.
\end{remarque}

\begin{proposition}[Caractérisation d'un homéomorphisme]\label{prop:carac-homeo}
  Soit $f\colon X\to Y$ une bijection. Les conditions suivantes sont
  équivalentes~:
  \begin{enumerate}[label=\textup{(\roman*)}]
    \item $f$ est un homéomorphisme.
    \item $f$ est continue et ouverte.
    \item $f$ est continue et fermée.
  \end{enumerate}
\end{proposition}

\begin{proof}
  $\mathrm{(i)\Leftrightarrow(ii)}$~: Si $f$ est un homéomorphisme, alors
  $f^{-1}$ est continue, c'est-à-dire, pour tout $U\in\T_X$,
  $(f^{-1})^{-1}(U) = f(U)\in\T_Y$~: $f$ est ouverte. Réciproquement, si $f$
  est bijective, continue et ouverte, alors pour tout $U\in\T_X$,
  $f(U)\in\T_Y$, i.e.\ $(f^{-1})^{-1}(U)\in\T_Y$, donc $f^{-1}$ est
  continue.

  $\mathrm{(i)\Leftrightarrow(iii)}$~: Analogue, en remplaçant
  \og{}ouvert\fg{} par \og{}fermé\fg{} et en utilisant la caractérisation
  \ref{cont:fermes} du 
ef{thm:carac-continuite}.
\end{proof}

% ============================================================
\section{Le théorème compact-Hausdorff}\label{sec:compact-hausdorff}
% ============================================================

\begin{theoreme}[Bijection continue d'un compact vers un Hausdorff]\label{thm:compact-hausdorff}
  Soit $f\colon X\to Y$ une bijection continue. Si $X$ est
  compact\index{espace compact} et $Y$ est séparé (Hausdorff)\index{espace
  de Hausdorff}, alors $f$ est un homéomorphisme.
\end{theoreme}

\begin{proof}
  D'après la 
ef{prop:carac-homeo}, il suffit de montrer que $f$ est
  fermée. Soit $F\subseteq X$ un fermé. Puisque $X$ est compact, $F$ est
  compact (fermé dans un compact). L'image $f(F)$ est compacte (image
  continue d'un compact). Puisque $Y$ est séparé, tout compact de $Y$ est
  fermé. Donc $f(F)$ est fermé dans $Y$.
\end{proof}

\begin{remarque}
  Ce théorème est extrêmement utile en pratique~: il suffit de construire une
  bijection continue d'un espace compact vers un espace séparé pour obtenir
  un homéomorphisme. On l'applique fréquemment en théorie des quotients.
\end{remarque}

% ============================================================
\section{Exemples d'homéomorphismes}\label{sec:exemples-homeo}
% ============================================================

\begin{exemple}\label{ex:homeo-01-R}
  \textbf{$(0,1)\cong\R$.}
  L'application $f\colon(0,1)\to\R$ définie par
  $f(x)=\tan\bigl(\pi(x-\tfrac{1}{2})\bigr)$ est un homéomorphisme. En
  effet~:
  \begin{itemize}
    \item $f$ est la composée $\tan\circ g$ où $g(x)=\pi(x-1/2)$ envoie
      $(0,1)$ sur $(-\pi/2,\pi/2)$. La fonction $g$ est un homéomorphisme
      (bijection affine continue à réciproque continue). La tangente est un
      homéomorphisme de $(-\pi/2,\pi/2)$ sur $\R$ (bijection continue de
      classe $\mathcal{C}^\infty$ avec réciproque $\arctan$ continue).
    \item Par composition (
ef{thm:composition-continue}), $f$ est un
      homéomorphisme.
  \end{itemize}
  Ceci montre que le caractère \og{}borné\fg{} n'est pas un invariant
  topologique.
\end{exemple}

\begin{exemple}\label{ex:homeo-S1-quotient}
  \textbf{$\mathbb{S}^1\cong[0,1]/\!\!\sim$.}
  Considérons l'espace quotient $[0,1]/\!\!\sim$ où $0\sim 1$ (les
  extrémités sont identifiées). L'application
  $f\colon[0,1]\to\mathbb{S}^1$ définie par $f(t)=e^{2\pi i t}$ est
  continue et surjective, et passe au quotient en une bijection continue
  $\bar{f}\colon[0,1]/\!\!\sim\;\to\mathbb{S}^1$. L'espace $[0,1]$ est
  compact, donc $[0,1]/\!\!\sim$ est compact (image continue d'un compact).
  Le cercle $\mathbb{S}^1\subseteq\R^2$ est séparé (sous-espace d'un
  Hausdorff). Par le 
ef{thm:compact-hausdorff}, $\bar{f}$ est un
  homéomorphisme.
\end{exemple}

\begin{exemple}\label{ex:homeo-Rn}
  \textbf{$\R^n\cong B(0,1)$ (boule ouverte unité).}
  L'application $f\colon B(0,1)\to\R^n$ définie par
  $f(x)=\frac{x}{1-\|x\|}$ est un homéomorphisme. Sa réciproque est
  $f^{-1}(y)=\frac{y}{1+\|y\|}$. Les deux applications sont continues
  (compositions d'opérations continues), et l'on vérifie directement que
  $f\circ f^{-1}=\mathrm{id}_{\R^n}$ et $f^{-1}\circ f=\mathrm{id}_{B(0,1)}$.
\end{exemple}

% ============================================================
\section{Plongements topologiques}\label{sec:plongements}
% ============================================================

\begin{definition}[Plongement topologique]\label{def:plongement}
  Soient $(X,\T_X)$ et $(Y,\T_Y)$ deux espaces topologiques. Un
  \emph{plongement}\index{plongement topologique} (ou
  \emph{embedding}\index{embedding|see{plongement topologique}}) est une
  application $f\colon X\to Y$ qui est un homéomorphisme de $X$ sur
  $f(X)$ muni de la topologie induite par $Y$.
\end{definition}

\begin{proposition}[Caractérisation d'un plongement]\label{prop:carac-plongement}
  Une application $f\colon X\to Y$ est un plongement si et seulement si $f$
  est injective, continue, et pour tout ouvert $U$ de $X$, il existe un ouvert
  $V$ de $Y$ tel que $f(U)=V\cap f(X)$.
\end{proposition}

\begin{proof}
  C'est la reformulation de la définition~: $f$ est un homéomorphisme sur
  son image si et seulement si $f$ est une bijection continue de $X$ sur
  $f(X)$ dont la réciproque $f^{-1}\colon f(X)\to X$ est continue. La
  continuité de $f^{-1}$ signifie que pour tout ouvert $U$ de $X$,
  $(f^{-1})^{-1}(U)=f(U)$ est ouvert dans la topologie induite sur $f(X)$,
  c'est-à-dire $f(U)=V\cap f(X)$ pour un certain $V\in\T_Y$.
\end{proof}

\begin{definition}[Propriété topologique]\label{def:propriete-topologique}
  Une \emph{propriété topologique}\index{propriété topologique} est une
  propriété des espaces topologiques qui est préservée par homéomorphisme.
  C'est un synonyme d'invariant topologique (
ef{def:invariant-topologique}).
\end{definition}

% ============================================================
\section{Lemme de recollement}\label{sec:lemme-recollement}
% ============================================================

\begin{theoreme}[Lemme de recollement (Pasting lemma)]\label{thm:pasting-lemma}
  Soit $X = A\cup B$ où $A$ et $B$ sont deux fermés\index{lemme de
  recollement} de $X$. Soient $f\colon A\to Y$ et $g\colon B\to Y$ deux
  applications continues telles que $f(x)=g(x)$ pour tout $x\in A\cap B$.
  Alors l'application $h\colon X\to Y$ définie par
  \[
    h(x) = \begin{cases}
      f(x) & \text{si } x\in A, \\
      g(x) & \text{si } x\in B,
    \end{cases}
  \]
  est continue.
\end{theoreme}

\begin{proof}
  Nous utilisons la caractérisation \ref{cont:fermes} du
  
ef{thm:carac-continuite}. Soit $F$ un fermé de $Y$. Alors
  \[
    h^{-1}(F) = \{x\in X : h(x)\in F\}
    = \{x\in A : f(x)\in F\} \cup \{x\in B : g(x)\in F\}
    = f^{-1}(F) \cup g^{-1}(F).
  \]
  Par continuité de $f$, l'ensemble $f^{-1}(F)$ est fermé dans $A$ pour la
  topologie induite. Puisque $A$ est fermé dans $X$, un fermé de $A$ pour la
  topologie induite est fermé dans $X$ (intersection d'un fermé de $X$ avec
  $A$, qui est fermé). Donc $f^{-1}(F)$ est fermé dans $X$. De même,
  $g^{-1}(F)$ est fermé dans $X$. Ainsi $h^{-1}(F) = f^{-1}(F)\cup g^{-1}(F)$
  est réunion finie de fermés, donc fermé dans $X$.
\end{proof}

\begin{remarque}
  L'hypothèse que $A$ et $B$ soient fermés est essentielle. Si l'on remplace
  \og{}fermés\fg{} par \og{}ouverts\fg{}, le résultat reste vrai (la preuve
  est analogue, en utilisant la caractérisation par les ouverts). En revanche,
  si l'un est ouvert et l'autre fermé, le lemme peut être en défaut.
\end{remarque}

\begin{corollaire}[Recollement sur un recouvrement ouvert]\label{cor:cor:recollement-ouvert}
  Soit $(U_i)_{i\in I}$ un recouvrement ouvert de $X$. Si
  $f_i\colon U_i\to Y$ sont des applications continues telles que
  $f_i|_{U_i\cap U_j}=f_j|_{U_i\cap U_j}$ pour tous $i,j\in I$, alors
  l'application $f\colon X\to Y$ définie par $f(x)=f_i(x)$ pour $x\in U_i$
  est bien définie et continue.
\end{corollaire}

\begin{proof}
  La bonne définition résulte de la condition de compatibilité. Pour la
  continuité, soit $V\in\T_Y$. Alors
  $f^{-1}(V) = \bigcup_{i\in I} f_i^{-1}(V)$. Chaque $f_i^{-1}(V)$ est
  ouvert dans $U_i$ (par continuité de $f_i$), donc ouvert dans $X$ (car
  $U_i$ est ouvert). Ainsi $f^{-1}(V)$ est réunion d'ouverts, donc ouvert.
\end{proof}

% ============================================================
\section{Contre-exemples}\label{sec:contre-exemples-ch4}
% ============================================================

\begin{exemple}\label{ex:bijection-continue-non-homeo}
  \textbf{Bijection continue qui n'est pas un homéomorphisme.}
  Considérons $X=[0,2\pi)$ muni de la topologie induite par $\R$ et
  $Y=\mathbb{S}^1$ le cercle unité dans $\R^2$. L'application
  $f\colon X\to Y$ définie par $f(t)=(\cos t,\sin t)$ est continue (composantes
  continues), bijective (chaque point du cercle a un unique antécédent dans
  $[0,2\pi)$), mais $f^{-1}$ n'est pas continue.

  En effet, considérons la suite $(y_n)$ dans $\mathbb{S}^1$ définie par
  $y_n = (\cos(2\pi-1/n),\sin(2\pi-1/n))$ pour $n\geq 1$. On a $y_n\to(1,0)$
  dans $\mathbb{S}^1$. Si $f^{-1}$ était continue, on aurait
  $f^{-1}(y_n)=2\pi-1/n\to f^{-1}(1,0)=0$. Mais $2\pi-1/n\to 2\pi\neq 0$.
  Donc $f^{-1}$ n'est pas continue.

  Cet exemple n'est pas en contradiction avec le 
ef{thm:compact-hausdorff}~:
  $X=[0,2\pi)$ n'est pas compact.
\end{exemple}

\begin{remarque}
  Si l'on considère le quotient $[0,2\pi]/\!\!\sim$ (où $0\sim 2\pi$),
  l'application induite est un homéomorphisme sur $\mathbb{S}^1$, car le
  quotient d'un compact est compact. C'est l'
ef{ex:homeo-S1-quotient}.
\end{remarque}

% ============================================================
\section{Invariance du domaine}\label{sec:invariance-domaine}
% ============================================================

Nous énonçons sans démonstration le théorème suivant, dont la preuve fait appel
à l'homologie singulière.

\begin{theoreme}[Invariance du domaine (Brouwer)]\label{thm:invariance-domaine}
  Soient $U$ un ouvert de $\R^n$ et $f\colon U\to\R^n$ une application
  continue et injective. Alors $f(U)$ est un ouvert de
  $\R^n$\index{invariance du domaine} et $f$ est un homéomorphisme de $U$
  sur $f(U)$.
\end{theoreme}

\begin{remarque}
  Ce théorème a des conséquences profondes~: il implique par exemple que $\R^m$
  et $\R^n$ ne sont homéomorphes que si $m=n$, c'est-à-dire que la
  \emph{dimension topologique} de $\R^n$ est un invariant topologique.
\end{remarque}

% ============================================================
\section{Exemples résolus}\label{sec:exemples-resolus-ch4}
% ============================================================

\begin{exemple}\label{ex:continuite-projections}
  \textbf{Continuité des projections.}
  Soient $(X,\T_X)$ et $(Y,\T_Y)$ deux espaces topologiques et
  $\pi_X\colon X\times Y\to X$ la projection sur le premier facteur, où
  $X\times Y$ est muni de la topologie produit. Montrons que $\pi_X$ est
  continue et ouverte.

  \textbf{Solution.} Pour la continuité, soit $U\in\T_X$. Alors
  $\pi_X^{-1}(U) = U\times Y$, qui est ouvert dans la topologie produit
  (c'est un ouvert de base).

  Pour le caractère ouvert, il suffit de vérifier sur une base. Un ouvert de
  base de la topologie produit est de la forme $U\times V$ avec $U\in\T_X$,
  $V\in\T_Y$. On a $\pi_X(U\times V)=U$ si $V\neq\varnothing$, et
  $\pi_X(\varnothing)=\varnothing$. Dans les deux cas, l'image est ouverte.
  Puisque tout ouvert est réunion d'ouverts de base, et que l'image préserve
  les réunions, $\pi_X$ est ouverte.
\end{exemple}

\begin{exemple}\label{ex:homeo-intervalle}
  \textbf{$[a,b]\cong[c,d]$ pour tous $a<b$ et $c<d$.}
  L'application $f\colon[a,b]\to[c,d]$ définie par
  $f(x) = c + \frac{(d-c)(x-a)}{b-a}$ est une bijection affine continue
  (restriction d'une fonction affine $\R\to\R$). Sa réciproque
  $f^{-1}(y) = a + \frac{(b-a)(y-c)}{d-c}$ est aussi affine et continue.
  Donc $f$ est un homéomorphisme.

  On peut aussi invoquer le 
ef{thm:compact-hausdorff}~: $[a,b]$ est
  compact, $[c,d]$ est séparé, et $f$ est bijective et continue.
\end{exemple}

\begin{exemple}\label{ex:exponentielle-non-homeo}
  \textbf{$\exp\colon\R\to(0,+\infty)$ est un homéomorphisme.}
  L'exponentielle $\exp\colon\R\to(0,+\infty)$ est continue, bijective, et sa
  réciproque $\ln\colon(0,+\infty)\to\R$ est continue. C'est donc un
  homéomorphisme.

  En revanche, $\exp\colon\R\to\R$ (vue comme application dans $\R$ entier)
  n'est pas surjective, donc n'est pas un homéomorphisme. C'est un
  plongement topologique de $\R$ dans $\R$.
\end{exemple}

% ============================================================
\section{Exercices}\label{sec:exercices-ch4}
% ============================================================

\begin{exercice}[$\bigstar$]
  Montrer qu'une application constante $f\colon X\to Y$ est continue pour
  toutes topologies sur $X$ et $Y$.
\end{exercice}

\begin{exercice}[$\bigstar$]
  Soit $f\colon(X,\T_X)\to(Y,\T_Y)$ continue et soit $A\subseteq X$. Montrer
  que la restriction $f|_A\colon A\to Y$ est continue lorsque $A$ est muni de
  la topologie induite.
\end{exercice}

\begin{exercice}[$\bigstar$]
  Montrer que l'application $f\colon\R\to\R$ définie par $f(x)=x^3$ est un
  homéomorphisme.
\end{exercice}

\begin{exercice}[$\bigstar\bigstar$]
  Soit $f\colon X\to Y$ une application. Montrer que $f$ est continue si et
  seulement si, pour tout $A\subseteq X$,
  $f(\overline{A})\subseteq\overline{f(A)}$. Montrer que $f$ est fermée si et
  seulement si, pour tout $A\subseteq X$,
  $\overline{f(A)}\subseteq f(\overline{A})$.
\end{exercice}

\begin{exercice}[$\bigstar\bigstar$]\label{exo:graphe-ferme}
  Soit $f\colon X\to Y$ une application continue et $Y$ un espace séparé.
  Montrer que le graphe $\Gamma_f = \{(x,f(x)) : x\in X\}\subseteq X\times Y$
  est fermé dans $X\times Y$.
\end{exercice}

\begin{exercice}[$\bigstar\bigstar$]
  Soient $X$ et $Y$ deux espaces topologiques, $f\colon X\to Y$ continue, et
  $\B$ une base de la topologie de $Y$. Montrer que
  $\{f^{-1}(B) : B\in\B\}$ est une sous-base de la topologie
  $\{f^{-1}(V) : V\in\T_Y\}$ sur $X$ (topologie initiale induite par $f$).
\end{exercice}

\begin{exercice}[$\bigstar\bigstar\bigstar$]
  Montrer que $\R$ et $\R^2$ ne sont pas homéomorphes.
  \emph{Indication~:} si $f\colon\R\to\R^2$ est un homéomorphisme, considérer
  la restriction de $f$ à $\R\setminus\{0\}$ et utiliser la connexité.
\end{exercice}

\begin{exercice}[$\bigstar\bigstar\bigstar$]
  Soit $f\colon[0,1]\to[0,1]$ une application continue. Montrer que $f$ admet
  un point fixe (théorème de point fixe de Brouwer en dimension $1$).
  \emph{Indication~:} considérer $g(x)=f(x)-x$ et appliquer le théorème des
  valeurs intermédiaires.
\end{exercice}

\begin{exercice}[$\bigstar\bigstar\bigstar$]
  Soit $X$ un espace compact et $Y$ un espace de Hausdorff. Montrer que toute
  application continue $f\colon X\to Y$ est fermée. En déduire le
  
ef{thm:compact-hausdorff} directement.
\end{exercice}

% ============================================================
\section*{Résumé du chapitre}
\addcontentsline{toc}{section}{Résumé du chapitre}
% ============================================================

\begin{itemize}
  \item Une application $f\colon X\to Y$ est \textbf{continue} si et seulement
    si l'image réciproque de tout ouvert est un ouvert. Les caractérisations
    équivalentes (fermés, adhérence, intérieur, voisinages, base, sous-base)
    sont résumées dans le 
ef{thm:carac-continuite}.
  \item La \textbf{composition} de deux applications continues est continue.
  \item Un \textbf{homéomorphisme} est une bijection continue dont la
    réciproque est continue. C'est équivalent à une bijection continue et
    ouverte, ou continue et fermée.
  \item Les \textbf{invariants topologiques} sont les propriétés préservées par
    homéomorphisme. Le caractère borné n'en est pas un~: $(0,1)\cong\R$.
  \item Le \textbf{théorème compact-Hausdorff}~: une bijection continue d'un
    compact dans un séparé est un homéomorphisme.
  \item Le \textbf{lemme de recollement}~: si $X=A\cup B$ avec $A,B$ fermés
    (ou ouverts), et si des applications continues coïncident sur $A\cap B$,
    l'application recollée est continue.
  \item Une bijection continue qui n'est pas un homéomorphisme existe dès que
    la source n'est pas compacte~: l'exemple
    $[0,2\pi)\to\mathbb{S}^1$.
  \item L'\textbf{invariance du domaine} (Brouwer)~: une injection continue
    d'un ouvert de $\R^n$ dans $\R^n$ est un plongement ouvert.
\end{itemize}
% chapters5_6.tex — Chapitres 5 et 6 : Axiomes de séparation et Compacité
% Partie 3 du cours de topologie générale.
% Conventions Bourbaki. Preuves complètes.

% ============================================================================
% CHAPITRE 5 : AXIOMES DE SÉPARATION
% ============================================================================

\chapter{Axiomes de séparation}\label{ch:separation}

\section*{Introduction}
\addcontentsline{toc}{section}{Introduction}

Les espaces métriques\index{espace métrique} jouissent de propriétés
remarquables~: les points distincts peuvent toujours être séparés par des
ouverts disjoints, les singletons sont fermés, les limites de suites
convergentes sont uniques. Ces propriétés ne sont pas automatiques dans un
espace topologique quelconque. Les \emph{axiomes de séparation}\index{axiomes
de séparation} constituent une hiérarchie de conditions qui permettent de
retrouver, dans un cadre purement topologique, les comportements
<<~raisonnables~>> auxquels l'intuition métrique nous habitue.

\medskip

La terminologie classique, due à Alexandroff et Hopf, utilise les notations
$T_0$, $T_1$, $T_2$, $T_3$, $T_4$ (du mot allemand \emph{Trennungsaxiom},
axiome de séparation). Chaque axiome impose une capacité de séparation
croissante. Ce chapitre définit ces axiomes, étudie leurs implications
mutuelles, et construit des contre-exemples montrant qu'aucune implication
ne se renverse.

\medskip

\begin{center}
\begin{tikzpicture}[
  node distance=1.8cm,
  every node/.style={draw, rounded corners, minimum width=2.2cm,
    minimum height=0.8cm, font=\small},
  arrow/.style={-{Stealth}, thick}
]
  \node (T4) {$T_4$ (normal)};
  \node (T3) [right=of T4] {$T_3$ (régulier)};
  \node (T2) [right=of T3] {$T_2$ (Hausdorff)};
  \node (T1) [right=of T2] {$T_1$};
  \node (T0) [right=of T1] {$T_0$};
  \draw[arrow] (T4) -- (T3);
  \draw[arrow] (T3) -- (T2);
  \draw[arrow] (T2) -- (T1);
  \draw[arrow] (T1) -- (T0);
  \node[draw=none, below=0.4cm of T2, font=\footnotesize\itshape,
    minimum width=0pt] {Aucune implication ne se renverse.};
\end{tikzpicture}
\captionof{figure}{Hiérarchie des axiomes de séparation.}
\label{fig:hierarchie-separation}
\end{center}

% ============================================================
\section{L'axiome $T_0$ (Kolmogorov)}\label{sec:T0}
% ============================================================

\begin{definition}[Espace $T_0$]\label{def:T0}
  Un espace topologique\index{espace T0@espace $T_0$}\index{Kolmogorov}
  $(X,\tau)$ est dit \emph{$T_0$} (ou espace de \emph{Kolmogorov}) si, pour
  tout couple de points distincts $x,y\in X$, il existe un ouvert $U\in\tau$
  contenant l'un des deux points mais pas l'autre~:
  \[
    \forall\, x\neq y\in X,\quad
    \exists\, U\in\tau,\quad
    \bigl(x\in U,\; y\notin U\bigr)
    \;\text{ou}\;
    \bigl(y\in U,\; x\notin U\bigr).
  \]
\end{definition}

\begin{remarque}[Interprétation]\label{rem:rem:T0-interpretation}
  L'axiome $T_0$ signifie que la topologie distingue topologiquement tout
  couple de points~: deux points distincts ne peuvent pas avoir exactement les
  mêmes voisinages ouverts.
\end{remarque}

\begin{exemple}[Exemples d'espaces $T_0$]\label{ex:T0}
  \begin{enumerate}
    \item Tout espace métrique est $T_0$.
    \item La \emph{topologie de Sierpi\'nski}\index{topologie de Sierpinski@topologie de Sierpi\'nski}
      sur $\{0,1\}$, dont les ouverts sont $\emptyset$, $\{0\}$ et $\{0,1\}$,
      est $T_0$ mais pas $T_1$.
    \item L'espace $(X,\{\emptyset, X\})$ muni de la topologie grossière n'est
      pas $T_0$ dès que $|X|\geq 2$.
  \end{enumerate}
\end{exemple}

% ============================================================
\section{L'axiome $T_1$}\label{sec:T1}
% ============================================================

\begin{definition}[Espace $T_1$]\label{def:T1}
  Un espace topologique\index{espace T1@espace $T_1$} $(X,\tau)$ est dit
  \emph{$T_1$} si, pour tout couple de points distincts $x,y\in X$, il existe
  un ouvert contenant $x$ mais pas $y$~:
  \[
    \forall\, x\neq y\in X,\quad
    \exists\, U\in\tau,\quad x\in U,\; y\notin U.
  \]
\end{definition}

\begin{remarque}[Différence avec $T_0$]\label{rem:rem:T1-vs-T0}
  Dans un espace $T_0$, on peut séparer \emph{l'un des deux} points de
  l'autre ; dans un espace $T_1$, chaque point peut être séparé de l'autre par
  un ouvert qui lui est propre. Tout espace $T_1$ est $T_0$, mais la réciproque
  est fausse (l'espace de Sierpi\'nski est $T_0$ mais pas $T_1$).
\end{remarque}

\begin{theoreme}[Caractérisation des espaces $T_1$]\label{thm:T1-singletons}
  Un espace topologique $(X,\tau)$ est $T_1$ si et seulement si tout singleton
  $\{x\}$ est un fermé\index{singleton fermé} de $X$.
\end{theoreme}

\begin{proof}
  $(\Rightarrow)$ Soit $x\in X$. Pour tout $y\neq x$, il existe un ouvert
  $U_y\in\tau$ tel que $y\in U_y$ et $x\notin U_y$. Alors
  \[
    X\setminus\{x\} = \bigcup_{y\neq x} U_y
  \]
  est une réunion d'ouverts, donc un ouvert. Ainsi $\{x\}$ est fermé.

  \medskip
  $(\Leftarrow)$ Soient $x\neq y$. Comme $\{y\}$ est fermé, $U=X\setminus\{y\}$
  est un ouvert contenant $x$ mais pas $y$.
\end{proof}

\begin{exemple}[Topologie cofinie]\label{ex:cofinie-T1}
  Soit $X$ un ensemble infini\index{topologie cofinie}. La \emph{topologie
  cofinie} sur $X$ est
  \[
    \tau_{\mathrm{cof}} = \{\emptyset\} \cup \{U\subseteq X : X\setminus U
    \text{ est fini}\}.
  \]
  Les singletons $\{x\}$ sont fermés (leur complémentaire $X\setminus\{x\}$
  est cofini donc ouvert). Par le théorème~\ref{thm:T1-singletons}, $(X,\tau_{\mathrm{cof}})$
  est $T_1$.

  Cependant, cet espace n'est \emph{pas} $T_2$~: si $U$ et $V$ sont deux
  ouverts non vides, alors $X\setminus U$ et $X\setminus V$ sont finis, donc
  $X\setminus(U\cap V) = (X\setminus U)\cup(X\setminus V)$ est fini, ce qui
  implique $U\cap V\neq\emptyset$. Il est donc impossible de trouver deux
  ouverts disjoints non vides.
\end{exemple}

% ============================================================
\section{L'axiome $T_2$ (Hausdorff)}\label{sec:T2}
% ============================================================

\begin{definition}[Espace de Hausdorff]\label{def:T2}
  Un espace topologique\index{espace de Hausdorff}\index{espace T2@espace $T_2$}
  $(X,\tau)$ est dit \emph{$T_2$} ou \emph{séparé} (ou espace de
  \emph{Hausdorff}) si, pour tout couple de points distincts $x,y\in X$, il
  existe des ouverts disjoints $U,V\in\tau$ tels que $x\in U$ et $y\in V$~:
  \[
    \forall\, x\neq y\in X,\quad
    \exists\, U,V\in\tau,\quad
    x\in U,\; y\in V,\; U\cap V=\emptyset.
  \]
\end{definition}

\begin{theoreme}[Unicité des limites dans les espaces de Hausdorff]\label{thm:unicite-limites}
  Dans un espace de Hausdorff\index{unicité des limites}, les limites de suites
  (plus généralement de filtres ou de filets) sont uniques.
\end{theoreme}

\begin{proof}
  Soit $(X,\tau)$ un espace de Hausdorff et soit $(x_n)_{n\in\N}$ une suite
  admettant deux limites $\ell_1$ et $\ell_2$ avec $\ell_1\neq\ell_2$. Par
  l'axiome $T_2$, il existe des ouverts disjoints $U$ et $V$ tels que
  $\ell_1\in U$ et $\ell_2\in V$. Par définition de la convergence, il existe
  $N_1\in\N$ tel que $x_n\in U$ pour tout $n\geq N_1$, et $N_2\in\N$ tel que
  $x_n\in V$ pour tout $n\geq N_2$. Pour $n\geq\max(N_1,N_2)$, on aurait
  $x_n\in U\cap V=\emptyset$, ce qui est absurde.

  \medskip
  Le même argument s'adapte aux filtres~: si un filtre $\mathcal{F}$ converge
  vers $\ell_1$ et $\ell_2$, alors $U\in\mathcal{F}$ et $V\in\mathcal{F}$,
  donc $U\cap V\in\mathcal{F}$, mais $U\cap V=\emptyset\notin\mathcal{F}$,
  contradiction.
\end{proof}

\begin{proposition}[Caractérisation par la diagonale]\label{prop:diagonale}
  Un espace topologique $(X,\tau)$ est de Hausdorff si et seulement si la
  diagonale\index{diagonale}
  \[
    \Delta_X = \{(x,x) : x\in X\}
  \]
  est un fermé de $X\times X$ muni de la topologie produit.
\end{proposition}

\begin{proof}
  $(\Rightarrow)$ Montrons que $(X\times X)\setminus\Delta_X$ est ouvert. Soit
  $(x,y)\in(X\times X)\setminus\Delta_X$, c'est-à-dire $x\neq y$. Par
  l'axiome $T_2$, il existe $U,V$ ouverts disjoints avec $x\in U$, $y\in V$.
  Alors $U\times V$ est un ouvert de $X\times X$ contenant $(x,y)$ et
  $(U\times V)\cap\Delta_X=\emptyset$ car si $(z,z)\in U\times V$ alors
  $z\in U\cap V=\emptyset$, contradiction.

  \medskip
  $(\Leftarrow)$ Si $\Delta_X$ est fermé, son complémentaire est ouvert. Pour
  $x\neq y$, le point $(x,y)\notin\Delta_X$ admet un voisinage ouvert de base
  $U\times V$ disjoint de $\Delta_X$, ce qui impose $U\cap V=\emptyset$.
\end{proof}

\begin{theoreme}[Compacts dans un espace de Hausdorff]\label{thm:compact-hausdorff-ferme}
  Tout sous-ensemble compact d'un espace de Hausdorff est fermé\index{compact!fermé dans Hausdorff}.
\end{theoreme}

\begin{proof}
  Soit $K$ un compact de l'espace de Hausdorff $(X,\tau)$. Montrons que
  $X\setminus K$ est ouvert. Soit $x\in X\setminus K$. Pour tout $y\in K$,
  il existe des ouverts disjoints $U_y$ et $V_y$ tels que $x\in U_y$ et
  $y\in V_y$. La famille $(V_y)_{y\in K}$ est un recouvrement ouvert de $K$~;
  par compacité, on en extrait un sous-recouvrement fini
  $V_{y_1},\ldots,V_{y_n}$. Posons
  \[
    U = U_{y_1}\cap\cdots\cap U_{y_n}.
  \]
  Alors $U$ est un ouvert contenant $x$. De plus, $U\cap K=\emptyset$~: en
  effet, si $z\in U\cap K$, alors $z\in V_{y_i}$ pour un certain $i$, d'où
  $z\in U_{y_i}\cap V_{y_i}=\emptyset$, contradiction. Ainsi
  $U\subseteq X\setminus K$, ce qui prouve que $X\setminus K$ est ouvert.
\end{proof}

% ============================================================
\section{Espaces réguliers et axiome $T_3$}\label{sec:T3}
% ============================================================

\begin{definition}[Espace régulier]\label{def:regulier}
  Un espace topologique\index{espace régulier}\index{espace T3@espace $T_3$}
  $(X,\tau)$ est dit \emph{régulier} si, pour tout fermé $F\subseteq X$ et
  tout point $x\notin F$, il existe des ouverts disjoints $U$ et $V$ tels que
  $x\in U$ et $F\subseteq V$~:
  \[
    \forall\, F \text{ fermé},\;\forall\, x\notin F,\quad
    \exists\, U,V\in\tau,\quad
    x\in U,\; F\subseteq V,\; U\cap V=\emptyset.
  \]
  Un espace est dit \emph{$T_3$} s'il est régulier et $T_1$.
\end{definition}

\begin{remarque}[Régulier vs.\ $T_3$]\label{rem:rem:regulier-T3}
  La convention peut varier selon les auteurs. Nous suivons la convention de
  Bourbaki~: un espace $T_3$ est régulier \emph{et} $T_1$. La régularité
  seule n'implique pas $T_1$ (un espace à un point muni de la topologie
  grossière est régulier mais cela est trivial).
\end{remarque}

\begin{proposition}[Caractérisation de la régularité]\label{prop:regulier-carac}
  Un espace $T_1$ est régulier (donc $T_3$) si et seulement si, pour tout
  point $x\in X$ et tout ouvert $U$ contenant $x$, il existe un ouvert $V$
  tel que $x\in V\subseteq\overline{V}\subseteq U$.
\end{proposition}

\begin{proof}
  $(\Rightarrow)$ Soit $x\in U$ avec $U$ ouvert. Le complémentaire
  $F=X\setminus U$ est fermé et $x\notin F$. Par régularité, il existe des
  ouverts disjoints $V\ni x$ et $W\supseteq F$. Comme $V\cap W=\emptyset$,
  on a $V\subseteq X\setminus W$. Or $X\setminus W$ est fermé et contenu dans
  $X\setminus F=U$, donc $\overline{V}\subseteq X\setminus W\subseteq U$.

  \medskip
  $(\Leftarrow)$ Soient $F$ fermé et $x\notin F$. Alors $U=X\setminus F$ est
  un ouvert contenant $x$. Par hypothèse, il existe un ouvert $V$ tel que
  $x\in V\subseteq\overline{V}\subseteq U$. Posons $W=X\setminus\overline{V}$.
  Alors $W$ est ouvert, $F\subseteq W$ (car $F\subseteq X\setminus U\subseteq
  X\setminus\overline{V}=W$), et $V\cap W=\emptyset$.
\end{proof}

% ============================================================
\section{Espaces normaux et axiome $T_4$}\label{sec:T4}
% ============================================================

\begin{definition}[Espace normal]\label{def:normal}
  Un espace topologique\index{espace normal}\index{espace T4@espace $T_4$}
  $(X,\tau)$ est dit \emph{normal} si, pour tous fermés disjoints $F_1,F_2
  \subseteq X$, il existe des ouverts disjoints $U_1$ et $U_2$ tels que
  $F_1\subseteq U_1$ et $F_2\subseteq U_2$~:
  \[
    \forall\, F_1,F_2 \text{ fermés},\; F_1\cap F_2=\emptyset,\quad
    \exists\, U_1,U_2\in\tau,\quad
    F_1\subseteq U_1,\; F_2\subseteq U_2,\; U_1\cap U_2=\emptyset.
  \]
  Un espace est dit \emph{$T_4$} s'il est normal et $T_1$.
\end{definition}

% ============================================================
\section{Implications entre axiomes de séparation}\label{sec:implications}
% ============================================================

\begin{theoreme}[Chaîne d'implications]\label{thm:chaine-implications}
  Pour un espace topologique, on a~:
  \[
    T_4 \implies T_3 \implies T_2 \implies T_1 \implies T_0.
  \]
  Aucune de ces implications ne se renverse.
\end{theoreme}

\begin{proof}
  \emph{$T_4\Rightarrow T_3$~:} Soit $(X,\tau)$ un espace $T_4$ (normal et $T_1$).
  Soient $F$ un fermé et $x\notin F$. Comme l'espace est $T_1$, le singleton
  $\{x\}$ est fermé. Les fermés $\{x\}$ et $F$ sont disjoints~; par normalité,
  il existe des ouverts disjoints $U\ni x$ et $V\supseteq F$. Donc l'espace
  est régulier, et comme il est $T_1$, il est $T_3$.

  \medskip
  \emph{$T_3\Rightarrow T_2$~:} Soit $(X,\tau)$ un espace $T_3$ (régulier et $T_1$).
  Soient $x\neq y$. Comme l'espace est $T_1$, $\{y\}$ est fermé et $x\notin\{y\}$.
  Par régularité, il existe des ouverts disjoints $U\ni x$ et $V\ni y$.

  \medskip
  \emph{$T_2\Rightarrow T_1$~:} Évident~: si $x$ et $y$ sont séparés par
  des ouverts disjoints $U\ni x$ et $V\ni y$, alors en particulier $U$ est un
  ouvert contenant $x$ mais pas $y$.

  \medskip
  \emph{$T_1\Rightarrow T_0$~:} Évident.
\end{proof}

\begin{theoreme}[Tout espace compact et Hausdorff est normal]\label{thm:compact-hausdorff-normal}
  Si $(X,\tau)$ est compact et $T_2$\index{compact!Hausdorff implique normal},
  alors il est $T_4$.
\end{theoreme}

\begin{proof}
  L'espace est $T_1$ car il est $T_2$. Montrons la normalité.

  \medskip
  \emph{Étape 1~: compact et point peuvent être séparés (régularité).}
  Soient $K$ un compact (donc fermé par le théorème~\ref{thm:compact-hausdorff-ferme})
  et $x\notin K$. Pour tout $y\in K$, il existe des ouverts disjoints $U_y\ni x$
  et $V_y\ni y$. Par compacité de $K$, on extrait un sous-recouvrement fini
  $V_{y_1},\ldots,V_{y_n}$. Alors $U=\bigcap_{i=1}^n U_{y_i}$ et
  $V=\bigcup_{i=1}^n V_{y_i}$ sont des ouverts disjoints séparant $x$ et $K$.

  \medskip
  \emph{Étape 2~: deux fermés disjoints peuvent être séparés.}
  Soient $F_1,F_2$ deux fermés disjoints. Comme $X$ est compact et les $F_i$
  sont fermés, les $F_i$ sont compacts. Pour chaque $x\in F_1$, l'étape~1
  fournit des ouverts disjoints $U_x\ni x$ et $V_x\supseteq F_2$. Par
  compacité de $F_1$, on extrait $U_{x_1},\ldots,U_{x_m}$ couvrant $F_1$.
  Posons $U=\bigcup_{j=1}^m U_{x_j}$ et $V=\bigcap_{j=1}^m V_{x_j}$. Alors
  $F_1\subseteq U$, $F_2\subseteq V$ et $U\cap V=\emptyset$.
\end{proof}

% ============================================================
\section{Contre-exemples}\label{sec:contre-exemples-separation}
% ============================================================

La construction de contre-exemples est essentielle pour vérifier qu'aucune
implication de la chaîne $T_4\Rightarrow\cdots\Rightarrow T_0$ ne se renverse.

\begin{exemple}[$T_1$ mais pas $T_2$~: topologie cofinie]\label{ex:cofinie-pas-T2}
  Soit $X$ un ensemble infini muni de la topologie cofinie\index{topologie
  cofinie!pas T2@pas $T_2$} $\tau_{\mathrm{cof}}$. Comme montré dans
  l'exemple~\ref{ex:cofinie-T1}, cet espace est $T_1$ mais pas $T_2$~: deux
  ouverts non vides s'intersectent toujours.
\end{exemple}

\begin{exemple}[$T_0$ mais pas $T_1$~: topologie du point
  particulier]\label{ex:point-particulier}
  Soit $X$ un ensemble de cardinal $\geq 2$ et soit $p\in X$ un point
  fixé\index{topologie du point particulier}. La \emph{topologie du point
  particulier} est
  \[
    \tau_p = \{U\subseteq X : p\in U\}\cup\{\emptyset\}.
  \]
  Cet espace est $T_0$~: pour $x\neq y$, si $x=p$ alors $\{p\}$ est ouvert et
  contient $x$ mais peut ne pas contenir $y$~; si ni $x$ ni $y$ n'est $p$,
  l'ouvert $\{p,x\}$ contient $x$ mais pas $y$. Cependant, il n'est pas $T_1$~:
  le singleton $\{p\}$ n'est pas fermé si $|X|\geq 2$, car
  $X\setminus\{p\}$ ne contient pas $p$ et n'est donc pas ouvert.
\end{exemple}

\begin{exemple}[$T_2$ mais pas $T_3$~: la droite à \og double origine\fg{}]\label{ex:double-origine}
  Considérons $X = (\R\setminus\{0\})\cup\{0^+,0^-\}$\index{droite à double
  origine} où $0^+$ et $0^-$ sont deux copies distinctes de l'origine. La
  topologie est engendrée par les ouverts de $\R\setminus\{0\}$ et les
  ensembles de la forme $((-\varepsilon,0)\cup\{0^+\}\cup(0,\varepsilon))$ et
  $((-\varepsilon,0)\cup\{0^-\}\cup(0,\varepsilon))$ pour $\varepsilon>0$.

  Cet espace est $T_2$~: les deux origines sont séparées par
  des ouverts disjoints (par exemple en choisissant des $\varepsilon$
  suffisamment petits, avec des voisinages alternés). Mais il n'est pas
  régulier~: le fermé $\{0^-\}$ et le point $0^+$ ne peuvent pas être séparés
  par des ouverts disjoints, car tout voisinage de $0^+$ et tout voisinage de
  $0^-$ contiennent un intervalle $(-\varepsilon,0)\cup(0,\varepsilon)$ en commun.
\end{exemple}

\begin{exemple}[$T_3$ mais pas $T_4$~: le plan de Sorgenfrey]\label{ex:sorgenfrey}
  La \emph{droite de Sorgenfrey}\index{droite de Sorgenfrey}\index{plan de
  Sorgenfrey} $\R_\ell$ est $\R$ muni de la topologie engendrée par les
  intervalles semi-ouverts $[a,b)$. Elle est $T_3$ (régulière et $T_1$). Le
  \emph{plan de Sorgenfrey} $\R_\ell\times\R_\ell$ est également $T_3$, mais
  il n'est \emph{pas} normal~: l'antidiagonale
  \[
    D = \{(x,-x) : x\in\R\}
  \]
  est un sous-espace fermé discret, et l'on peut montrer (par un argument de
  cardinalité utilisant le lemme de Jones) que deux fermés disjoints non
  dénombrables de $D$ ne peuvent pas être séparés par des ouverts disjoints
  dans $\R_\ell\times\R_\ell$.
\end{exemple}

% ============================================================
\section{Exemples résolus}\label{sec:exemples-separation}
% ============================================================

\begin{exemple}[Tout espace métrique est $T_4$]\label{ex:metrique-T4}
  Soit $(X,d)$ un espace métrique\index{espace métrique!est $T_4$}. Montrons
  qu'il est normal. Soient $F_1$ et $F_2$ deux fermés disjoints. Pour tout
  $x\in X$, posons $d(x,F_i)=\inf\{d(x,y):y\in F_i\}$. Comme $F_1\cap
  F_2=\emptyset$ et les $F_i$ sont fermés, pour $x\in F_1$ on a
  $d(x,F_2)>0$, et réciproquement. Posons
  \[
    U_1 = \bigl\{x\in X : d(x,F_1)<d(x,F_2)\bigr\},\qquad
    U_2 = \bigl\{x\in X : d(x,F_2)<d(x,F_1)\bigr\}.
  \]
  Les fonctions $x\mapsto d(x,F_i)$ sont continues (même $1$-lipschitziennes),
  donc $U_1$ et $U_2$ sont ouverts. De plus $F_1\subseteq U_1$, $F_2\subseteq
  U_2$ et $U_1\cap U_2=\emptyset$. L'espace est $T_1$ car les singletons sont
  fermés dans un espace métrique. Donc $(X,d)$ est $T_4$.
\end{exemple}

\begin{exemple}[Produit d'espaces de Hausdorff]\label{ex:produit-hausdorff}
  Soit $(X_i,\tau_i)_{i\in I}$ une famille d'espaces de Hausdorff\index{espace
  de Hausdorff!produit}. Alors $X=\prod_{i\in I}X_i$ muni de la topologie
  produit est de Hausdorff.

  En effet, soient $x=(x_i)_{i\in I}$ et $y=(y_i)_{i\in I}$ deux points
  distincts de $X$. Il existe $j\in I$ tel que $x_j\neq y_j$. Par l'axiome
  $T_2$ dans $X_j$, il existe des ouverts disjoints $U_j\ni x_j$ et $V_j\ni
  y_j$. Alors $\pi_j^{-1}(U_j)$ et $\pi_j^{-1}(V_j)$ sont des ouverts
  disjoints de $X$ séparant $x$ et $y$.
\end{exemple}

\begin{exemple}[Sous-espace d'un espace régulier]\label{ex:sous-espace-regulier}
  Tout sous-espace d'un espace $T_3$ est $T_3$\index{espace T3@espace $T_3$!sous-espace}.

  Soit $(X,\tau)$ un espace $T_3$ et $A\subseteq X$. L'espace $A$ est $T_1$
  car la propriété $T_1$ est héréditaire. Montrons la régularité. Soient
  $x\in A$ et $F$ un fermé relatif de $A$, avec $x\notin F$. On a
  $F=A\cap G$ pour un fermé $G$ de $X$, et $x\notin G$. Par régularité de
  $X$, il existe des ouverts disjoints $U\ni x$ et $V\supseteq G$ dans $X$.
  Alors $U\cap A$ et $V\cap A$ sont des ouverts disjoints de $A$ avec
  $x\in U\cap A$ et $F\subseteq V\cap A$.
\end{exemple}

% ============================================================
\section{Exercices}\label{sec:exo-separation}
% ============================================================

\begin{exercice}[Topologie grossière et $T_0$]\label{exo:exo:grossiere-T0}%
  {\normalfont\textbf{$\star$}}
  Montrer qu'un espace muni de la topologie grossière n'est $T_0$ que s'il
  contient au plus un point.
\end{exercice}

\begin{exercice}[Hausdorff et graphe fermé]\label{exo:exo:hausdorff-graphe}%
  {\normalfont\textbf{$\star$}}
  Soit $f\colon X\to Y$ une application continue, où $Y$ est de Hausdorff.
  Montrer que le graphe $\Gamma_f=\{(x,f(x)):x\in X\}$ est un fermé de
  $X\times Y$.
\end{exercice}

\begin{exercice}[Topologie discrète]\label{exo:exo:discrete-T4}%
  {\normalfont\textbf{$\star$}}
  Montrer que tout espace muni de la topologie discrète est $T_4$.
\end{exercice}

\begin{exercice}[Limite de suites en $T_1$]\label{exo:exo:limite-T1}%
  {\normalfont\textbf{$\star\star$}}
  Soit $(X,\tau)$ un espace $T_1$. Montrer que si une suite $(x_n)$ converge
  vers $\ell$ et si $x_n=c$ pour une infinité d'indices, alors $c=\ell$.

  \emph{Indication~:} utiliser le fait que les singletons sont fermés.
\end{exercice}

\begin{exercice}[Image continue et séparation]\label{exo:exo:image-continue-T2}%
  {\normalfont\textbf{$\star\star$}}
  \begin{enumerate}
    \item Montrer que si $f\colon X\to Y$ est continue et injective, et si $Y$
      est $T_i$ ($i=0,1,2$), alors $X$ est $T_i$.
    \item En déduire que la propriété d'être $T_2$ est un invariant topologique.
  \end{enumerate}
\end{exercice}

\begin{exercice}[Régularité et bases de voisinages fermés]\label{exo:exo:regulier-base-fermee}%
  {\normalfont\textbf{$\star\star$}}
  Montrer qu'un espace $T_1$ est régulier si et seulement si tout point admet
  une base de voisinages formée de fermés.
\end{exercice}

\begin{exercice}[Normalité et lemme d'Urysohn (aperçu)]\label{exo:exo:urysohn-apercu}%
  {\normalfont\textbf{$\star\star\star$}}
  Soit $(X,\tau)$ un espace $T_4$. Soient $F_0,F_1$ deux fermés disjoints.
  En utilisant la normalité de manière itérée, construire pour tout rationnel
  dyadique $r=k/2^n$ ($0\leq r\leq 1$) un ouvert $U_r$ tel que~:
  \[
    F_0\subseteq U_r,\quad \overline{U_r}\cap F_1=\emptyset,\quad
    r<s\implies \overline{U_r}\subseteq U_s.
  \]
  Puis définir $f(x)=\inf\{r:x\in U_r\}$ et montrer que $f\colon X\to[0,1]$
  est continue avec $f|_{F_0}=0$ et $f|_{F_1}=1$.
\end{exercice}

\begin{exercice}[Compact et Hausdorff~: bijectivité continue]\label{exo:exo:compact-hausdorff-homeo}%
  {\normalfont\textbf{$\star\star\star$}}
  Soit $f\colon X\to Y$ une application continue et bijective, où $X$ est
  compact et $Y$ est de Hausdorff. Montrer que $f$ est un homéomorphisme.

  \emph{Indication~:} montrer que $f$ est fermée en utilisant le
  théorème~\ref{thm:compact-hausdorff-ferme}.
\end{exercice}

% ============================================================
\section*{Résumé du chapitre}\label{sec:resume-separation}
\addcontentsline{toc}{section}{Résumé du chapitre}
% ============================================================

\begin{itemize}
  \item Les axiomes $T_0$ à $T_4$ forment une hiérarchie croissante de
    capacités de séparation~: $T_4\Rightarrow T_3\Rightarrow T_2\Rightarrow
    T_1\Rightarrow T_0$.
  \item Un espace est $T_1$ si et seulement si ses singletons sont fermés.
  \item Un espace de Hausdorff ($T_2$) garantit l'unicité des limites et la
    fermeture des compacts.
  \item $T_2$ est équivalent à la fermeture de la diagonale dans le produit.
  \item Tout espace compact et Hausdorff est normal ($T_4$).
  \item Tout espace métrique est $T_4$.
  \item Les propriétés $T_0$, $T_1$, $T_2$ et $T_3$ sont héréditaires et
    stables par produit. La normalité n'est pas héréditaire en général (le plan
    de Sorgenfrey en est un contre-exemple).
\end{itemize}

% ============================================================================
% CHAPITRE 6 : COMPACITÉ ET THÉORÈME DE TYCHONOFF
% ============================================================================

\chapter{Compacité et théorème de Tychonoff}\label{ch:compacite}

\section*{Introduction}
\addcontentsline{toc}{section}{Introduction}

Le théorème de Heine-Borel\index{théorème de Heine-Borel} affirme qu'un
sous-ensemble de $\R^n$ est compact si et seulement s'il est fermé et borné.
Ce résultat, pilier de l'analyse classique, repose sur la structure
euclidienne. En topologie générale, la notion de compacité se formule
intrinsèquement à l'aide de recouvrements ouverts, sans référence à une
distance ni à un bornage.

\medskip

La compacité est, avec la connexité, l'une des propriétés topologiques les
plus fondamentales. Elle permet d'extraire des informations finies à partir de
données infinies~: convergence de sous-suites, atteinte des bornes d'une
fonction continue, passage de propriétés locales à des propriétés globales.

\medskip

Le résultat culminant de ce chapitre est le \emph{théorème de Tychonoff}~: un
produit arbitraire d'espaces compacts est compact. Ce théorème, dont la
démonstration utilise de manière essentielle le théorème de la sous-base
d'Alexander, est l'un des résultats les plus importants de la topologie
générale. Nous terminons par la \emph{compactification d'Alexandroff}, qui
montre comment <<~rendre compact~>> un espace qui ne l'est pas.

% ============================================================
\section{Définition et premières propriétés}\label{sec:compact-def}
% ============================================================

\begin{definition}[Espace compact]\label{def:compact}
  Un espace topologique\index{compact} $(X,\tau)$ est dit \emph{compact} si
  tout recouvrement ouvert de $X$ admet un sous-recouvrement fini~: pour toute
  famille $(U_i)_{i\in I}$ d'ouverts telle que $X=\bigcup_{i\in I}U_i$, il
  existe une partie finie $J\subseteq I$ telle que $X=\bigcup_{j\in J}U_j$.
\end{definition}

\begin{remarque}[Convention]\label{rem:rem:compact-hausdorff}
  Contrairement à certains auteurs (notamment Bourbaki), nous ne supposons
  \emph{pas} qu'un espace compact est nécessairement séparé. Lorsque nous
  voudrons un espace à la fois compact et de Hausdorff, nous le préciserons
  explicitement.
\end{remarque}

\begin{definition}[Sous-ensemble compact]\label{def:sous-ensemble-compact}
  Un sous-ensemble $K$ d'un espace topologique $(X,\tau)$ est dit
  \emph{compact}\index{compact!sous-ensemble} si le sous-espace $(K,\tau|_K)$
  est compact, c'est-à-dire si tout recouvrement de $K$ par des ouverts de $X$
  admet un sous-recouvrement fini.
\end{definition}

% ============================================================
\section{Propriété de l'intersection finie}\label{sec:FIP}
% ============================================================

\begin{definition}[Propriété de l'intersection finie]\label{def:FIP}
  Une famille $(F_i)_{i\in I}$ de sous-ensembles d'un ensemble $X$ vérifie la
  \emph{propriété de l'intersection finie}\index{propriété de l'intersection
  finie}\index{FIP|see{propriété de l'intersection finie}} (\emph{FIP}, de
  l'anglais \emph{finite intersection property}) si, pour toute partie finie
  $J\subseteq I$, on a $\bigcap_{j\in J}F_j\neq\emptyset$.
\end{definition}

\begin{theoreme}[Caractérisation de la compacité par la FIP]\label{thm:compact-FIP}
  Un espace topologique $(X,\tau)$ est compact si et seulement si toute famille
  de fermés vérifiant la FIP a une intersection non vide\index{compact!caractérisation par FIP}~:
  \[
    \text{pour toute famille } (F_i)_{i\in I} \text{ de fermés avec la FIP},
    \quad \bigcap_{i\in I}F_i\neq\emptyset.
  \]
\end{theoreme}

\begin{proof}
  Il s'agit d'une reformulation par passage au complémentaire.

  \medskip
  $(\Rightarrow)$ Supposons $X$ compact et soit $(F_i)_{i\in I}$ une famille
  de fermés avec la FIP telle que $\bigcap_{i\in I}F_i=\emptyset$. Alors
  $(X\setminus F_i)_{i\in I}$ est un recouvrement ouvert de $X$. Par
  compacité, il existe $J\subseteq I$ fini tel que
  $X=\bigcup_{j\in J}(X\setminus F_j)=X\setminus\bigcap_{j\in J}F_j$, d'où
  $\bigcap_{j\in J}F_j=\emptyset$, contredisant la FIP.

  \medskip
  $(\Leftarrow)$ Soit $(U_i)_{i\in I}$ un recouvrement ouvert de $X$ sans
  sous-recouvrement fini. Alors les fermés $F_i=X\setminus U_i$ vérifient la
  FIP (car aucune sous-famille finie des $U_i$ ne recouvre $X$), mais
  $\bigcap_{i\in I}F_i=X\setminus\bigcup_{i\in I}U_i=\emptyset$,
  contradiction.
\end{proof}

% ============================================================
\section{Compacité et opérations topologiques}\label{sec:compact-operations}
% ============================================================

\begin{theoreme}[Fermé d'un compact]\label{thm:ferme-compact}
  Tout sous-ensemble fermé d'un espace compact est compact\index{compact!fermé dans compact}.
\end{theoreme}

\begin{proof}
  Soit $(X,\tau)$ compact et $F\subseteq X$ fermé. Soit $(U_i)_{i\in I}$ un
  recouvrement de $F$ par des ouverts de $X$. Alors
  $(U_i)_{i\in I}\cup\{X\setminus F\}$ est un recouvrement ouvert de $X$. Par
  compacité, on en extrait un sous-recouvrement fini~:
  \[
    X = U_{i_1}\cup\cdots\cup U_{i_n}\cup(X\setminus F).
  \]
  En particulier, $F\subseteq U_{i_1}\cup\cdots\cup U_{i_n}$.
\end{proof}

\begin{theoreme}[Compact dans un espace de Hausdorff]\label{thm:compact-dans-hausdorff}
  Tout sous-ensemble compact d'un espace de Hausdorff est fermé\index{compact!fermé dans Hausdorff}.
\end{theoreme}

\begin{proof}
  C'est exactement le théorème~\ref{thm:compact-hausdorff-ferme} du
  chapitre~\ref{ch:separation}, dont nous rappelons la preuve. Soit $K$ compact
  dans l'espace de Hausdorff $X$. Pour $x\in X\setminus K$, on construit pour
  chaque $y\in K$ des ouverts disjoints $U_y\ni x$ et $V_y\ni y$. Par
  compacité de $K$, on extrait un sous-recouvrement fini $V_{y_1},\ldots,V_{y_n}$.
  L'ouvert $U=\bigcap_{i=1}^n U_{y_i}$ contient $x$ et est disjoint de $K$.
  Ainsi $X\setminus K$ est ouvert.
\end{proof}

\begin{theoreme}[Image continue d'un compact]\label{thm:image-continue-compact}
  Soit $f\colon X\to Y$ une application continue. Si $X$ est compact, alors
  $f(X)$ est compact\index{compact!image continue}.
\end{theoreme}

\begin{proof}
  Soit $(V_i)_{i\in I}$ un recouvrement ouvert de $f(X)$. Alors
  $(f^{-1}(V_i))_{i\in I}$ est un recouvrement ouvert de $X$ (par continuité
  de $f$). Par compacité de $X$, il existe $J\subseteq I$ fini tel que
  $X=\bigcup_{j\in J}f^{-1}(V_j)$. Donc $f(X)\subseteq\bigcup_{j\in J}V_j$.
\end{proof}

\begin{corollaire}[Théorème des bornes atteintes]\label{cor:cor:bornes-atteintes}
  Soit $f\colon X\to\R$ une application continue, où $X$ est un espace compact
  non vide\index{théorème des bornes atteintes}. Alors $f$ est bornée et
  atteint ses bornes~: il existe $a,b\in X$ tels que
  \[
    f(a) = \inf_{x\in X} f(x),\qquad f(b) = \sup_{x\in X} f(x).
  \]
\end{corollaire}

\begin{proof}
  Par le théorème~\ref{thm:image-continue-compact}, $f(X)$ est une partie
  compacte de $\R$. Comme $\R$ est de Hausdorff, $f(X)$ est fermé
  (théorème~\ref{thm:compact-dans-hausdorff}). Comme $f(X)$ est compact dans
  $\R$, il est borné (tout recouvrement par des intervalles ouverts admet un
  sous-recouvrement fini). Donc $f(X)$ est un fermé borné de $\R$, ce qui
  implique $\sup f(X)\in f(X)$ et $\inf f(X)\in f(X)$.
\end{proof}

% ============================================================
\section{Compacité séquentielle et compacité par point
  d'accumulation}\label{sec:compact-sequentiel}
% ============================================================

\begin{definition}[Compacité séquentielle]\label{def:compact-seq}
  Un espace topologique $(X,\tau)$ est dit \emph{séquentiellement
  compact}\index{compact!séquentiellement} si toute suite $(x_n)_{n\in\N}$
  d'éléments de $X$ admet une sous-suite convergente.
\end{definition}

\begin{definition}[Compacité par point d'accumulation]\label{def:compact-accum}
  Un espace topologique $(X,\tau)$ est dit \emph{compact par point
  d'accumulation}\index{compact!par point d'accumulation} (ou possède la
  \emph{propriété de Bolzano-Weierstrass}\index{Bolzano-Weierstrass}) si tout
  sous-ensemble infini de $X$ admet un point d'accumulation.
\end{definition}

\begin{proposition}[Relations entre les notions de compacité]\label{prop:relations-compacite}
  Pour un espace topologique quelconque~:
  \[
    \text{compact} \implies \text{compact par point d'accumulation}.
  \]
  De plus, dans un espace $T_1$~:
  \[
    \text{séquentiellement compact} \implies \text{compact par point d'accumulation}.
  \]
  En général, aucune de ces implications ne se renverse.
\end{proposition}

\begin{proof}
  \emph{Compact $\Rightarrow$ compact par point d'accumulation.}
  Supposons $X$ compact et soit $A\subseteq X$ infini sans point d'accumulation.
  Alors $A$ est fermé (car $A'\subseteq A$ est vide, où $A'$ désigne l'ensemble
  des points d'accumulation). Pour tout $x\in X$, il existe un ouvert $U_x$
  contenant $x$ tel que $U_x\cap A$ est fini (au plus $\{x\}$ si $x\in A$,
  vide sinon). La famille $(U_x)_{x\in X}$ recouvre $X$~; par compacité, on
  extrait un sous-recouvrement fini $U_{x_1},\ldots,U_{x_n}$. Alors
  $A\subseteq\bigcup_{i=1}^n(U_{x_i}\cap A)$ est une réunion finie
  d'ensembles finis, donc fini, contradiction.

  \medskip
  \emph{Séquentiellement compact $\Rightarrow$ compact par point d'accumulation
  (dans un $T_1$).}
  Soit $A\subseteq X$ infini. On peut extraire une suite $(x_n)$ d'éléments
  deux à deux distincts de $A$. Par compacité séquentielle, il existe une
  sous-suite $(x_{n_k})$ convergeant vers un point $\ell\in X$. Montrons que
  $\ell$ est un point d'accumulation de $A$. Soit $U$ un ouvert contenant
  $\ell$. Il existe $K$ tel que $x_{n_k}\in U$ pour tout $k\geq K$. Comme les
  $x_{n_k}$ sont deux à deux distincts, $(U\setminus\{\ell\})\cap A\neq\emptyset$
  (car au plus un des $x_{n_k}$ peut être égal à $\ell$, puisque l'espace est
  $T_1$ et les $x_{n_k}$ sont distincts).
\end{proof}

\begin{theoreme}[Équivalence dans les espaces métriques]\label{thm:compact-metrique}
  Pour un espace métrique $(X,d)$\index{compact!dans un espace métrique}, les
  conditions suivantes sont équivalentes~:
  \begin{enumerate}
    \item\label{it:cpt1} $X$ est compact.
    \item\label{it:cpt2} $X$ est séquentiellement compact.
    \item\label{it:cpt3} $X$ est complet et totalement borné.
  \end{enumerate}
\end{theoreme}

\begin{proof}
  \emph{\ref{it:cpt1}$\Rightarrow$\ref{it:cpt2}~:} Soit $(x_n)$ une suite
  dans $X$ compact. L'ensemble $A=\{x_n:n\in\N\}$ est, s'il est infini,
  doté d'un point d'accumulation $\ell$ (par la proposition précédente). On
  construit alors par récurrence une sous-suite convergeant vers $\ell$~: ayant
  choisi $x_{n_k}\in B(\ell,1/k)$, on choisit $n_{k+1}>n_k$ tel que
  $x_{n_{k+1}}\in B(\ell,1/(k+1))$, ce qui est possible car tout voisinage de
  $\ell$ contient une infinité de termes de la suite. Si $A$ est fini, une
  valeur est prise une infinité de fois et fournit une sous-suite constante.

  \medskip
  \emph{\ref{it:cpt2}$\Rightarrow$\ref{it:cpt3}~:}

  \emph{Complétude.} Soit $(x_n)$ une suite de Cauchy. Par compacité
  séquentielle, elle admet une sous-suite $(x_{n_k})$ convergeant vers
  $\ell\in X$. Une suite de Cauchy admettant une sous-suite convergente
  converge elle-même (vers $\ell$).

  \emph{Totale bornitude.} Supposons que $X$ n'est pas totalement borné~: il
  existe $\varepsilon>0$ tel qu'aucun ensemble fini de boules ouvertes de rayon
  $\varepsilon$ ne recouvre $X$. On construit par récurrence une suite
  $(x_n)$~: $x_0\in X$ arbitraire, et $x_{n+1}\in X\setminus\bigcup_{i=0}^n
  B(x_i,\varepsilon)$. Alors $d(x_m,x_n)\geq\varepsilon$ pour $m\neq n$,
  donc aucune sous-suite n'est de Cauchy, a fortiori aucune n'est convergente,
  contredisant la compacité séquentielle.

  \medskip
  \emph{\ref{it:cpt3}$\Rightarrow$\ref{it:cpt1}~:} Supposons $X$ complet et
  totalement borné. Soit $(U_i)_{i\in I}$ un recouvrement ouvert sans
  sous-recouvrement fini. On construit une suite de Cauchy sans sous-suite
  convergente, contradiction.

  Pour $\varepsilon=1$, la totale bornitude fournit un recouvrement fini par
  des boules $B(a_1^{(0)},1),\ldots,B(a_{k_0}^{(0)},1)$. L'une d'elles,
  disons $B_0=B(a^{(0)},1)$, ne peut être recouverte par un nombre fini des
  $U_i$. On itère~: $B_0\cap X$ est totalement borné (car sous-ensemble d'un
  espace totalement borné), et on le recouvre par des boules de rayon $1/2$.
  L'une d'elles, $B_1$, n'est pas recouvrable par un nombre fini des $U_i$.
  Par récurrence, on obtient des boules $B_n$ de rayon $1/2^n$ avec
  $\overline{B_{n+1}}\subseteq \overline{B_n}$, et les centres forment une
  suite de Cauchy. Par complétude, cette suite converge vers $\ell\in X$. Alors
  $\ell\in U_{i_0}$ pour un certain $i_0$, et il existe $N$ tel que
  $B_N\subseteq U_{i_0}$, ce qui contredit le fait que $B_N$ ne peut être
  recouvert par un nombre fini des $U_i$.
\end{proof}

% ============================================================
\section{Théorème de la sous-base d'Alexander}\label{sec:alexander}
% ============================================================

\begin{definition}[Sous-base]\label{def:sous-base-rappel}
  Soit $(X,\tau)$ un espace topologique\index{sous-base}. Une famille
  $\mathcal{S}\subseteq\tau$ est une \emph{sous-base} de $\tau$ si les
  intersections finies d'éléments de $\mathcal{S}$ forment une base de $\tau$.
\end{definition}

\begin{theoreme}[Théorème de la sous-base d'Alexander]\label{thm:alexander}
  Un espace topologique $(X,\tau)$ est compact si et seulement si tout
  recouvrement de $X$ par des éléments d'une sous-base $\mathcal{S}$
  de $\tau$ admet un sous-recouvrement fini\index{théorème d'Alexander}.
\end{theoreme}

\begin{proof}
  Le sens direct est évident (un recouvrement par des éléments de $\mathcal{S}$
  est un cas particulier de recouvrement ouvert).

  \medskip
  Pour la réciproque, supposons que $X$ n'est pas compact. Parmi tous les
  recouvrements ouverts de $X$ sans sous-recouvrement fini, considérons
  l'ensemble $\Sigma$ de tels recouvrements, ordonné par inclusion. Nous allons
  appliquer le lemme de Zorn pour obtenir un recouvrement maximal.

  Soit $(\mathcal{U}_\alpha)_{\alpha\in A}$ une chaîne dans $\Sigma$. La
  réunion $\mathcal{U}=\bigcup_{\alpha\in A}\mathcal{U}_\alpha$ est un
  recouvrement ouvert (car chaque $\mathcal{U}_\alpha$ recouvre $X$) qui n'a
  pas de sous-recouvrement fini (sinon, par finitude, ce sous-recouvrement
  serait contenu dans un $\mathcal{U}_\alpha$, contradiction). Donc
  $\mathcal{U}\in\Sigma$, et par le lemme de Zorn, $\Sigma$ admet un élément
  maximal $\mathcal{U}_0$.

  \medskip
  \emph{Propriétés de $\mathcal{U}_0$~:}
  \begin{enumerate}
    \item $\mathcal{U}_0$ recouvre $X$ et n'admet pas de sous-recouvrement fini.
    \item Par maximalité, pour tout ouvert $W\notin\mathcal{U}_0$,
      $\mathcal{U}_0\cup\{W\}$ admet un sous-recouvrement fini, c'est-à-dire
      qu'il existe $U_1,\ldots,U_n\in\mathcal{U}_0$ tels que
      $X=W\cup U_1\cup\cdots\cup U_n$.
  \end{enumerate}

  Montrons que $\mathcal{U}_0\cap\mathcal{S}$ recouvre $X$, ce qui fournira
  la contradiction souhaitée (car $\mathcal{U}_0\cap\mathcal{S}$ serait un
  recouvrement par des éléments de $\mathcal{S}$ sans sous-recouvrement fini).

  Soit $x\in X$. Il existe $U\in\mathcal{U}_0$ avec $x\in U$. Comme
  $\mathcal{S}$ est une sous-base, il existe $S_1,\ldots,S_m\in\mathcal{S}$
  tels que $x\in S_1\cap\cdots\cap S_m\subseteq U$. Affirmons qu'au moins un
  $S_k$ appartient à $\mathcal{U}_0$. Sinon, pour chaque $k$,
  $S_k\notin\mathcal{U}_0$, et par maximalité, il existe des
  $U_1^{(k)},\ldots,U_{n_k}^{(k)}\in\mathcal{U}_0$ tels que
  $X=S_k\cup U_1^{(k)}\cup\cdots\cup U_{n_k}^{(k)}$. Alors
  \[
    X \setminus \bigl(S_1\cap\cdots\cap S_m\bigr)
    = (X\setminus S_1)\cup\cdots\cup(X\setminus S_m)
    \subseteq \bigcup_{k=1}^m\bigcup_{j=1}^{n_k} U_j^{(k)}.
  \]
  Donc $X \subseteq (S_1\cap\cdots\cap S_m) \cup \bigcup_{k,j}U_j^{(k)}
  \subseteq U \cup \bigcup_{k,j}U_j^{(k)}$, fournissant un sous-recouvrement
  fini de $\mathcal{U}_0$ (car $U\in\mathcal{U}_0$ et les $U_j^{(k)}\in
  \mathcal{U}_0$), contradiction.

  Donc il existe $S_k\in\mathcal{U}_0\cap\mathcal{S}$ avec $x\in S_k$,
  prouvant que $\mathcal{U}_0\cap\mathcal{S}$ recouvre $X$.
\end{proof}

% ============================================================
\section{Théorème de Tychonoff}\label{sec:tychonoff}
% ============================================================

\begin{theoreme}[Théorème de Tychonoff]\label{thm:tychonoff}
  Soit $(X_i)_{i\in I}$ une famille d'espaces topologiques\index{théorème de
  Tychonoff}. Le produit $X=\prod_{i\in I}X_i$, muni de la topologie produit,
  est compact si et seulement si chaque $X_i$ est compact.
\end{theoreme}

\begin{proof}
  \emph{Sens direct ($\Rightarrow$).} Si $X$ est compact, alors chaque
  projection $\pi_i\colon X\to X_i$ est continue et surjective (car les $X_i$
  sont non vides), donc $X_i=\pi_i(X)$ est compact comme image continue d'un
  compact (théorème~\ref{thm:image-continue-compact}).

  \medskip
  \emph{Sens réciproque ($\Leftarrow$).} Supposons chaque $X_i$ compact.
  Rappelons que la topologie produit est engendrée par la sous-base
  \[
    \mathcal{S} = \bigl\{\pi_i^{-1}(U_i) : i\in I,\; U_i\in\tau_i\bigr\}.
  \]
  Par le théorème d'Alexander~(\ref{thm:alexander}), il suffit de montrer que
  tout recouvrement de $X$ par des éléments de $\mathcal{S}$ admet un
  sous-recouvrement fini.

  Soit $(\pi_{i_\alpha}^{-1}(U_{i_\alpha}))_{\alpha\in A}$ un recouvrement de
  $X$ par des éléments de $\mathcal{S}$. Supposons par l'absurde qu'il n'admet
  pas de sous-recouvrement fini. Nous allons montrer que le recouvrement n'est
  pas surjectif, i.e., qu'il ne recouvre pas $X$.

  Pour chaque $i\in I$, considérons l'ensemble
  \[
    A_i = \{\alpha\in A : i_\alpha = i\}.
  \]
  La famille $(U_{i_\alpha})_{\alpha\in A_i}$ est une famille d'ouverts de
  $X_i$. Affirmons que pour au moins un indice $i$, cette famille ne recouvre
  pas $X_i$. En effet, si pour tout $i$, $(U_{i_\alpha})_{\alpha\in A_i}$
  recouvrait $X_i$, alors par compacité de $X_i$, on pourrait en extraire un
  sous-recouvrement fini, et la réunion de ces sous-recouvrements finis (portant
  sur un nombre fini d'indices $i$ effectivement utilisés dans le recouvrement)
  fournirait un sous-recouvrement fini de $X$, contradiction.

  Plus précisément, notons $I_0=\{i\in I:\exists\,\alpha\in A,\;i_\alpha=i\}$
  l'ensemble des indices effectivement utilisés. Si pour tout $i\in I_0$, les
  $(U_{i_\alpha})_{\alpha\in A_i}$ recouvraient $X_i$, par compacité on
  extrairait des sous-recouvrements finis. La réunion de ces sous-recouvrements
  finis (un pour chaque $i\in I_0$) donnerait un recouvrement fini de $X$ par
  des éléments de $\mathcal{S}$... mais $I_0$ pourrait être infini.

  Reformulons l'argument. Supposons que pour tout $i\in I$, la famille
  $(U_{i_\alpha})_{\alpha\in A_i}$ recouvre $X_i$. Si $A_i=\emptyset$, cela
  est vacuement vrai (ce cas ne pose pas de problème car $\pi_i^{-1}(X_i)=X$
  n'intervient pas). Considérons $i\in I$ fixé. Si $(U_{i_\alpha})_{\alpha\in
  A_i}$ recouvre $X_i$ et $X_i$ est compact, il existe un sous-ensemble fini
  $B_i\subseteq A_i$ tel que $X_i=\bigcup_{\alpha\in B_i}U_{i_\alpha}$. Alors
  $X=\pi_i^{-1}(X_i)=\bigcup_{\alpha\in B_i}\pi_i^{-1}(U_{i_\alpha})$, ce
  qui est un sous-recouvrement fini, contradiction.

  Donc il existe $i_0\in I$ tel que $(U_{i_\alpha})_{\alpha\in A_{i_0}}$ ne
  recouvre pas $X_{i_0}$. Si $A_{i_0}=\emptyset$, c'est immédiat. Sinon, il
  existe $y_{i_0}\in X_{i_0}\setminus\bigcup_{\alpha\in A_{i_0}}U_{i_\alpha}$.
  Par l'axiome du choix, pour $i\neq i_0$, choisissons $y_i\in X_i$ quelconque.
  Le point $y=(y_i)_{i\in I}\in X$ n'appartient à aucun
  $\pi_{i_\alpha}^{-1}(U_{i_\alpha})$ avec $i_\alpha=i_0$ (car
  $y_{i_0}\notin U_{i_\alpha}$).

  Mais $y$ pourrait appartenir à un $\pi_{i_\alpha}^{-1}(U_{i_\alpha})$ avec
  $i_\alpha\neq i_0$. Nous devons raffiner~: il faut choisir les $y_i$ de
  manière à éviter \emph{tous} les éléments du recouvrement.

  Reprenons. Pour tout $i\in I_0$, si $(U_{i_\alpha})_{\alpha\in A_i}$
  recouvre $X_i$, on obtient un sous-recouvrement fini de $X$, contradiction.
  Donc pour tout $i\in I_0$, il existe
  $y_i\in X_i\setminus\bigcup_{\alpha\in A_i}U_{i_\alpha}$. Pour $i\notin
  I_0$, choisissons $y_i\in X_i$ arbitraire. Le point $y=(y_i)_{i\in I}$
  satisfait~: pour tout $\alpha\in A$, on a $i_\alpha\in I_0$ et
  $y_{i_\alpha}\notin U_{i_\alpha}$, donc $y\notin\pi_{i_\alpha}^{-1}
  (U_{i_\alpha})$. Ainsi $y$ n'est recouvert par aucun élément, contredisant
  le fait que la famille est un recouvrement. Cette contradiction achève la
  preuve.
\end{proof}

\begin{remarque}[Rôle de l'axiome du choix]\label{rem:rem:tychonoff-AC}
  La preuve ci-dessus utilise le lemme de Zorn (dans le théorème d'Alexander)
  et l'axiome du choix (pour construire le point $y$). En fait, le théorème de
  Tychonoff est \emph{équivalent} à l'axiome du choix dans la théorie des
  ensembles ZF\index{axiome du choix}.
\end{remarque}

% ============================================================
\section{Compactification}\label{sec:compactification}
% ============================================================

\begin{definition}[Compactification]\label{def:compactification}
  Une \emph{compactification}\index{compactification} d'un espace topologique
  $(X,\tau)$ est un couple $(Y,\varphi)$ où $Y$ est un espace compact et
  $\varphi\colon X\to Y$ est un plongement (homéomorphisme sur son image) tel
  que $\varphi(X)$ est dense dans $Y$.
\end{definition}

\begin{definition}[Compactification d'Alexandroff]\label{def:alexandroff}
  Soit $(X,\tau)$ un espace topologique non compact\index{compactification!d'Alexandroff}\index{Alexandroff}.
  La \emph{compactification d'Alexandroff} (ou \emph{compactification par un
  point}) est l'espace $X^*=X\cup\{\infty\}$, où $\infty\notin X$, muni de la
  topologie
  \[
    \tau^* = \tau \cup \{(X\setminus K)\cup\{\infty\} : K\subseteq X
    \text{ compact et fermé}\}.
  \]
  Les ouverts de $X^*$ sont les ouverts de $X$ et les ensembles de la forme
  $(X\setminus K)\cup\{\infty\}$ où $K$ est un sous-ensemble compact fermé
  de $X$.
\end{definition}

\begin{theoreme}[Propriétés de la compactification d'Alexandroff]\label{thm:alexandroff-proprietes}
  Soit $(X,\tau)$ un espace topologique.
  \begin{enumerate}
    \item\label{it:alex1} $\tau^*$ est bien une topologie sur $X^*$.
    \item\label{it:alex2} $X^*$ est compact.
    \item\label{it:alex3} $X$ est un sous-espace ouvert dense de $X^*$.
    \item\label{it:alex4} $X^*$ est de Hausdorff si et seulement si $X$ est
      localement compact et de Hausdorff.
  \end{enumerate}
\end{theoreme}

\begin{proof}
  \emph{\ref{it:alex1}~:} Vérifions les axiomes d'une topologie.
  $\emptyset\in\tau\subseteq\tau^*$ et $X^*=(X\setminus\emptyset)\cup\{\infty\}
  \in\tau^*$ (car $\emptyset$ est compact et fermé).

  \emph{Stabilité par intersection finie~:} L'intersection de deux ouverts de
  $\tau$ est dans $\tau$. L'intersection d'un ouvert $U\in\tau$ et d'un ouvert
  $(X\setminus K)\cup\{\infty\}$ est $U\cap(X\setminus K)\in\tau$ si
  $\infty\notin U$ (ce qui est le cas). L'intersection de $(X\setminus
  K_1)\cup\{\infty\}$ et $(X\setminus K_2)\cup\{\infty\}$ est
  $(X\setminus(K_1\cup K_2))\cup\{\infty\}\in\tau^*$ (car $K_1\cup K_2$ est
  compact et fermé).

  \emph{Stabilité par réunion quelconque~:} analogue.

  \medskip
  \emph{\ref{it:alex2}~:} Soit $(W_\alpha)_{\alpha\in A}$ un recouvrement
  ouvert de $X^*$. Il existe $\alpha_0$ tel que $\infty\in W_{\alpha_0}$, donc
  $W_{\alpha_0}=(X\setminus K)\cup\{\infty\}$ pour un compact fermé $K$. La
  famille $(W_\alpha\cap X)_{\alpha\neq\alpha_0}$ est un recouvrement ouvert de
  $K$. Par compacité de $K$, on en extrait un sous-recouvrement fini
  $W_{\alpha_1}\cap X,\ldots,W_{\alpha_n}\cap X$. Alors
  $W_{\alpha_0},W_{\alpha_1},\ldots,W_{\alpha_n}$ recouvrent $X^*$.

  \medskip
  \emph{\ref{it:alex3}~:} La topologie induite sur $X$ par $\tau^*$ coïncide
  avec $\tau$. De plus, $X$ est dense~: tout ouvert non vide de $X^*$ contenant
  $\infty$ est de la forme $(X\setminus K)\cup\{\infty\}$ avec $K$ compact
  fermé et $K\neq X$ (car $X$ n'est pas compact), donc rencontre $X$.

  \medskip
  \emph{\ref{it:alex4}~:} ($\Rightarrow$) Si $X^*$ est $T_2$, alors $X$,
  comme sous-espace d'un Hausdorff, est $T_2$. Pour la compacité locale, soit
  $x\in X$. Comme $X^*$ est $T_2$, il existe des ouverts disjoints $U\ni x$ et
  $V\ni\infty$ dans $X^*$. On a $V=(X\setminus K)\cup\{\infty\}$ pour un
  compact fermé $K$, et $U\cap V=\emptyset$ implique $U\subseteq K$. Comme
  $K$ est compact et $\overline{U}^X\subseteq K$, le point $x$ admet un
  voisinage compact $\overline{U}^X$.

  ($\Leftarrow$) Si $X$ est localement compact et $T_2$, la séparation de deux
  points de $X$ est assurée par $T_2$. Pour séparer $x\in X$ de $\infty$~: par
  compacité locale, il existe un ouvert $U$ et un compact $K$ tels que
  $x\in U\subseteq K$. Comme $X$ est $T_2$, $K$ est fermé. Alors $U$ et
  $(X\setminus K)\cup\{\infty\}$ sont des ouverts disjoints séparant $x$ et
  $\infty$.
\end{proof}

\begin{center}
\begin{tikzpicture}[scale=1.0]
  % The real line
  \draw[thick, blue!70!black] (-4,0) -- (4,0);
  \draw[thick, blue!70!black, -{Stealth}] (-4,0) -- (-4.5,0);
  \draw[thick, blue!70!black, -{Stealth}] (4,0) -- (4.5,0);
  \node[below, font=\small] at (0,-0.2) {$\R$};

  % Arrow
  \draw[-{Stealth}, thick, gray] (0,-0.8) -- (0,-1.8)
    node[midway, right, font=\small] {compactification};

  % The circle (one-point compactification)
  \draw[thick, red!70!black] (0,-3.5) circle (1.5);
  \fill[red!70!black] (0,-2) circle (2pt);
  \node[above right, font=\small] at (0.1,-2) {$\infty$};
  \node[below, font=\small] at (0,-5.2) {$\R^*\simeq S^1$};

  % Mark a point on the circle
  \fill[blue!70!black] (-1.5,-3.5) circle (2pt);
  \node[left, font=\small] at (-1.7,-3.5) {$0$};
\end{tikzpicture}
\captionof{figure}{La compactification d'Alexandroff de $\R$ est
  homéomorphe au cercle $S^1$.}
\label{fig:compactification-R}
\end{center}

\begin{exemple}[Compactification d'Alexandroff de $\R^n$]\label{ex:alexandroff-Rn}
  La compactification d'Alexandroff\index{compactification!de $\R^n$} de
  $\R^n$ est homéomorphe à la sphère $S^n$. En particulier, $\R^*\simeq S^1$
  et $(\R^2)^*\simeq S^2$.
\end{exemple}

\begin{exemple}[Compactification de $\N$]\label{ex:alexandroff-N}
  Considérons $\N$ muni de la topologie discrète\index{compactification!de
  $\N$}. Sa compactification d'Alexandroff est $\N^*=\N\cup\{\infty\}$, où les
  ouverts contenant $\infty$ sont les complémentaires d'ensembles finis. Cet
  espace est homéomorphe au sous-espace $\{0\}\cup\{1/n:n\geq 1\}$ de $\R$
  (avec la topologie induite).
\end{exemple}

% ============================================================
\section{Contre-exemples}\label{sec:contre-exemples-compacite}
% ============================================================

\begin{exemple}[$[0,1)$ n'est pas compact]\label{ex:intervalle-ouvert}
  L'intervalle $[0,1)$\index{compact!contre-exemples} n'est pas compact~: le
  recouvrement $\bigl([0,1-1/n)\bigr)_{n\geq 2}$ est un recouvrement ouvert
  sans sous-recouvrement fini, car toute sous-famille finie a pour réunion
  $[0,1-1/N)$ pour un certain $N$, qui ne contient pas les points de
  $[1-1/N,1)$.
\end{exemple}

\begin{exemple}[{$\Q\cap[0,1]$ n'est pas compact}]\label{ex:Q-pas-compact}
  Le sous-espace $\Q\cap[0,1]$ de $\R$ n'est pas compact. En effet,
  choisissons un irrationnel $\alpha\in(0,1)$ (par exemple $\alpha=1/\sqrt{2}$)
  et considérons le recouvrement
  \[
    \bigl\{[0,\alpha-1/n)\cap\Q : n\geq 1\bigr\}
    \cup \bigl\{(\alpha+1/n,1]\cap\Q : n\geq 1\bigr\}.
  \]
  Ce recouvrement ouvert (dans la topologie induite) n'admet pas de
  sous-recouvrement fini, car toute sous-famille finie laisse non couvert un
  voisinage rationnel de $\alpha$.
\end{exemple}

\begin{exemple}[Produit $\{0,1\}^{\N}$ est compact]\label{ex:cantor-compact}
  L'espace $\{0,1\}$ muni de la topologie discrète est compact (il est fini).
  Par le théorème de Tychonoff\index{espace de Cantor}, le produit
  $\{0,1\}^{\N}$ est compact. Cet espace est homéomorphe à l'ensemble de
  Cantor $C\subseteq[0,1]$, fournissant un compact métrisable non
  dénombrable totalement discontinu.
\end{exemple}

% ============================================================
\section{Exemples résolus}\label{sec:exemples-compacite}
% ============================================================

\begin{exemple}[{Compacité de $[0,1]$}]\label{ex:compact-01}
  Montrons directement que $[0,1]$ est compact\index{compact!de $[0,1]$}. Soit
  $(U_i)_{i\in I}$ un recouvrement ouvert de $[0,1]$. Posons
  \[
    s = \sup\bigl\{t\in[0,1] : [0,t] \text{ est recouvert par un nombre fini
    des } U_i\bigr\}.
  \]
  Comme $0$ appartient à un certain $U_{i_0}$, on a $s>0$. Montrons que $s=1$
  et que $[0,s]$ est finiment recouvert.

  Il existe $j\in I$ tel que $s\in U_j$. Comme $U_j$ est ouvert, il contient
  un intervalle $(s-\varepsilon,s+\varepsilon)\cap[0,1]$ pour un certain
  $\varepsilon>0$. Par définition de la borne supérieure, il existe $t\in
  (s-\varepsilon,s)$ tel que $[0,t]$ est recouvert par $U_{i_1},\ldots,U_{i_n}$.
  Alors $[0,\min(s+\varepsilon,1)]$ est recouvert par $U_{i_1},\ldots,U_{i_n},
  U_j$. Si $s<1$, cela contredit la définition de $s$ (car $\min(s+\varepsilon,
  1)>s$). Donc $s=1$ et $[0,1]$ est finiment recouvert.
\end{exemple}

\begin{exemple}[Application du théorème de Tychonoff~: le cube de
  Hilbert]\label{ex:cube-hilbert}
  Le \emph{cube de Hilbert}\index{cube de Hilbert} est l'espace
  $[0,1]^{\N}=\prod_{n=0}^{\infty}[0,1]$ muni de la topologie produit.
  Chaque facteur $[0,1]$ étant compact, le théorème de Tychonoff assure que
  $[0,1]^{\N}$ est compact. Cet espace est métrisable (par la distance
  $d(x,y)=\sum_{n=0}^{\infty}|x_n-y_n|/2^n$) et a la propriété remarquable
  d'être un \emph{espace universel}~: tout espace métrique séparable se plonge
  dans $[0,1]^{\N}$.
\end{exemple}

\begin{exemple}[Compacité et suites de fonctions]\label{ex:compact-fonctions}
  Soit $K$ un espace compact et $(f_n)_{n\in\N}$ une suite décroissante de
  fonctions continues $f_n\colon K\to\R$ convergeant simplement vers $0$.
  Montrons que la convergence est uniforme (\emph{théorème de
  Dini}\index{théorème de Dini}).

  Soit $\varepsilon>0$. Pour chaque $n$, posons $U_n=\{x\in K:f_n(x)<
  \varepsilon\}$. Par continuité, $U_n$ est ouvert. Comme $(f_n)$ est
  décroissante, $U_n\subseteq U_{n+1}$. Comme $f_n\to 0$ simplement,
  $\bigcup_{n\in\N}U_n=K$. Par compacité de $K$, il existe $N$ tel que
  $K=U_N$, c'est-à-dire $f_N(x)<\varepsilon$ pour tout $x\in K$. Comme $(f_n)$
  est décroissante et positive, $f_n(x)<\varepsilon$ pour tout $n\geq N$ et
  tout $x\in K$.
\end{exemple}

\begin{exemple}[Recouvrement ouvert et nombre de Lebesgue]\label{ex:lebesgue}
  Soit $(X,d)$ un espace métrique compact et $(U_i)_{i\in I}$ un recouvrement
  ouvert\index{nombre de Lebesgue}. Il existe $\delta>0$ (\emph{nombre de
  Lebesgue du recouvrement}) tel que tout sous-ensemble de diamètre
  $<\delta$ est contenu dans un $U_i$.

  Supposons le contraire~: pour tout $n\geq 1$, il existe un ensemble $A_n$ de
  diamètre $<1/n$ non contenu dans aucun $U_i$. Choisissons $x_n\in A_n$. Par
  compacité séquentielle, $(x_n)$ admet une sous-suite $(x_{n_k})$ convergeant
  vers $\ell\in X$. Il existe $i_0\in I$ et $r>0$ tels que
  $B(\ell,r)\subseteq U_{i_0}$. Pour $k$ assez grand, $x_{n_k}\in B(\ell,r/2)$
  et $\operatorname{diam}(A_{n_k})<r/2$, donc $A_{n_k}\subseteq B(\ell,r)
  \subseteq U_{i_0}$, contradiction.
\end{exemple}

% ============================================================
\section{Exercices}\label{sec:exo-compacite}
% ============================================================

\begin{exercice}[Ensembles compacts de $\R$]\label{exo:exo:compacts-R}%
  {\normalfont\textbf{$\star$}}
  Montrer que les sous-ensembles compacts de $\R$ (muni de la topologie
  usuelle) sont exactement les fermés bornés.
\end{exercice}

\begin{exercice}[Image réciproque et compacité]\label{exo:exo:image-reciproque}%
  {\normalfont\textbf{$\star$}}
  Soit $f\colon X\to Y$ continue, avec $Y$ compact. L'image réciproque $f^{-1}(K)$
  d'un compact $K\subseteq Y$ est-elle nécessairement compacte ? Justifier.
\end{exercice}

\begin{exercice}[Union et intersection de compacts]\label{exo:exo:union-intersection}%
  {\normalfont\textbf{$\star$}}
  Soit $(X,\tau)$ un espace topologique.
  \begin{enumerate}
    \item Montrer que la réunion finie de sous-ensembles compacts est compacte.
    \item Si $X$ est de Hausdorff, montrer que l'intersection quelconque de
      sous-ensembles compacts est compacte.
    \item Donner un exemple montrant que l'intersection de compacts n'est pas
      nécessairement compacte si $X$ n'est pas de Hausdorff.
  \end{enumerate}
\end{exercice}

\begin{exercice}[Compacité et topologie cofinie]\label{exo:exo:cofinie-compact}%
  {\normalfont\textbf{$\star$}}
  Montrer que tout espace muni de la topologie cofinie est compact.
\end{exercice}

\begin{exercice}[Tube lemma]\label{exo:exo:tube-lemma}%
  {\normalfont\textbf{$\star\star$}}
  Soient $X,Y$ des espaces topologiques avec $Y$ compact. Soient $x_0\in X$ et
  $W$ un ouvert de $X\times Y$ contenant $\{x_0\}\times Y$. Montrer qu'il
  existe un ouvert $U$ de $X$ contenant $x_0$ tel que $U\times Y\subseteq W$.
\end{exercice}

\begin{exercice}[Compacité du spectre d'un anneau]\label{exo:exo:spectre}%
  {\normalfont\textbf{$\star\star$}}
  Soit $A$ un anneau commutatif unitaire. Le \emph{spectre} $\operatorname{Spec}(A)$
  est l'ensemble des idéaux premiers de $A$, muni de la topologie de Zariski.
  Montrer que $\operatorname{Spec}(A)$ est compact (quasi-compact dans la
  terminologie de certains auteurs).

  \emph{Indication~:} utiliser la caractérisation par la FIP.
\end{exercice}

\begin{exercice}[Tychonoff et ultrafiltres]\label{exo:exo:tychonoff-ultrafiltres}%
  {\normalfont\textbf{$\star\star\star$}}
  Donner une preuve alternative du théorème de Tychonoff en utilisant les
  ultrafiltres.

  \emph{Indication~:} montrer qu'un espace est compact si et seulement si tout
  ultrafiltre converge. Puis, pour un produit, montrer que l'image d'un
  ultrafiltre par chaque projection converge.
\end{exercice}

\begin{exercice}[Compactification de Čech-Stone (aperçu)]\label{exo:exo:cech-stone}%
  {\normalfont\textbf{$\star\star\star$}}
  Soit $X$ un espace complètement régulier ($T_{3\frac{1}{2}}$). Notons
  $C_b(X)$ l'ensemble des fonctions continues bornées $f\colon X\to\R$. Pour
  chaque $f\in C_b(X)$, notons $I_f=[\ \inf f,\sup f\ ]$. Considérons le
  plongement
  \[
    e\colon X\to\prod_{f\in C_b(X)}I_f,\qquad e(x)=\bigl(f(x)\bigr)_{f\in C_b(X)}.
  \]
  \begin{enumerate}
    \item Montrer que $e$ est un plongement (injectif, continu, et ouvert sur
      son image).
    \item Le produit $\prod_{f}I_f$ est compact par Tychonoff. La
      compactification de \v{C}ech-Stone est $\beta X=\overline{e(X)}$.
      Montrer que $\beta X$ est un espace compact de Hausdorff.
    \item Montrer que toute fonction continue bornée $g\colon X\to\R$ admet un
      prolongement continu $\bar{g}\colon\beta X\to\R$.
  \end{enumerate}
\end{exercice}

\begin{exercice}[Compacité locale et prolongement]\label{exo:exo:locale-prolongement}%
  {\normalfont\textbf{$\star\star$}}
  Soit $X$ un espace localement compact et de Hausdorff, et soit $X^*$ sa
  compactification d'Alexandroff. Montrer que toute application continue
  $f\colon X\to Y$, où $Y$ est de Hausdorff, telle que pour tout compact
  $K\subseteq Y$, $f^{-1}(K)$ est compact, se prolonge en une application
  continue $f^*\colon X^*\to Y^*$ (compactification d'Alexandroff de $Y$).
\end{exercice}

% ============================================================
\section*{Résumé du chapitre}\label{sec:resume-compacite}
\addcontentsline{toc}{section}{Résumé du chapitre}
% ============================================================

\begin{itemize}
  \item Un espace est \emph{compact} si tout recouvrement ouvert admet un
    sous-recouvrement fini, de manière équivalente si toute famille de fermés
    avec la FIP a une intersection non vide.
  \item Tout fermé d'un compact est compact ; tout compact d'un Hausdorff est
    fermé.
  \item L'image continue d'un compact est compacte, d'où le théorème des bornes
    atteintes.
  \item Dans un espace métrique~: compact $\Leftrightarrow$ séquentiellement
    compact $\Leftrightarrow$ complet et totalement borné.
  \item Le \emph{théorème d'Alexander}~: la compacité se teste sur une
    sous-base.
  \item Le \emph{théorème de Tychonoff}~: un produit (arbitraire) d'espaces
    compacts est compact. Ce résultat est équivalent à l'axiome du choix.
  \item La \emph{compactification d'Alexandroff} ajoute un point à un espace
    non compact pour le rendre compact. Elle est de Hausdorff si et seulement
    si l'espace de départ est localement compact et de Hausdorff.
  \item La compacité est le cadre naturel de nombreux résultats fondamentaux~:
    théorème de Dini, lemme du tube, existence d'un nombre de Lebesgue, etc.
\end{itemize}
%=============================================================================
% CHAPITRE 7 : Connexité et Connexité par Arcs
%=============================================================================

\chapter{Connexité et connexité par arcs}
\label{chap:connexite}
\index{connexité}

\medskip
\noindent
Les sous-ensembles connexes de~$\mathbb{R}$ sont exactement les intervalles : c'est un
résultat fondamental de l'analyse réelle qui repose sur la propriété de la borne
supérieure. Comment étendre cette notion aux espaces topologiques généraux ?
La connexité traduit l'idée intuitive qu'un espace est formé « d'un seul
morceau » ; la connexité par arcs en donne un renforcement naturel, en exigeant
que deux points quelconques puissent être joints par un chemin continu.
Ce chapitre développe ces deux notions, établit leurs relations, et fournit
les contre-exemples classiques qui en délimitent les frontières.

%-----------------------------------------------------------------------------
\section{Espaces connexes}
\label{sec:espaces-connexes}
\index{espace connexe}

\begin{definition}[Espace connexe]
\label{def:connexe}
\index{connexe!espace}
Un espace topologique $(X,\mathcal{T})$ est dit \emph{connexe} s'il n'existe pas de
partition de~$X$ en deux ouverts non vides disjoints, c'est-à-dire s'il n'existe
pas $U,V \in \mathcal{T}$ tels que
\[
U \neq \varnothing, \quad V \neq \varnothing, \quad U \cap V = \varnothing, \quad
U \cup V = X.
\]
Une partie $A \subset X$ est \emph{connexe} si le sous-espace $(A,\mathcal{T}_A)$ est connexe.
\end{definition}

\begin{proposition}[Caractérisations équivalentes de la connexité]
\label{prop:equiv-connexe}
\index{connexe!caractérisations}
Soit $(X,\mathcal{T})$ un espace topologique. Les assertions suivantes sont équivalentes :
\begin{enumerate}[label=(\roman*)]
  \item $X$ est connexe.
  \item Les seules parties de~$X$ qui sont à la fois ouvertes et fermées (dites
        \emph{ouvertes-fermées}\index{ouvert-fermé} ou \emph{clopen}\index{clopen})
        sont $\varnothing$ et~$X$.
  \item Toute application continue $f \colon X \to \{0,1\}$ (muni de la topologie discrète)
        est constante.
  \item Pour toute partition $X = A \cup B$ avec $A \cap \overline{B} = \varnothing$
        et $\overline{A} \cap B = \varnothing$, on a $A = \varnothing$ ou $B = \varnothing$.
\end{enumerate}
\end{proposition}

\begin{proof}
\noindent\textbf{(i) $\Leftrightarrow$ (ii).}
Si $A \subset X$ est ouvert-fermé avec $A \neq \varnothing$ et $A \neq X$, alors
$U = A$ et $V = X \setminus A$ forment une partition en deux ouverts non vides, donc
$X$ n'est pas connexe. Réciproquement, si $X = U \cup V$ est une partition en
ouverts non vides disjoints, alors $U$ est ouvert et fermé (car $U = X \setminus V$),
avec $U \neq \varnothing$ et $U \neq X$.

\noindent\textbf{(ii) $\Leftrightarrow$ (iii).}
Soit $f \colon X \to \{0,1\}$ continue. Alors $f^{-1}(\{0\})$ est ouvert-fermé dans~$X$.
Par~(ii), $f^{-1}(\{0\}) = \varnothing$ ou $f^{-1}(\{0\}) = X$, donc $f$ est constante.
Réciproquement, si $A$ est ouvert-fermé non trivial, la fonction indicatrice
$\mathbf{1}_{A}$ est continue non constante vers~$\{0,1\}$.

\noindent\textbf{(i) $\Leftrightarrow$ (iv).}
Si $X = A \cup B$ est une partition avec $A \cap \overline{B} = \varnothing$ et
$\overline{A} \cap B = \varnothing$, alors $A$ est ouvert (car $A = X \setminus \overline{B}$
n'est possible que si $\overline{B} \subset B$, d'où $B$ fermé, donc $A$ ouvert) et
de même $B$ est ouvert. On se ramène à~(i). La réciproque est immédiate puisque
des ouverts disjoints vérifient les conditions de séparation.
\end{proof}

\begin{remarque}
\label{rem:connexe-non-vide}
Par convention, l'espace vide~$\varnothing$ n'est pas connexe dans la tradition
de Bourbaki, bien que certains auteurs l'admettent. Nous adoptons ici la convention
que les espaces connexes sont non vides.
\end{remarque}

%-----------------------------------------------------------------------------
\section{Image continue d'un connexe}
\label{sec:image-connexe}

\begin{theoreme}[Image continue d'un connexe]
\label{thm:image-connexe}
\index{connexe!image continue}
Soit $f \colon X \to Y$ une application continue entre espaces topologiques.
Si $X$ est connexe, alors $f(X)$ est connexe dans~$Y$.
\end{theoreme}

\begin{proof}
Supposons par l'absurde que $f(X)$ ne soit pas connexe. Il existe alors
$U, V$ ouverts dans $f(X)$ (pour la topologie induite) tels que
\[
U \neq \varnothing, \quad V \neq \varnothing, \quad U \cap V = \varnothing, \quad
U \cup V = f(X).
\]
Par continuité de~$f$, les ensembles $f^{-1}(U)$ et $f^{-1}(V)$ sont ouverts
dans~$X$. De plus :
\begin{itemize}
  \item $f^{-1}(U) \neq \varnothing$ et $f^{-1}(V) \neq \varnothing$
        (car $U \cap f(X) \neq \varnothing$ et $V \cap f(X) \neq \varnothing$) ;
  \item $f^{-1}(U) \cap f^{-1}(V) = f^{-1}(U \cap V) = f^{-1}(\varnothing)
        = \varnothing$ ;
  \item $f^{-1}(U) \cup f^{-1}(V) = f^{-1}(U \cup V) = f^{-1}(f(X)) = X$.
\end{itemize}
On obtient une partition de~$X$ en deux ouverts non vides, ce qui contredit
la connexité de~$X$.
\end{proof}

\begin{corollaire}[Théorème des valeurs intermédiaires]
\label{cor:valeurs-intermediaires}
\index{théorème!des valeurs intermédiaires}
Soit $f \colon [a,b] \to \mathbb{R}$ une application continue. Pour tout
$\gamma$ compris entre $f(a)$ et $f(b)$, il existe $c \in [a,b]$ tel que $f(c) = \gamma$.
\end{corollaire}

\begin{proof}
L'intervalle $[a,b]$ est connexe dans~$\mathbb{R}$ (car c'est un intervalle).
Par le théorème~\ref{thm:image-connexe}, $f([a,b])$ est une partie connexe
de~$\mathbb{R}$, donc un intervalle (les connexes de~$\mathbb{R}$ sont les
intervalles). Comme $f(a)$ et $f(b)$ appartiennent à $f([a,b])$ et que $\gamma$
est compris entre eux, on a $\gamma \in f([a,b])$.
\end{proof}

\begin{proposition}
\label{prop:connexe-reunion}
\index{connexe!réunion}
Soit $(A_i)_{i \in I}$ une famille de parties connexes d'un espace topologique~$X$
telle que $\bigcap_{i \in I} A_i \neq \varnothing$. Alors $\bigcup_{i \in I} A_i$
est connexe.
\end{proposition}

\begin{proof}
Soit $x_0 \in \bigcap_{i \in I} A_i$ et soit $f \colon \bigcup_{i \in I} A_i \to \{0,1\}$
continue. Pour tout $i \in I$, la restriction $f|_{A_i}$ est continue sur le
connexe~$A_i$, donc constante (proposition~\ref{prop:equiv-connexe}). Comme
$x_0 \in A_i$, on a $f|_{A_i} \equiv f(x_0)$ pour tout~$i$, d'où $f$ est
constante sur $\bigcup_{i \in I} A_i$.
\end{proof}

\begin{proposition}
\label{prop:adherence-connexe}
\index{connexe!adhérence}
Si $A$ est une partie connexe de~$X$, alors $\overline{A}$ est connexe.
Plus généralement, si $A \subset B \subset \overline{A}$, alors $B$ est connexe.
\end{proposition}

\begin{proof}
Soit $f \colon \overline{A} \to \{0,1\}$ continue. Alors $f|_A$ est continue sur
le connexe~$A$, donc constante égale à une valeur $c \in \{0,1\}$. Comme $A$ est
dense dans~$\overline{A}$ et $\{c\}$ est fermé dans~$\{0,1\}$, on a
$f^{-1}(\{c\}) \supset \overline{A}$, donc $f \equiv c$ sur~$\overline{A}$.
Le même raisonnement s'applique à toute partie~$B$ avec $A \subset B \subset \overline{A}$.
\end{proof}

%-----------------------------------------------------------------------------
\section{Composantes connexes}
\label{sec:composantes-connexes}
\index{composante connexe}

\begin{definition}[Composante connexe]
\label{def:composante-connexe}
Soit $X$ un espace topologique et $x \in X$. La \emph{composante connexe}
de~$x$ dans~$X$ est la plus grande partie connexe de~$X$ contenant~$x$ :
\[
C(x) = \bigcup \{ A \subset X : A \text{ connexe}, \, x \in A \}.
\]
\end{definition}

\begin{proposition}
\label{prop:partition-composantes}
Les composantes connexes forment une partition de~$X$ en fermés connexes.
\end{proposition}

\begin{proof}
Par la proposition~\ref{prop:connexe-reunion}, $C(x)$ est connexe (réunion de
connexes ayant le point~$x$ en commun). Si $C(x) \cap C(y) \neq \varnothing$,
alors $C(x) \cup C(y)$ est connexe, donc $C(x) \cup C(y) \subset C(x)$ par
maximalité, d'où $C(x) = C(y)$. Les composantes forment donc une partition.
Par la proposition~\ref{prop:adherence-connexe}, $\overline{C(x)}$ est connexe et
contient~$x$, donc $\overline{C(x)} \subset C(x)$ par maximalité : chaque
composante est fermée.
\end{proof}

\begin{remarque}
\label{rem:composante-non-ouverte}
Les composantes connexes ne sont pas nécessairement ouvertes. Par exemple, dans
$\mathbb{Q}$ muni de la topologie induite par~$\mathbb{R}$, chaque composante
connexe est un singleton $\{q\}$, qui n'est pas ouvert.
\end{remarque}

\begin{definition}[Espace totalement discontinu]
\label{def:totalement-discontinu}
\index{totalement discontinu}
Un espace topologique~$X$ est dit \emph{totalement discontinu} si ses composantes
connexes sont réduites à des singletons, c'est-à-dire si les seules parties
connexes de~$X$ sont les singletons et l'ensemble vide.
\end{definition}

\begin{exemple}
\label{ex:Q-totalement-discontinu}
\index{rationnels!totalement discontinu}
L'ensemble~$\mathbb{Q}$ des nombres rationnels, muni de la topologie induite
par~$\mathbb{R}$, est totalement discontinu. En effet, si $A \subset \mathbb{Q}$
contient deux points distincts $p < q$, il existe un irrationnel
$\alpha \in {]p,q[}$, et alors
$A = (A \cap {]{-\infty},\alpha[}) \cup (A \cap {]\alpha,{+\infty}[})$
est une partition en deux ouverts non vides de~$A$.
\end{exemple}

\begin{exemple}[Ensemble de Cantor]
\label{ex:cantor-totalement-discontinu}
\index{Cantor!ensemble de}
L'ensemble triadique de Cantor $\mathcal{C} \subset [0,1]$, défini comme
l'ensemble des réels de $[0,1]$ admettant un développement en base~$3$ sans le
chiffre~$1$, est un compact totalement discontinu sans point isolé.
C'est même un espace parfait\index{espace parfait} : fermé et sans point isolé.
\end{exemple}

%-----------------------------------------------------------------------------
\section{Connexité par arcs}
\label{sec:connexite-arcs}
\index{connexité par arcs}

\begin{definition}[Chemin, espace connexe par arcs]
\label{def:connexe-par-arcs}
\index{chemin}
\index{connexe par arcs}
Soit $X$ un espace topologique.
\begin{enumerate}[label=(\roman*)]
  \item Un \emph{chemin} dans~$X$ d'origine~$a$ et d'extrémité~$b$ est une
        application continue $\gamma \colon [0,1] \to X$ telle que
        $\gamma(0) = a$ et $\gamma(1) = b$.
  \item L'espace~$X$ est dit \emph{connexe par arcs} si, pour tous $a, b \in X$,
        il existe un chemin de~$a$ à~$b$ dans~$X$.
\end{enumerate}
\end{definition}

\begin{theoreme}[Connexe par arcs implique connexe]
\label{thm:arcs-implique-connexe}
\index{connexe par arcs!implique connexe}
Tout espace connexe par arcs est connexe.
\end{theoreme}

\begin{proof}
Soit $X$ connexe par arcs et supposons par l'absurde que $X = U \cup V$ avec
$U, V$ ouverts non vides disjoints. Soient $a \in U$ et $b \in V$. Par hypothèse,
il existe un chemin $\gamma \colon [0,1] \to X$ avec $\gamma(0) = a$ et
$\gamma(1) = b$. Alors
\[
[0,1] = \gamma^{-1}(U) \cup \gamma^{-1}(V)
\]
est une partition de $[0,1]$ en deux ouverts non vides ($0 \in \gamma^{-1}(U)$
et $1 \in \gamma^{-1}(V)$). Ceci contredit la connexité de~$[0,1]$.
\end{proof}

\begin{remarque}
\label{rem:reciproque-fausse}
La réciproque est fausse : il existe des espaces connexes qui ne sont pas connexes
par arcs. Le contre-exemple classique est la courbe du topologiste
(exemple~\ref{ex:courbe-topologiste}).
\end{remarque}

%-----------------------------------------------------------------------------
\section{Contre-exemple : la courbe du topologiste}
\label{sec:courbe-topologiste}
\index{courbe du topologiste}
\index{sinus du topologiste|see{courbe du topologiste}}

\begin{exemple}[Courbe du topologiste]
\label{ex:courbe-topologiste}
Considérons dans~$\mathbb{R}^2$ l'ensemble
\[
S = \Bigl\{ \Bigl(x, \sin\frac{1}{x}\Bigr) : x \in {]0,1]} \Bigr\}
\]
et sa fermeture
\[
\overline{S} = S \cup (\{0\} \times [-1,1]).
\]
L'ensemble~$\overline{S}$ est la \emph{courbe du topologiste} (ou \emph{sinus du
topologiste fermé}).

\begin{center}
\begin{tikzpicture}[scale=3]
  % Axes
  \draw[->] (-0.2,0) -- (1.3,0) node[right] {$x$};
  \draw[->] (0,-1.3) -- (0,1.3) node[above] {$y$};
  % Segment vertical {0} x [-1,1]
  \draw[blue, very thick] (0,-1) -- (0,1);
  \fill[blue] (0,-1) circle (0.5pt);
  \fill[blue] (0,1) circle (0.5pt);
  % Courbe sin(1/x) - échantillonnage dense près de 0
  \draw[red, thick, domain=0.01:1, samples=2000, smooth]
    plot (\x, {sin(1/\x r)});
  % Labels
  \node[left, blue] at (0,1) {$(0,1)$};
  \node[left, blue] at (0,-1) {$(0,-1)$};
  \node[below right] at (1,0) {$1$};
  \node[below left] at (0,0) {$0$};
  % Légende
  \node[red, right] at (0.6, -1.1) {$y = \sin(1/x)$};
  \node[blue, left] at (-0.05, 0) {$\{0\}\times[-1,1]$};
\end{tikzpicture}
\end{center}

\noindent\textbf{$\overline{S}$ est connexe.}
L'ensemble~$S$ est l'image de l'intervalle ${]0,1]}$ par l'application continue
$x \mapsto (x, \sin(1/x))$, donc $S$ est connexe
(théorème~\ref{thm:image-connexe}). Par la
proposition~\ref{prop:adherence-connexe}, $\overline{S}$ est connexe.

\noindent\textbf{$\overline{S}$ n'est pas connexe par arcs.}
Montrons qu'il n'existe pas de chemin de $(0,0)$ à $(1, \sin 1)$ dans~$\overline{S}$.
Supposons par l'absurde qu'un tel chemin
$\gamma = (\gamma_1, \gamma_2) \colon [0,1] \to \overline{S}$ existe, avec
$\gamma(0) = (0,0)$. Posons $t_0 = \sup \{ t \in [0,1] : \gamma_1(t) = 0 \}$.
Par continuité, $\gamma_1(t_0) = 0$.
Pour $t > t_0$ proche de~$t_0$, on a $\gamma_1(t) > 0$, donc
$\gamma_2(t) = \sin(1/\gamma_1(t))$. Or, lorsque $t \to t_0^+$, $\gamma_1(t) \to 0^+$
et $\sin(1/\gamma_1(t))$ oscille entre $-1$ et~$1$ sans converger, ce qui contredit
la continuité de~$\gamma_2$ en~$t_0$.
\end{exemple}

%-----------------------------------------------------------------------------
\section{Connexité locale}
\label{sec:connexite-locale}
\index{connexité locale}

\begin{definition}[Localement connexe, localement connexe par arcs]
\label{def:localement-connexe}
\index{localement connexe}
\index{localement connexe par arcs}
Un espace topologique~$X$ est :
\begin{enumerate}[label=(\roman*)]
  \item \emph{localement connexe} si, pour tout $x \in X$ et tout ouvert $U$
        contenant~$x$, il existe un ouvert connexe $V$ tel que $x \in V \subset U$ ;
  \item \emph{localement connexe par arcs} si, pour tout $x \in X$ et tout ouvert $U$
        contenant~$x$, il existe un ouvert connexe par arcs $V$ tel que
        $x \in V \subset U$.
\end{enumerate}
\end{definition}

\begin{proposition}
\label{prop:loc-connexe-composantes-ouvertes}
Un espace~$X$ est localement connexe si et seulement si les composantes connexes
de tout ouvert de~$X$ sont ouvertes.
\end{proposition}

\begin{proof}
\emph{Condition suffisante.} Soit $x \in X$ et $U$ un ouvert contenant~$x$. La
composante connexe de~$x$ dans~$U$ est un ouvert connexe contenant~$x$ et contenu
dans~$U$.

\emph{Condition nécessaire.} Soit $U$ ouvert dans~$X$ et $C$ une composante connexe
de~$U$. Pour tout $x \in C$, il existe un ouvert connexe $V_x$ avec
$x \in V_x \subset U$. Comme $V_x$ est connexe et rencontre~$C$ (en~$x$), on a
$V_x \subset C$. Donc $C = \bigcup_{x \in C} V_x$ est ouvert.
\end{proof}

\begin{theoreme}[Connexe + localement connexe par arcs $\Rightarrow$ connexe par arcs]
\label{thm:connexe-loc-arcs}
\index{connexe par arcs!et localement connexe par arcs}
Soit $X$ un espace topologique connexe et localement connexe par arcs.
Alors $X$ est connexe par arcs.
\end{theoreme}

\begin{proof}
Fixons $x_0 \in X$ et posons
\[
A = \{ x \in X : \text{il existe un chemin de } x_0 \text{ à } x \text{ dans } X \}.
\]
Montrons que $A$ est ouvert-fermé non vide, ce qui entraînera $A = X$ par connexité.

\noindent\emph{$A$ est non vide.} Le chemin constant $\gamma(t) = x_0$ montre
$x_0 \in A$.

\noindent\emph{$A$ est ouvert.} Soit $a \in A$. Comme $X$ est localement connexe
par arcs, il existe un ouvert connexe par arcs $V$ contenant~$a$. Pour tout
$v \in V$, il existe un chemin de~$a$ à~$v$ dans~$V$, et un chemin de~$x_0$ à~$a$
dans~$X$. En concaténant ces chemins, on obtient un chemin de~$x_0$ à~$v$, donc
$v \in A$. Ainsi $V \subset A$.

\noindent\emph{$A$ est fermé.} Soit $b \in X \setminus A$. Il existe un ouvert
connexe par arcs $V$ contenant~$b$. Si $V \cap A \neq \varnothing$, soit $a \in V \cap A$.
Alors il existe un chemin de~$x_0$ à~$a$ et un chemin de~$a$ à~$b$ dans~$V$, d'où
un chemin de~$x_0$ à~$b$ : contradiction avec $b \notin A$. Donc
$V \subset X \setminus A$, et $X \setminus A$ est ouvert.

Par connexité de~$X$ et le fait que $A$ est ouvert-fermé non vide, $A = X$.
\end{proof}

%-----------------------------------------------------------------------------
\section{Produit d'espaces connexes}
\label{sec:produit-connexes}

\begin{theoreme}[Produit fini de connexes]
\label{thm:produit-connexes}
\index{connexe!produit}
Soit $(X_i)_{1 \leqslant i \leqslant n}$ une famille finie d'espaces topologiques
non vides. Alors $\prod_{i=1}^{n} X_i$ (muni de la topologie produit) est connexe
si et seulement si chaque~$X_i$ est connexe.
\end{theoreme}

\begin{proof}
\emph{Condition nécessaire.} Chaque projection
$\pi_i \colon \prod_{j} X_j \to X_i$ est continue et surjective. Si le produit est
connexe, son image~$X_i$ l'est par le théorème~\ref{thm:image-connexe}.

\emph{Condition suffisante.} Il suffit de traiter le cas $n = 2$ et de procéder par
récurrence. Soient $X$ et~$Y$ connexes. Fixons un point $(a,b) \in X \times Y$.
Pour tout $y \in Y$, l'ensemble $X \times \{y\}$ est homéomorphe à~$X$, donc
connexe, et passe par~$(a,y)$. L'ensemble $\{a\} \times Y$ est homéomorphe à~$Y$,
donc connexe. Alors
\[
C_y = (\{a\} \times Y) \cup (X \times \{y\})
\]
est connexe (réunion de deux connexes ayant $(a,y)$ en commun). Finalement,
\[
X \times Y = \bigcup_{y \in Y} C_y
\]
est connexe comme réunion de connexes ayant le point $(a,b)$ en commun
(car $(a,b) \in C_y$ pour tout~$y$, puisque $(a,b) \in \{a\} \times Y \subset C_y$).
\end{proof}

\begin{theoreme}[Produit arbitraire de connexes]
\label{thm:produit-arbitraire-connexes}
\index{connexe!produit arbitraire}
Soit $(X_i)_{i \in I}$ une famille quelconque d'espaces topologiques non vides.
Alors $\prod_{i \in I} X_i$ (muni de la topologie produit) est connexe si et
seulement si chaque~$X_i$ est connexe.
\end{theoreme}

\begin{proof}
La condition nécessaire découle des projections comme ci-dessus. Pour la condition
suffisante, fixons $a = (a_i)_{i \in I} \in \prod X_i$. Pour toute partie finie
$J \subset I$, l'ensemble
\[
E_J = \Bigl\{ x \in \prod_{i \in I} X_i : x_i = a_i \text{ pour } i \notin J \Bigr\}
\]
est homéomorphe à $\prod_{j \in J} X_j$, qui est connexe (cas fini). De plus,
$a \in E_J$ pour tout~$J$, donc $A = \bigcup_J E_J$ est connexe. Montrons que
$A$ est dense dans $\prod X_i$. Un ouvert de base est de la forme
$\prod_{i \in I} U_i$ avec $U_i = X_i$ sauf pour un nombre fini d'indices
$J_0$. On peut choisir $x \in A$ avec $x_j \in U_j$ pour $j \in J_0$ et $x_i = a_i$
sinon, ce qui montre $A \cap \prod U_i \neq \varnothing$. Donc $\overline{A} = \prod X_i$
et le résultat découle de la proposition~\ref{prop:adherence-connexe}.
\end{proof}

%-----------------------------------------------------------------------------
\section{Exemples et illustrations}
\label{sec:exemples-connexite}

\begin{center}
\begin{tikzpicture}[scale=0.8]
  % Espace connexe
  \begin{scope}[shift={(-4,0)}]
    \draw[thick, fill=blue!10] plot[smooth cycle, tension=0.7]
      coordinates {(0,1.5) (1.2,1) (1.5,0) (1,-1) (-0.5,-1.2) (-1.5,-0.3) (-1,1)};
    \node at (0,-1.8) {Connexe};
  \end{scope}
  % Espace non connexe
  \begin{scope}[shift={(3,0)}]
    \draw[thick, fill=red!10] plot[smooth cycle, tension=0.7]
      coordinates {(-2,0.8) (-1.2,1) (-0.8,0) (-1.5,-0.8) (-2.2,0)};
    \draw[thick, fill=red!10] plot[smooth cycle, tension=0.7]
      coordinates {(0.8,0.8) (1.8,0.5) (1.5,-0.5) (0.5,-0.6) (0.3,0.2)};
    \node at (0,-1.8) {Non connexe};
  \end{scope}
\end{tikzpicture}
\end{center}

\begin{exemple}
\label{ex:connexes-R}
\index{connexe!dans $\mathbb{R}$}
\begin{enumerate}[label=(\alph*)]
  \item Les parties connexes de~$\mathbb{R}$ sont exactement les intervalles
        (y compris $\varnothing$, les singletons, et $\mathbb{R}$ entier).
  \item $\mathbb{R}^n$ est connexe par arcs pour tout $n \geqslant 1$ : pour tous
        $a, b \in \mathbb{R}^n$, le segment $\gamma(t) = (1-t)a + tb$ est un chemin
        de~$a$ à~$b$.
  \item La sphère $S^n = \{x \in \mathbb{R}^{n+1} : \|x\| = 1\}$ est connexe
        par arcs pour $n \geqslant 1$.
\end{enumerate}
\end{exemple}

\begin{exemple}
\label{ex:GL-connexe}
\index{groupe linéaire!connexité}
Le groupe $\mathrm{GL}_n(\mathbb{R})$ n'est pas connexe pour $n \geqslant 1$ :
l'application $\det \colon \mathrm{GL}_n(\mathbb{R}) \to \mathbb{R}^*$ est continue
et surjective, et $\mathbb{R}^* = {]{-\infty},0[} \cup {]0,{+\infty}[}$ n'est pas
connexe. Les deux composantes connexes sont $\mathrm{GL}_n^+(\mathbb{R})$
(matrices de déterminant positif) et $\mathrm{GL}_n^-(\mathbb{R})$ (déterminant négatif).
\end{exemple}

\begin{exemple}
\label{ex:connexes-divers}
\begin{enumerate}[label=(\alph*)]
  \item Tout espace muni de la topologie grossière est connexe.
  \item Un espace discret à plus d'un point est totalement discontinu.
  \item L'ensemble $\{(x,y) \in \mathbb{R}^2 : xy \neq 0\}$ a quatre composantes
        connexes, correspondant aux quatre quadrants ouverts.
\end{enumerate}
\end{exemple}

%-----------------------------------------------------------------------------
\section{Exercices}
\label{sec:exo-connexite}

\begin{exercice}[$\star$]
\label{exo:connexe-image}
Soit $f \colon \mathbb{R} \to \mathbb{R}$ continue. Montrer que le graphe
$\Gamma_f = \{(x, f(x)) : x \in \mathbb{R}\}$ est connexe par arcs dans~$\mathbb{R}^2$.
\end{exercice}

\begin{exercice}[$\star$]
\label{exo:connexe-fermeture}
Soit $A$ une partie connexe d'un espace topologique~$X$ et $B$ tel que
$A \subset B \subset \overline{A}$. Montrer que $B$ est connexe.
\end{exercice}

\begin{exercice}[$\star$]
\label{exo:connexe-intersection}
Donner un exemple de deux parties connexes $A, B$ de~$\mathbb{R}^2$ telles que
$A \cap B$ ne soit pas connexe.
\end{exercice}

\begin{exercice}[$\star\star$]
\label{exo:composante-connexe-fermee}
Soit $X$ un espace topologique et $C$ une composante connexe de~$X$.
\begin{enumerate}[label=(\alph*)]
  \item Montrer que $C$ est fermée.
  \item Donner un exemple où $C$ n'est pas ouverte.
  \item Montrer que si $X$ est localement connexe, alors $C$ est ouverte.
\end{enumerate}
\end{exercice}

\begin{exercice}[$\star\star$]
\label{exo:R-prive-point}
Montrer que $\mathbb{R}$ privé d'un point n'est pas connexe, mais que
$\mathbb{R}^n$ privé d'un point est connexe pour $n \geqslant 2$.
\end{exercice}

\begin{exercice}[$\star\star$]
\label{exo:connexe-quotient}
Soit $X$ un espace connexe et $\sim$ une relation d'équivalence sur~$X$.
Montrer que $X/{\sim}$ est connexe.
\end{exercice}

\begin{exercice}[$\star\star$]
\label{exo:cantor-totalement-discontinu}
Montrer que l'ensemble de Cantor~$\mathcal{C}$ est totalement discontinu.
\emph{Indication :} montrer que pour tous $x \neq y$ dans~$\mathcal{C}$, il existe
un ouvert-fermé de~$\mathcal{C}$ contenant~$x$ mais pas~$y$.
\end{exercice}

\begin{exercice}[$\star\star\star$]
\label{exo:connexe-par-arcs-ouvert}
Soit $U$ un ouvert connexe de~$\mathbb{R}^n$. Montrer que $U$ est connexe par arcs.
\emph{Indication :} $\mathbb{R}^n$ est localement connexe par arcs.
\end{exercice}

\begin{exercice}[$\star\star\star$]
\label{exo:peigne}
\index{peigne (espace)}
On considère le \emph{peigne} $P \subset \mathbb{R}^2$ défini par
\[
P = ([0,1] \times \{0\}) \cup \Bigl( \bigcup_{n=1}^{\infty} \bigl\{\tfrac{1}{n}\bigr\} \times [0,1] \Bigr) \cup (\{0\} \times [0,1]).
\]
\begin{enumerate}[label=(\alph*)]
  \item Montrer que $P$ est connexe par arcs.
  \item Montrer que $P$ n'est pas localement connexe en $(0,1)$.
\end{enumerate}
\end{exercice}

\begin{exercice}[$\star\star\star$]
\label{exo:topologiste-variante}
Soit $T = \{(x, \sin(1/x)) : x > 0\} \cup \{(0,0)\} \subset \mathbb{R}^2$.
\begin{enumerate}[label=(\alph*)]
  \item Montrer que $T$ est connexe.
  \item Montrer que $T$ n'est pas connexe par arcs.
  \item Comparer avec $\overline{S}$ de l'exemple~\ref{ex:courbe-topologiste}.
\end{enumerate}
\end{exercice}

%-----------------------------------------------------------------------------
\section*{Résumé du chapitre~\thechapter}
\addcontentsline{toc}{section}{Résumé du chapitre~\thechapter}

\begin{itemize}
  \item Un espace est \textbf{connexe} s'il ne peut être partitionné en deux ouverts
        non vides ; de manière équivalente, ses seuls ouverts-fermés sont
        $\varnothing$ et~$X$.
  \item L'\textbf{image continue} d'un espace connexe est connexe, d'où le
        théorème des valeurs intermédiaires.
  \item Les \textbf{composantes connexes} partitionnent $X$ en fermés connexes
        maximaux.
  \item Un espace est \textbf{connexe par arcs} si deux points quelconques sont
        reliés par un chemin continu. Connexe par arcs $\Rightarrow$ connexe, mais
        pas réciproquement (courbe du topologiste).
  \item \textbf{Connexe + localement connexe par arcs} $\Rightarrow$ connexe par arcs.
  \item Le \textbf{produit} (topologie produit) d'espaces connexes est connexe.
  \item $\mathbb{Q}$ et l'ensemble de Cantor sont \textbf{totalement discontinus}.
\end{itemize}


%=============================================================================
% CHAPITRE 8 : Espaces Produits et Topologie Quotient
%=============================================================================

\chapter{Espaces produits et topologie quotient}
\label{chap:produit-quotient}
\index{topologie produit}
\index{topologie quotient}

\medskip
\noindent
La construction de nouveaux espaces topologiques à partir d'espaces donnés est
l'une des opérations fondamentales de la topologie. Ce chapitre présente deux
constructions essentielles : la \emph{topologie produit}, qui permet de munir un
produit cartésien d'une topologie naturelle, et la \emph{topologie quotient}, qui
permet d'identifier des points pour obtenir de nouveaux espaces. Ces outils sont
indispensables : le tore, la bouteille de Klein, l'espace projectif et bien
d'autres surfaces classiques sont définis comme des espaces quotients.

%-----------------------------------------------------------------------------
\section{Topologie produit : produits finis}
\label{sec:topologie-produit-fini}
\index{topologie produit!produit fini}

\begin{definition}[Topologie produit, cas fini]
\label{def:topologie-produit-finie}
Soient $(X_1, \mathcal{T}_1), \ldots, (X_n, \mathcal{T}_n)$ des espaces topologiques.
La \emph{topologie produit} sur $X = \prod_{i=1}^n X_i$ est la topologie engendrée
par la base
\[
\mathcal{B} = \Bigl\{ \prod_{i=1}^n U_i : U_i \in \mathcal{T}_i \text{ pour tout } i \Bigr\}.
\]
On note $\pi_i \colon X \to X_i$ la \emph{projection canonique}\index{projection canonique}
sur le $i$-ème facteur, définie par $\pi_i(x_1, \ldots, x_n) = x_i$.
\end{definition}

\begin{remarque}
\label{rem:produit-fini-base}
Pour un produit fini, la famille ci-dessus est bien une base de topologie (elle est
stable par intersections finies). La topologie produit coïncide avec la topologie
engendrée par les ensembles $\pi_i^{-1}(U_i)$ pour $U_i \in \mathcal{T}_i$ :
c'est la topologie initiale pour la famille $(\pi_i)_{1 \leqslant i \leqslant n}$.
\end{remarque}

%-----------------------------------------------------------------------------
\section{Topologie produit : cas général}
\label{sec:topologie-produit-general}
\index{topologie produit!cas général}
\index{topologie boîte}

\begin{definition}[Topologie produit, cas général]
\label{def:topologie-produit-generale}
Soit $(X_i, \mathcal{T}_i)_{i \in I}$ une famille d'espaces topologiques.
La \emph{topologie produit} (ou \emph{topologie de Tychonoff}\index{Tychonoff!topologie de})
sur $X = \prod_{i \in I} X_i$ est la topologie engendrée par la sous-base
\[
\mathcal{S} = \bigl\{ \pi_i^{-1}(U_i) : i \in I, \, U_i \in \mathcal{T}_i \bigr\}.
\]
Les ouverts de base sont les ensembles de la forme $\prod_{i \in I} U_i$ où
$U_i \in \mathcal{T}_i$ pour tout~$i$ et $U_i = X_i$ sauf pour un \emph{nombre
fini} d'indices.
\end{definition}

\begin{definition}[Topologie boîte]
\label{def:topologie-boite}
\index{topologie boîte}
La \emph{topologie boîte} sur $\prod_{i \in I} X_i$ est la topologie engendrée par la base
\[
\mathcal{B}_{\mathrm{box}} = \Bigl\{ \prod_{i \in I} U_i : U_i \in \mathcal{T}_i
\text{ pour tout } i \Bigr\}
\]
où \emph{tous} les $U_i$ sont des ouverts quelconques (sans restriction de finitude).
\end{definition}

\begin{remarque}
\label{rem:boite-vs-produit}
Pour un produit fini, les topologies boîte et produit coïncident. Pour un produit
infini, la topologie boîte est strictement plus fine que la topologie produit.
Comme nous le verrons (section~\ref{sec:boite-problemes}), la topologie boîte
possède de mauvaises propriétés : c'est la topologie produit qui est la « bonne »
topologie sur un produit infini.
\end{remarque}

%-----------------------------------------------------------------------------
\section{Propriété universelle de la topologie produit}
\label{sec:propriete-universelle-produit}
\index{propriété universelle!topologie produit}

\begin{theoreme}[La topologie produit est la plus grossière rendant les projections continues]
\label{thm:produit-plus-grossiere}
La topologie produit sur $\prod_{i \in I} X_i$ est la topologie la plus grossière
sur $\prod_{i \in I} X_i$ rendant toutes les projections $\pi_i$ continues.
\end{theoreme}

\begin{proof}
Soit $\mathcal{T}$ une topologie sur $X = \prod_{i \in I} X_i$ rendant toutes les
$\pi_i$ continues. Alors pour tout $i \in I$ et tout $U_i \in \mathcal{T}_i$, on a
$\pi_i^{-1}(U_i) \in \mathcal{T}$. Comme la topologie produit $\mathcal{T}_{\mathrm{prod}}$
est engendrée par ces ensembles, on a $\mathcal{T}_{\mathrm{prod}} \subset \mathcal{T}$.
Réciproquement, les projections sont continues pour $\mathcal{T}_{\mathrm{prod}}$
par construction.
\end{proof}

\begin{theoreme}[Propriété universelle du produit]
\label{thm:universelle-produit}
\index{propriété universelle!produit}
Soit $(X_i)_{i \in I}$ une famille d'espaces topologiques et $X = \prod_{i \in I} X_i$
muni de la topologie produit. Pour tout espace topologique~$Y$ et toute famille
d'applications continues $(f_i \colon Y \to X_i)_{i \in I}$, il existe une unique
application continue $f \colon Y \to X$ telle que $\pi_i \circ f = f_i$ pour tout
$i \in I$.
\[
\begin{tikzcd}
  & Y \arrow[dl, "f_i"'] \arrow[d, dashed, "\exists!\, f"] \\
  X_i & X \arrow[l, "\pi_i"]
\end{tikzcd}
\]
\end{theoreme}

\begin{proof}
L'application $f$ est nécessairement définie par $f(y) = (f_i(y))_{i \in I}$.
Vérifions sa continuité. Il suffit de vérifier que $f^{-1}(\pi_i^{-1}(U_i))$
est ouvert dans~$Y$ pour tout $i \in I$ et tout $U_i$ ouvert dans~$X_i$. Or
$f^{-1}(\pi_i^{-1}(U_i)) = (\pi_i \circ f)^{-1}(U_i) = f_i^{-1}(U_i)$,
qui est ouvert car $f_i$ est continue.
\end{proof}

\begin{corollaire}
\label{cor:continuite-produit}
\index{continuité!dans un produit}
Une application $f \colon Y \to \prod_{i \in I} X_i$ est continue si et seulement
si chaque composante $\pi_i \circ f \colon Y \to X_i$ est continue.
\end{corollaire}

%-----------------------------------------------------------------------------
\section{Projections et lemme du tube}
\label{sec:projections-tube}

\begin{proposition}[Les projections sont ouvertes]
\label{prop:projection-ouverte}
\index{projection!ouverte}
Les projections $\pi_i \colon \prod_{j \in I} X_j \to X_i$ sont des applications
ouvertes.
\end{proposition}

\begin{proof}
Il suffit de vérifier sur les ouverts de base. Si $U = \prod_{j \in I} U_j$ est
un ouvert de base (avec $U_j = X_j$ sauf pour un nombre fini d'indices), alors
$\pi_i(U) = U_i$ qui est ouvert dans~$X_i$.
\end{proof}

\begin{remarque}
\label{rem:projection-non-fermee}
\index{projection!non fermée}
Les projections ne sont pas nécessairement fermées. Par exemple, la projection
$\pi_1 \colon \mathbb{R}^2 \to \mathbb{R}$ envoie l'hyperbole fermée
$\{(x,y) \in \mathbb{R}^2 : xy = 1\}$ sur $\mathbb{R} \setminus \{0\}$, qui
n'est pas fermé.
\end{remarque}

\begin{lemme}[Lemme du tube]
\label{lem:tube}
\index{lemme du tube}
\index{tube (lemme du)}
Soient $X$ et~$Y$ des espaces topologiques avec $Y$ compact. Soit $x_0 \in X$
et $N$ un ouvert de $X \times Y$ contenant la « tranche »
$\{x_0\} \times Y$. Alors il existe un ouvert $W$ de~$X$ contenant~$x_0$ tel que
$\{x_0\} \times Y \subset W \times Y \subset N$.
(L'ouvert $W \times Y$ est appelé un \emph{tube} autour de $\{x_0\} \times Y$.)
\end{lemme}

\begin{proof}
Pour chaque $y \in Y$, le point $(x_0, y) \in N$. Comme $N$ est ouvert, il existe
des ouverts $U_y \subset X$ et $V_y \subset Y$ tels que
$(x_0, y) \in U_y \times V_y \subset N$.

La famille $(V_y)_{y \in Y}$ est un recouvrement ouvert de~$Y$. Comme $Y$ est
compact, il existe $y_1, \ldots, y_m \in Y$ tels que
$Y = V_{y_1} \cup \cdots \cup V_{y_m}$.

Posons $W = U_{y_1} \cap \cdots \cap U_{y_m}$, qui est un ouvert de~$X$ contenant~$x_0$.
Montrons que $W \times Y \subset N$. Soit $(x, y) \in W \times Y$. Il existe
$k \in \{1, \ldots, m\}$ tel que $y \in V_{y_k}$. Comme $x \in W \subset U_{y_k}$,
on a $(x, y) \in U_{y_k} \times V_{y_k} \subset N$.
\end{proof}

%-----------------------------------------------------------------------------
\section{Topologie quotient}
\label{sec:topologie-quotient}
\index{topologie quotient}

\begin{definition}[Topologie quotient]
\label{def:topologie-quotient}
Soit $(X, \mathcal{T})$ un espace topologique, $\sim$ une relation d'équivalence
sur~$X$, et $p \colon X \to X/{\sim}$ la surjection canonique.
La \emph{topologie quotient} sur $X/{\sim}$ est
\[
\mathcal{T}_{X/\sim} = \bigl\{ U \subset X/{\sim} : p^{-1}(U) \in \mathcal{T} \bigr\}.
\]
L'application $p$ est alors continue et surjective ; elle est appelée
\emph{application quotient}\index{application quotient}.
\end{definition}

\begin{remarque}
\label{rem:quotient-plus-fine}
La topologie quotient est la topologie la plus fine sur $X/{\sim}$ rendant $p$ continue.
\end{remarque}

\begin{definition}[Application quotient]
\label{def:application-quotient}
\index{application quotient}
Plus généralement, une surjection continue $q \colon X \to Y$ est une
\emph{application quotient} si la topologie de~$Y$ est la topologie quotient
induite par~$q$, c'est-à-dire si
\[
V \text{ ouvert dans } Y \iff q^{-1}(V) \text{ ouvert dans } X.
\]
\end{definition}

\begin{definition}[Ensemble saturé]
\label{def:sature}
\index{saturé (ensemble)}
Soit $q \colon X \to Y$ une surjection. Une partie $A \subset X$ est dite
\emph{saturée} (par rapport à~$q$) si $A = q^{-1}(q(A))$, c'est-à-dire si
$A$ est une réunion de fibres de~$q$.
\end{definition}

\begin{proposition}[Conditions suffisantes pour être une application quotient]
\label{prop:conditions-quotient}
\index{application quotient!conditions}
Soit $q \colon X \to Y$ une surjection continue. Chacune des conditions suivantes
entraîne que $q$ est une application quotient :
\begin{enumerate}[label=(\roman*)]
  \item $q$ est ouverte ;
  \item $q$ est fermée ;
  \item $X$ est compact et $Y$ est séparé (Hausdorff).
\end{enumerate}
\end{proposition}

\begin{proof}
\noindent\textbf{(i).} Soit $V \subset Y$ tel que $q^{-1}(V)$ est ouvert. Comme $q$ est
ouverte et surjective, $V = q(q^{-1}(V))$ est ouvert.

\noindent\textbf{(ii).} Soit $V \subset Y$ tel que $q^{-1}(V)$ est ouvert. Alors
$q^{-1}(Y \setminus V) = X \setminus q^{-1}(V)$ est fermé. Comme $q$ est fermée,
$q(q^{-1}(Y \setminus V)) = Y \setminus V$ est fermé, donc $V$ est ouvert.

\noindent\textbf{(iii).} Découle de~(ii), car une application continue d'un
compact dans un séparé est fermée.
\end{proof}

%-----------------------------------------------------------------------------
\section{Propriété universelle de la topologie quotient}
\label{sec:propriete-universelle-quotient}
\index{propriété universelle!topologie quotient}

\begin{theoreme}[Propriété universelle du quotient]
\label{thm:universelle-quotient}
Soit $q \colon X \to Y$ une application quotient. Soit $Z$ un espace topologique
et $f \colon X \to Z$ une application continue constante sur les fibres de~$q$
(c'est-à-dire : $q(x) = q(x') \Rightarrow f(x) = f(x')$). Alors il existe une unique
application continue $\bar{f} \colon Y \to Z$ telle que $f = \bar{f} \circ q$.
\[
\begin{tikzcd}
  X \arrow[r, "q"] \arrow[dr, "f"'] & Y \arrow[d, dashed, "\exists!\, \bar{f}"] \\
  & Z
\end{tikzcd}
\]
\end{theoreme}

\begin{proof}
L'application $\bar{f}$ est définie par $\bar{f}(q(x)) = f(x)$, ce qui est licite
car $f$ est constante sur les fibres. L'unicité est immédiate par surjectivité de~$q$.
Pour la continuité : si $W$ est ouvert dans~$Z$, alors
$q^{-1}(\bar{f}^{-1}(W)) = f^{-1}(W)$ est ouvert dans~$X$ (par continuité de~$f$).
Par définition de la topologie quotient, $\bar{f}^{-1}(W)$ est ouvert dans~$Y$.
\end{proof}

%-----------------------------------------------------------------------------
\section{Composition d'applications quotient}
\label{sec:composition-quotient}

\begin{theoreme}[Composition d'applications quotient]
\label{thm:composition-quotient}
\index{application quotient!composition}
Soient $q_1 \colon X \to Y$ et $q_2 \colon Y \to Z$ des applications quotient.
Alors $q_2 \circ q_1 \colon X \to Z$ est une application quotient.
\end{theoreme}

\begin{proof}
L'application $q_2 \circ q_1$ est continue et surjective. Soit $W \subset Z$ tel
que $(q_2 \circ q_1)^{-1}(W) = q_1^{-1}(q_2^{-1}(W))$ est ouvert dans~$X$.
Comme $q_1$ est quotient, $q_2^{-1}(W)$ est ouvert dans~$Y$.
Comme $q_2$ est quotient, $W$ est ouvert dans~$Z$.
\end{proof}

%-----------------------------------------------------------------------------
\section{Exemples classiques de quotients}
\label{sec:exemples-quotients}

\begin{exemple}[Le cercle $S^1$]
\label{ex:cercle-quotient}
\index{cercle $S^1$!comme quotient}
Définissons sur $[0,1]$ la relation $\sim$ par : $x \sim y$ si et seulement si
$x = y$ ou $\{x,y\} = \{0,1\}$. L'espace quotient $[0,1]/{\sim}$ est
homéomorphe au cercle $S^1 = \{z \in \mathbb{C} : |z| = 1\}$.

En effet, l'application $f \colon [0,1] \to S^1$ définie par
$f(t) = e^{2\pi i t}$ est continue, surjective, constante sur les classes de~$\sim$
($f(0) = f(1) = 1$), et induit un homéomorphisme $\bar{f} \colon [0,1]/{\sim} \to S^1$
(bijectif continu d'un compact vers un séparé, donc homéomorphisme).
\end{exemple}

\begin{exemple}[Le tore $\mathbb{T}^2$]
\label{ex:tore-quotient}
\index{tore!comme quotient}
Sur le carré $[0,1]^2$, on définit la relation d'équivalence engendrée par
$(x,0) \sim (x,1)$ pour tout $x \in [0,1]$ et $(0,y) \sim (1,y)$ pour tout
$y \in [0,1]$. L'espace quotient est le \emph{tore} $\mathbb{T}^2$, homéomorphe
à $S^1 \times S^1$.

\begin{center}
\begin{tikzpicture}[scale=2.5]
  % Carré avec identifications
  \fill[blue!5] (0,0) rectangle (1,1);
  \draw[red, very thick, ->] (0,0) -- (0.55,0);
  \draw[red, very thick] (0.55,0) -- (1,0);
  \draw[red, very thick, ->] (0,1) -- (0.55,1);
  \draw[red, very thick] (0.55,1) -- (1,1);
  \draw[blue, very thick, ->] (0,0) -- (0,0.55);
  \draw[blue, very thick] (0,0.55) -- (0,1);
  \draw[blue, very thick, ->] (1,0) -- (1,0.55);
  \draw[blue, very thick] (1,0.55) -- (1,1);
  % Sommets
  \fill (0,0) circle (1pt); \fill (1,0) circle (1pt);
  \fill (0,1) circle (1pt); \fill (1,1) circle (1pt);
  % Labels
  \node[below left] at (0,0) {$A$};
  \node[below right] at (1,0) {$A$};
  \node[above left] at (0,1) {$A$};
  \node[above right] at (1,1) {$A$};
  \node[below, red] at (0.5,0) {$a$};
  \node[above, red] at (0.5,1) {$a$};
  \node[left, blue] at (0,0.5) {$b$};
  \node[right, blue] at (1,0.5) {$b$};
  % Flèche et résultat
  \draw[->, thick] (1.3,0.5) -- (1.7,0.5);
  \node at (1.5,0.7) {\small quotient};
  % Tore schématique
  \begin{scope}[shift={(2.7,0.5)}]
    \draw[thick] (0,0) ellipse (0.7 and 0.4);
    \draw[thick] (-0.35,0) arc (180:360:0.35 and 0.15);
    \draw[thick, dashed] (-0.35,0) arc (180:0:0.35 and 0.15);
    \node[below] at (0,-0.55) {$\mathbb{T}^2$};
  \end{scope}
\end{tikzpicture}
\end{center}
\end{exemple}

\begin{exemple}[Le ruban de Möbius]
\label{ex:mobius-quotient}
\index{Möbius!ruban de}
Sur le rectangle $[0,1] \times [-1,1]$, on identifie $(0,y) \sim (1,-y)$ pour
tout $y \in [-1,1]$. Le quotient est le \emph{ruban de Möbius}.

\begin{center}
\begin{tikzpicture}[scale=2]
  \fill[green!5] (0,-0.7) rectangle (1,0.7);
  \draw[red, very thick, ->] (0,-0.7) -- (0,0.05);
  \draw[red, very thick] (0,0.05) -- (0,0.7);
  \draw[blue, very thick, ->] (1,0.7) -- (1,-0.05);
  \draw[blue, very thick] (1,-0.05) -- (1,-0.7);
  \draw[thick] (0,0.7) -- (1,0.7);
  \draw[thick] (0,-0.7) -- (1,-0.7);
  % Les flèches opposées indiquent l'inversion
  \node[left, red] at (0,0) {$a$};
  \node[right, blue] at (1,0) {$a$};
  % Note
  \node[below] at (0.5,-0.9) {Ruban de Möbius};
\end{tikzpicture}
\end{center}
\end{exemple}

\begin{exemple}[La bouteille de Klein]
\label{ex:klein-quotient}
\index{Klein!bouteille de}
Sur $[0,1]^2$, on identifie $(x,0) \sim (x,1)$ et $(0,y) \sim (1,1-y)$.
Le quotient est la \emph{bouteille de Klein}, surface non orientable qui ne peut
se plonger dans~$\mathbb{R}^3$.

\begin{center}
\begin{tikzpicture}[scale=2.5]
  \fill[yellow!10] (0,0) rectangle (1,1);
  \draw[red, very thick, ->] (0,0) -- (0.55,0);
  \draw[red, very thick] (0.55,0) -- (1,0);
  \draw[red, very thick, ->] (0,1) -- (0.55,1);
  \draw[red, very thick] (0.55,1) -- (1,1);
  \draw[blue, very thick, ->] (0,0) -- (0,0.55);
  \draw[blue, very thick] (0,0.55) -- (0,1);
  \draw[blue, very thick, <-] (1,0.45) -- (1,1);
  \draw[blue, very thick] (1,0) -- (1,0.45);
  \node[below, red] at (0.5,0) {$a$};
  \node[above, red] at (0.5,1) {$a$};
  \node[left, blue] at (0,0.5) {$b$};
  \node[right, blue] at (1,0.5) {$b$};
  \node[below] at (0.5,-0.2) {Bouteille de Klein};
\end{tikzpicture}
\end{center}
\end{exemple}

\begin{exemple}[L'espace projectif réel $\mathbb{R}P^2$]
\label{ex:RP2-quotient}
\index{espace projectif réel $\mathbb{R}P^2$}
L'espace projectif réel $\mathbb{R}P^2$ peut être défini de plusieurs manières
équivalentes :
\begin{enumerate}[label=(\alph*)]
  \item Comme quotient de $\mathbb{R}^3 \setminus \{0\}$ par la relation
        $x \sim \lambda x$ pour tout $\lambda \in \mathbb{R}^*$ : les points
        de~$\mathbb{R}P^2$ sont les droites vectorielles de~$\mathbb{R}^3$.
  \item Comme quotient de~$S^2$ par la relation antipodale $x \sim -x$.
  \item Comme quotient du disque $D^2$ en identifiant les points antipodaux
        du bord : $(x,y) \sim (-x,-y)$ pour $(x,y) \in S^1 = \partial D^2$.
\end{enumerate}

\begin{center}
\begin{tikzpicture}[scale=2]
  % Disque
  \fill[orange!10] (0,0) circle (1);
  \draw[thick] (0,0) circle (1);
  % Identifications antipodales sur le bord
  \foreach \a in {0,30,60,90,120,150} {
    \fill[red] ({\a}:1) circle (1.5pt);
    \fill[red] ({\a+180}:1) circle (1.5pt);
    \draw[red, dashed, thin] ({\a}:1) -- ({\a+180}:1);
  }
  % Flèches d'identification
  \draw[blue, very thick, ->] (90:1) arc (90:0:1);
  \draw[blue, very thick, <-] (-90:1) arc (-90:0:1);
  \draw[blue, very thick, ->] (90:1) arc (90:180:1);
  \draw[blue, very thick, <-] (270:1) arc (270:180:1);
  \node[below] at (0,-1.3) {$\mathbb{R}P^2$ : points antipodaux identifiés};
\end{tikzpicture}
\end{center}
\end{exemple}

%-----------------------------------------------------------------------------
\section{La topologie boîte et ses défauts}
\label{sec:boite-problemes}
\index{topologie boîte!défauts}

Nous justifions ici pourquoi la topologie produit, et non la topologie boîte, est
la construction appropriée pour les produits infinis.

\begin{exemple}[La topologie boîte n'est pas compacte]
\label{ex:boite-non-compacte}
\index{topologie boîte!non compacte}
Considérons $X = \prod_{n=1}^{\infty} \{0,1\}$ où chaque facteur $\{0,1\}$ porte
la topologie discrète. Pour la topologie produit, $X$ est compact par le
théorème de Tychonoff\index{Tychonoff!théorème de} (produit de compacts). Pour la
topologie boîte, $X$ est discret (tout singleton $\{x\}$ est ouvert comme
$\prod_{n} \{x_n\}$), donc non compact (puisque $X$ est infini non dénombrable).
\end{exemple}

\begin{exemple}[La topologie boîte n'est pas connexe]
\label{ex:boite-non-connexe}
\index{topologie boîte!non connexe}
Posons $X = \prod_{n=1}^{\infty} \mathbb{R}$ avec la topologie boîte. L'ensemble
\[
U = \prod_{n=1}^{\infty} {]{-1/n}, 1/n[}
\]
est ouvert dans la topologie boîte mais n'est pas ouvert dans la topologie produit.
On peut montrer que $\mathbb{R}^{\mathbb{N}}$ muni de la topologie boîte n'est pas
connexe, alors qu'il l'est pour la topologie produit
(théorème~\ref{thm:produit-arbitraire-connexes}).

Plus précisément, soit $B$ l'ensemble des suites bornées et $U$ l'ensemble des
suites non bornées dans $\mathbb{R}^{\mathbb{N}}$. On peut vérifier que $B$ et $U$
sont tous deux ouverts pour la topologie boîte, donnant une déconnexion de
$\mathbb{R}^{\mathbb{N}}$.
\end{exemple}

\begin{remarque}
\label{rem:boite-continuite}
La propriété universelle (théorème~\ref{thm:universelle-produit}) échoue pour la
topologie boîte : une application $f \colon Y \to \prod_{i \in I} X_i$ peut avoir
toutes ses composantes $\pi_i \circ f$ continues sans être elle-même continue pour
la topologie boîte. Par exemple, $f \colon \mathbb{R} \to \mathbb{R}^{\mathbb{N}}$
définie par $f(t) = (t, t, t, \ldots)$ a des composantes continues (l'identité),
mais $f$ n'est pas continue pour la topologie boîte : l'image réciproque de
$\prod_{n} {]{-1/n}, 1/n[}$ est $\{0\}$, qui n'est pas ouvert dans~$\mathbb{R}$.
\end{remarque}

%-----------------------------------------------------------------------------
\section{Topologie somme (coproduit)}
\label{sec:topologie-somme}
\index{topologie somme}
\index{coproduit topologique}
\index{réunion disjointe}

\begin{definition}[Somme topologique]
\label{def:topologie-somme}
Soit $(X_i, \mathcal{T}_i)_{i \in I}$ une famille d'espaces topologiques.
La \emph{somme topologique} (ou \emph{réunion disjointe topologique}) est
l'ensemble $\bigsqcup_{i \in I} X_i = \bigcup_{i \in I} X_i \times \{i\}$ muni
de la topologie
\[
\mathcal{T} = \Bigl\{ U \subset \bigsqcup_{i \in I} X_i :
U \cap (X_i \times \{i\}) \text{ est ouvert dans } X_i \times \{i\}
\text{ pour tout } i \in I \Bigr\}.
\]
Les injections canoniques $\iota_i \colon X_i \to \bigsqcup X_i$, $x \mapsto (x,i)$,
sont des homéomorphismes sur leur image, et chaque $\iota_i(X_i)$ est ouvert-fermé
dans $\bigsqcup X_i$.
\end{definition}

\begin{theoreme}[Propriété universelle du coproduit]
\label{thm:universelle-coproduit}
\index{propriété universelle!coproduit}
Soit $(X_i)_{i \in I}$ une famille d'espaces topologiques et
$\bigsqcup_{i \in I} X_i$ leur somme topologique. Pour tout espace~$Y$ et toute
famille d'applications continues $(f_i \colon X_i \to Y)_{i \in I}$, il existe une
unique application continue $f \colon \bigsqcup X_i \to Y$ telle que
$f \circ \iota_i = f_i$ pour tout~$i$.
\[
\begin{tikzcd}
  X_i \arrow[r, "\iota_i"] \arrow[dr, "f_i"'] & \bigsqcup X_i \arrow[d, dashed, "\exists!\, f"] \\
  & Y
\end{tikzcd}
\]
\end{theoreme}

\begin{proof}
On définit $f(x, i) = f_i(x)$. Pour tout ouvert $V \subset Y$,
\[
f^{-1}(V) \cap (X_i \times \{i\}) = \iota_i(f_i^{-1}(V))
\]
est ouvert dans $X_i \times \{i\}$ car $f_i$ est continue et $\iota_i$ est un
homéomorphisme. Par définition de la topologie somme, $f^{-1}(V)$ est ouvert.
\end{proof}

%-----------------------------------------------------------------------------
\section{Exemples résolus}
\label{sec:exemples-produit-quotient}

\begin{exemple}[Convergence dans un produit]
\label{ex:convergence-produit}
\index{convergence!dans un produit}
Soit $(x^{(n)})_{n \geqslant 1}$ une suite dans $\prod_{i \in I} X_i$ muni de la
topologie produit, et $x \in \prod_{i \in I} X_i$. Montrons que
$x^{(n)} \to x$ si et seulement si $\pi_i(x^{(n)}) \to \pi_i(x)$ pour tout $i \in I$.

\emph{Condition nécessaire.} Chaque $\pi_i$ est continue, donc $\pi_i(x^{(n)}) \to \pi_i(x)$.

\emph{Condition suffisante.} Soit $U$ un ouvert de base du produit contenant~$x$.
Alors $U = \bigcap_{k=1}^{m} \pi_{i_k}^{-1}(U_{i_k})$ pour des ouverts
$U_{i_k}$ avec $\pi_{i_k}(x) \in U_{i_k}$. Pour chaque~$k$, il existe $N_k$ tel
que $\pi_{i_k}(x^{(n)}) \in U_{i_k}$ pour $n \geqslant N_k$. Pour
$n \geqslant \max(N_1, \ldots, N_m)$, on a $x^{(n)} \in U$.
\end{exemple}

\begin{exemple}[Le cylindre comme quotient]
\label{ex:cylindre-quotient}
\index{cylindre!comme quotient}
Le cylindre $S^1 \times [0,1]$ peut être obtenu comme quotient de
$[0,1]^2$ en identifiant seulement $(0,y) \sim (1,y)$ pour tout $y \in [0,1]$
(sans identifier les bords horizontaux). C'est un cas intermédiaire entre le
rectangle (pas d'identification) et le tore (identification des deux paires de
côtés opposés).
\end{exemple}

\begin{exemple}[Quotient par un sous-espace]
\label{ex:quotient-sous-espace}
\index{quotient!par un sous-espace}
Soit $X$ un espace topologique et $A \subset X$ un sous-espace non vide. L'espace
quotient $X/A$ est obtenu en identifiant tous les points de~$A$ à un seul point,
noté~$[A]$. La relation est : $x \sim y$ si et seulement si $x = y$ ou
$\{x, y\} \subset A$.

Par exemple :
\begin{enumerate}[label=(\alph*)]
  \item $[0,1] / \{0,1\} \cong S^1$ (on identifie les extrémités).
  \item $D^2 / S^1 \cong S^2$ (on contracte le bord du disque en un point).
  \item $S^n / S^{n-1} \cong S^n$ (en identifiant l'équateur,
        on obtient un bouquet de deux sphères, non $S^n$; attention à la
        formulation exacte de l'identification).
\end{enumerate}
\end{exemple}

%-----------------------------------------------------------------------------
\section{Exercices}
\label{sec:exo-produit-quotient}

\begin{exercice}[$\star$]
\label{exo:produit-hausdorff}
Montrer que $\prod_{i \in I} X_i$ est séparé (Hausdorff) si et seulement si
chaque~$X_i$ est séparé.
\end{exercice}

\begin{exercice}[$\star$]
\label{exo:produit-denombrable-metrisable}
Montrer que si $(X_n, d_n)_{n \geqslant 1}$ est une famille dénombrable d'espaces
métriques, alors $\prod_{n=1}^{\infty} X_n$ muni de la topologie produit est
métrisable. \emph{Indication :} considérer la distance
$d(x,y) = \sum_{n=1}^{\infty} \frac{1}{2^n} \min(d_n(x_n, y_n), 1)$.
\end{exercice}

\begin{exercice}[$\star$]
\label{exo:projection-non-fermee}
Trouver un fermé $F \subset \mathbb{R}^2$ dont la projection sur le premier axe
n'est pas fermée.
\end{exercice}

\begin{exercice}[$\star\star$]
\label{exo:tube-lemme-application}
Soient $X$ un espace topologique et $Y$ un espace compact. Montrer que la
projection $\pi_1 \colon X \times Y \to X$ est une application fermée.
\emph{Indication :} utiliser le lemme du tube~\ref{lem:tube}.
\end{exercice}

\begin{exercice}[$\star\star$]
\label{exo:quotient-non-hausdorff}
Soit $X = \mathbb{R}$ et la relation d'équivalence $x \sim y$ si et seulement si
$x - y \in \mathbb{Q}$. Montrer que $\mathbb{R}/\mathbb{Q}$ n'est pas séparé.
En fait, montrer que la seule topologie sur cet espace est la topologie grossière.
\end{exercice}

\begin{exercice}[$\star\star$]
\label{exo:tore-produit}
Montrer en détail que le tore $\mathbb{T}^2$ (exemple~\ref{ex:tore-quotient})
est homéomorphe à $S^1 \times S^1$.
\end{exercice}

\begin{exercice}[$\star\star$]
\label{exo:quotient-ouvert}
Soit $G$ un groupe topologique et $H$ un sous-groupe de~$G$. On munit $G/H$
(ensemble des classes à gauche $gH$) de la topologie quotient.
\begin{enumerate}[label=(\alph*)]
  \item Montrer que la projection $\pi \colon G \to G/H$ est une application ouverte.
  \item En déduire que $G/H$ est séparé si et seulement si $H$ est fermé dans~$G$.
\end{enumerate}
\end{exercice}

\begin{exercice}[$\star\star\star$]
\label{exo:boite-pas-premier-denombrable}
Montrer que $\mathbb{R}^{\mathbb{N}}$ muni de la topologie boîte ne satisfait pas
le premier axiome de dénombrabilité.
\emph{Indication :} montrer que le point $0 = (0, 0, \ldots)$ n'admet pas de base
dénombrable de voisinages.
\end{exercice}

\begin{exercice}[$\star\star\star$]
\label{exo:cone}
\index{cône topologique}
Le \emph{cône} sur un espace topologique~$X$ est $CX = (X \times [0,1]) / (X \times \{1\})$.
\begin{enumerate}[label=(\alph*)]
  \item Montrer que $CX$ est contractile (c'est-à-dire a le type d'homotopie d'un
        point).
  \item Montrer que $CS^{n-1} \cong D^n$ (le cône sur la sphère est homéomorphe
        au disque).
\end{enumerate}
\end{exercice}

\begin{exercice}[$\star\star\star$]
\label{exo:suspension}
\index{suspension topologique}
La \emph{suspension} d'un espace~$X$ est $SX = (X \times [-1,1]) / {\sim}$, où
$(x,-1) \sim (x',-1)$ et $(x,1) \sim (x',1)$ pour tous $x, x' \in X$.
\begin{enumerate}[label=(\alph*)]
  \item Montrer que $SS^n \cong S^{n+1}$.
  \item Montrer que si $X$ est compact et connexe par arcs, alors $SX$ est
        connexe par arcs.
\end{enumerate}
\end{exercice}

%-----------------------------------------------------------------------------
\section*{Résumé du chapitre~\thechapter}
\addcontentsline{toc}{section}{Résumé du chapitre~\thechapter}

\begin{itemize}
  \item La \textbf{topologie produit} est la topologie la plus grossière rendant les
        projections continues. Les ouverts de base sont les produits d'ouverts, avec
        la contrainte que seul un nombre fini de facteurs diffèrent de l'espace total.
  \item La \textbf{propriété universelle} du produit : $f \colon Y \to \prod X_i$ est
        continue si et seulement si chaque $\pi_i \circ f$ est continue.
  \item La \textbf{topologie boîte} coïncide avec la topologie produit pour les
        produits finis, mais donne un espace aux propriétés indésirables pour les
        produits infinis (non compact, non connexe, perte de la propriété universelle).
  \item Le \textbf{lemme du tube} : si $Y$ est compact et $N$ est un ouvert contenant
        $\{x_0\} \times Y$, alors $N$ contient un « tube » $W \times Y$.
  \item Les \textbf{projections} sont ouvertes mais pas nécessairement fermées ;
        elles deviennent fermées lorsque le facteur est compact.
  \item La \textbf{topologie quotient} sur $X/{\sim}$ est la topologie la plus fine
        rendant la surjection canonique continue. Sa propriété universelle caractérise
        les applications continues depuis le quotient.
  \item Les \textbf{espaces quotients classiques} : $S^1 = [0,1]/{\sim}$,
        $\mathbb{T}^2 = [0,1]^2/{\sim}$, ruban de Möbius, bouteille de Klein,
        $\mathbb{R}P^2$.
  \item La \textbf{somme topologique} (coproduit) est la réunion disjointe munie de
        la topologie la plus fine rendant les injections continues.
\end{itemize}
%%%%%%%%%%%%%%%%%%%%%%%%%%%%%%%%%%%%%%%%%%%%%%%%%%%%%%%%%%%%
%% PARTIE 5 — Chapitres 9–10, Annexe, Bibliographie, Index
%%%%%%%%%%%%%%%%%%%%%%%%%%%%%%%%%%%%%%%%%%%%%%%%%%%%%%%%%%%%

%% ============================================================
%%  CHAPITRE 9 — Lemme d'Urysohn et Théorème d'Extension de Tietze
%% ============================================================
\chapter{Lemme d'Urysohn et théorème d'extension de Tietze}
\label{chap:urysohn-tietze}
\index{lemme d'Urysohn}\index{théorème de Tietze}

\medskip
\noindent
\textit{Comment séparer deux fermés disjoints non seulement par des ouverts,
mais par une fonction continue à valeurs réelles~? Le lemme d'Urysohn
répond à cette question fondamentale pour les espaces normaux et ouvre
la voie au théorème d'extension de Tietze, pierre angulaire de l'analyse
sur les espaces topologiques. Ces résultats fournissent en outre le
théorème de métrisation d'Urysohn, qui caractérise une large classe
d'espaces métrisables.}

%% -------------------------------------------------------
%%  9.1  Espaces complètement réguliers
%% -------------------------------------------------------
\section{Espaces complètement réguliers}
\label{sec:completement-regulier}
\index{espace complètement régulier}\index{espace de Tychonoff}\index{T@$T_{3\frac{1}{2}}$}

\begin{definition}[Espace complètement régulier]
\label{def:completement-regulier}
Un espace topologique $X$ est dit \emph{complètement régulier}
(ou \emph{$T_{3\frac{1}{2}}$}, ou \emph{espace de Tychonoff}\index{Tychonoff!espace de})
s'il est $T_1$ et si, pour tout fermé $F\subset X$ et tout point
$x\in X\setminus F$, il existe une application continue
$f\colon X\to [0,1]$ telle que $f(x)=0$ et $f|_F\equiv 1$.
\end{definition}

\begin{remarque}
\label{rem:hierarchie-separation-T3half}
On a les implications strictes :
\[
T_4 \;\Longrightarrow\; T_{3\tfrac{1}{2}} \;\Longrightarrow\; T_3
  \;\Longrightarrow\; T_2 \;\Longrightarrow\; T_1 \;\Longrightarrow\; T_0.
\]
La complète régularité est exactement la condition nécessaire et suffisante
pour qu'un espace $T_1$ admette un plongement dans un cube
$[0,1]^I$\index{cube de Tychonoff} (théorème de plongement de Tychonoff).
\end{remarque}

\begin{proposition}
\label{prop:sous-espace-completement-regulier}
Tout sous-espace d'un espace complètement régulier est complètement
régulier. Tout produit d'espaces complètement réguliers est complètement
régulier.\index{produit!complète régularité}
\end{proposition}

\begin{proof}
Soit $A\subset X$ un sous-espace d'un espace complètement régulier $X$.
Soient $F$ un fermé de $A$ et $a\in A\setminus F$. Il existe un fermé
$G$ de $X$ tel que $F=G\cap A$. Comme $a\notin G$, il existe
$f\colon X\to[0,1]$ continue avec $f(a)=0$ et $f|_G\equiv 1$.
La restriction $f|_A$ convient.

Pour le produit $X=\prod_{i\in I}X_i$, soient $F$ fermé de $X$ et
$x\notin F$. Il existe un ouvert élémentaire
$U=\bigcap_{k=1}^n \pi_{i_k}^{-1}(U_{i_k})$ contenant $x$ et disjoint
de $F$. Pour chaque $k$, le fermé $X_{i_k}\setminus U_{i_k}$ est séparé
de $\pi_{i_k}(x)$ par une fonction continue $f_k\colon X_{i_k}\to[0,1]$.
La fonction $f=\max_k(f_k\circ\pi_{i_k})$ sépare $x$ de $F$.
\end{proof}

\begin{exemple}
\label{ex:metrique-completement-regulier}
Tout espace métrique $(X,d)$ est complètement régulier. En effet,
si $F$ est fermé et $x\notin F$, la fonction
$f(y)=\min\!\bigl(1,\,d(y,F)^{-1}\,d(y,x)\bigr)$... plus simplement,
$f(y)=\frac{d(y,F)}{d(y,F)+d(y,x)}$ est continue, vaut $0$ sur $F$
et $1$ en $x$. En inversant les rôles :
$g(y)=\min\!\bigl(1,\,d(y,F)\big/\delta\bigr)$ où
$\delta=d(x,F)>0$ donne $g(x)\ge 1$ et $g|_F\equiv 0$, que l'on
renverse en $1-g$ pour satisfaire la convention.
\end{exemple}

%% -------------------------------------------------------
%%  9.2  Lemme d'Urysohn
%% -------------------------------------------------------
\section{Lemme d'Urysohn}
\label{sec:lemme-urysohn}
\index{lemme d'Urysohn}

\begin{theoreme}[Lemme d'Urysohn, 1925]
\label{thm:urysohn}
Soit $X$ un espace topologique normal\index{espace normal} ($T_4$).
Pour tous fermés disjoints $A,B\subset X$, il existe une application
continue $f\colon X\to [0,1]$ telle que $f|_A\equiv 0$ et $f|_B\equiv 1$.
\end{theoreme}

\begin{proof}
La démonstration procède en trois étapes.

\medskip\noindent
\textbf{Étape 1 : construction d'une famille d'ouverts indexée par les
rationnels dyadiques.}\index{rationnels dyadiques}

Notons $D=\bigl\{\,k/2^n : k,n\in\mathbf{N},\; 0\le k\le 2^n\bigr\}
\cap [0,1]$ l'ensemble des rationnels dyadiques de $[0,1]$.
Nous allons construire, pour chaque $r\in D$, un ouvert $U_r$ de $X$
vérifiant :
\begin{enumerate}[label=(\roman*)]
  \item $A\subset U_r$ pour tout $r\in D$ ;
  \item $U_0 \supset A$ et $X\setminus B \supset U_1$ ;
  \item si $r<s$ dans $D$, alors $\overline{U_r}\subset U_s$.
\end{enumerate}

\emph{Initialisation.} Puisque $X$ est normal et $A\cap B=\varnothing$,
$A$ et $B$ sont fermés, l'ouvert $V=X\setminus B$ contient $A$.
Par normalité, il existe un ouvert $U_1$ tel que
$A\subset U_1\subset\overline{U_1}\subset V=X\setminus B$.
Appliquons de nouveau la normalité au couple $(A,\,X\setminus U_1)$ :
il existe un ouvert $U_0$ tel que
$A\subset U_0\subset\overline{U_0}\subset U_1$.

\emph{Récurrence.} Supposons les ouverts $U_r$ construits pour tous les
rationnels dyadiques de dénominateur $\le 2^n$ et vérifiant la condition
d'emboîtement. Pour chaque nouveau rationnel de la forme
$r=\frac{2m+1}{2^{n+1}}$ avec $0\le m < 2^n$, posons
$r^{-}=\frac{m}{2^n}$ et $r^{+}=\frac{m+1}{2^n}$; on a $r^{-}<r<r^{+}$
et, par hypothèse de récurrence,
$\overline{U_{r^{-}}}\subset U_{r^{+}}$.
Comme $\overline{U_{r^{-}}}$ est fermé et $U_{r^{+}}$ ouvert,
la normalité fournit un ouvert $U_r$ tel que
\[
\overline{U_{r^{-}}} \subset U_r \subset \overline{U_r} \subset U_{r^{+}}.
\]
Par récurrence, la famille $(U_r)_{r\in D}$ est construite et satisfait
$r<s \Rightarrow \overline{U_r}\subset U_s$.

\medskip\noindent
\textbf{Étape 2 : définition de la fonction $f$.}

Pour $x\in X$, posons
\[
f(x) \;=\;
\begin{cases}
\inf\,\{\,r\in D : x\in U_r\} & \text{si } x\in U_1,\\
1 & \text{si } x\notin U_1.
\end{cases}
\]
On a immédiatement $0\le f(x)\le 1$ pour tout $x$.
\begin{itemize}
  \item Si $x\in A$, alors $x\in U_r$ pour tout $r\in D$ (par (i)),
    donc $f(x)=\inf D\cap[0,1]=0$.
  \item Si $x\in B$, alors $x\notin U_1$ (car $\overline{U_1}\cap B
    =\varnothing$), donc $f(x)=1$.
\end{itemize}
Ainsi $f|_A\equiv 0$ et $f|_B\equiv 1$.

\medskip\noindent
\textbf{Étape 3 : continuité de $f$.}

Il suffit de montrer que $f^{-1}\bigl((-\infty,a)\bigr)$ et
$f^{-1}\bigl((a,+\infty)\bigr)$ sont ouverts pour tout $a\in\mathbf{R}$,
car les intervalles de cette forme engendrent la topologie de
$\mathbf{R}$.

\emph{Cas 1 :} $f^{-1}\bigl((-\infty,a)\bigr)=\{x\in X: f(x)<a\}$.
Si $f(x)<a$, il existe $r\in D$ tel que $f(x)\le r<a$ et $x\in U_r$.
Donc
\[
f^{-1}\bigl((-\infty,a)\bigr) = \bigcup_{\substack{r\in D\\ r<a}} U_r,
\]
qui est ouvert comme réunion d'ouverts.

\emph{Cas 2 :} $f^{-1}\bigl((a,+\infty)\bigr)=\{x\in X: f(x)>a\}$.
Si $f(x)>a$, il existe $s\in D$ tel que $a<s\le f(x)$ et
$x\notin \overline{U_s}$ (car si $x\in\overline{U_s}$ pour tout
$s>a$ dans $D$, alors $x\in U_t$ pour tout $t>s$ par l'emboîtement,
ce qui donnerait $f(x)\le s$). Ainsi $x\in X\setminus\overline{U_s}$
qui est ouvert. Plus précisément :
\[
f^{-1}\bigl((a,+\infty)\bigr) = \bigcup_{\substack{s\in D\\ s>a}}
  \bigl(X\setminus \overline{U_s}\bigr),
\]
qui est ouvert. En effet, si $f(x)>a$, choisissons $r,s\in D$ avec
$a<r<s<f(x)$. Alors $x\notin U_r$ (sinon $f(x)\le r$), donc
$x\notin\overline{U_r}$ est impossible directement ; procédons autrement.
Si $x\notin U_s$, alors en particulier $x\notin\overline{U_r}$
(car $\overline{U_r}\subset U_s$), d'où $x\in X\setminus\overline{U_r}$,
ouvert contenu dans $\{f>a\}$ puisque pour tout $y\notin\overline{U_r}$,
$y\notin U_r$, donc $f(y)\ge r>a$.

Ainsi $f$ est continue, ce qui achève la démonstration.
\end{proof}

% --- TikZ : Visualisation de la construction d'Urysohn ---
\begin{figure}[ht]
\centering
\begin{tikzpicture}[scale=1.1, >=stealth]
  % Fermés A et B
  \fill[blue!15] (-4.5,-1.2) rectangle (-3,-1.2+2.4);
  \fill[red!15] (3,-1.2) rectangle (4.5,-1.2+2.4);
  \node at (-3.75,0) {$A$};
  \node at (3.75,0) {$B$};
  % Ouverts emboîtés
  \draw[blue!70, thick, rounded corners=6pt]
    (-4.8,-1.5) rectangle (-2.2,1.5);
  \node[blue!70, above] at (-3.5,1.5) {$U_0$};
  \draw[blue!50, thick, rounded corners=6pt]
    (-5.0,-1.7) rectangle (-0.8,1.7);
  \node[blue!50, above] at (-2.9,1.7) {$U_{1/4}$};
  \draw[green!60!black, thick, rounded corners=6pt]
    (-5.2,-1.9) rectangle (0.5,1.9);
  \node[green!60!black, above] at (-2.3,1.9) {$U_{1/2}$};
  \draw[orange!70, thick, rounded corners=6pt]
    (-5.4,-2.1) rectangle (1.8,2.1);
  \node[orange!70, above] at (-1.8,2.1) {$U_{3/4}$};
  \draw[red!50, thick, rounded corners=6pt]
    (-5.6,-2.3) rectangle (2.7,2.3);
  \node[red!50, above] at (-1.4,2.3) {$U_1$};
  % Axe des valeurs
  \draw[->] (-5,-3) -- (5,-3) node[right] {$f$};
  \foreach \x/\lab in {-4/0, -1.5/1/4, 0.5/1/2, 2/3/4, 3.5/1}{
    \draw (\x,-3.1)--(\x,-2.9);
  }
  \node[below] at (-4,-3.1) {$0$};
  \node[below] at (-1.5,-3.1) {$\tfrac{1}{4}$};
  \node[below] at (0.5,-3.1) {$\tfrac{1}{2}$};
  \node[below] at (2,-3.1) {$\tfrac{3}{4}$};
  \node[below] at (3.5,-3.1) {$1$};
\end{tikzpicture}
\caption{Construction du lemme d'Urysohn : ouverts emboîtés
$(U_r)_{r\in D}$ entre $A$ (valeur $0$) et $B$ (valeur $1$).}
\label{fig:urysohn-construction}
\end{figure}

\begin{corollaire}
\label{cor:normal-implique-completement-regulier}
Tout espace normal\index{espace normal!complètement régulier} est
complètement régulier.
\end{corollaire}

\begin{proof}
Soit $X$ normal, $F$ fermé, $x\notin F$. Comme $X$ est $T_1$,
$\{x\}$ est fermé. Les fermés disjoints $\{x\}$ et $F$ sont séparés
par une fonction d'Urysohn $f\colon X\to[0,1]$ avec $f(x)=0$ et
$f|_F\equiv 1$.
\end{proof}

%% -------------------------------------------------------
%%  9.3  Théorème d'extension de Tietze
%% -------------------------------------------------------
\section{Théorème d'extension de Tietze}
\label{sec:tietze}
\index{théorème de Tietze}\index{extension!théorème de Tietze}

\begin{theoreme}[Tietze--Urysohn]
\label{thm:tietze}
Soit $X$ un espace topologique normal et $A\subset X$ un sous-ensemble
fermé. Toute application continue $f\colon A\to [a,b]$
(resp.\ $f\colon A\to\mathbf{R}$) admet une extension continue
$\tilde{f}\colon X\to [a,b]$ (resp.\ $\tilde{f}\colon X\to\mathbf{R}$).
\end{theoreme}

\begin{proof}
Il suffit de traiter le cas $f\colon A\to[-1,1]$, le cas général
$[a,b]$ s'en déduisant par transformation affine, et le cas non borné
par composition avec un homéomorphisme $\mathbf{R}\xrightarrow{\;\sim\;}(-1,1)$.

\medskip\noindent
\textbf{Étape 1 : lemme d'approximation.}

Soit $g\colon A\to [-1,1]$ continue. Posons
\[
  B^{-}=g^{-1}\!\bigl([-1,-\tfrac{1}{3}]\bigr),\qquad
  B^{+}=g^{-1}\!\bigl([\tfrac{1}{3},1]\bigr).
\]
Ce sont deux fermés disjoints de $A$, donc de $X$ (car $A$ est fermé
dans $X$). Par le lemme d'Urysohn (théorème~\ref{thm:urysohn}),
il existe $h\colon X\to[-\frac{1}{3},\frac{1}{3}]$ continue avec
$h|_{B^{-}}\equiv -\frac{1}{3}$ et $h|_{B^{+}}\equiv \frac{1}{3}$.
Alors pour tout $a\in A$ :
\[
|g(a)-h(a)|\le \tfrac{2}{3}.
\]
En effet, si $g(a)\in[-1,-\frac{1}{3}]$, alors $h(a)=-\frac{1}{3}$
et $|g(a)-h(a)|\le 1-\frac{1}{3}=\frac{2}{3}$ ;
si $g(a)\in[-\frac{1}{3},\frac{1}{3}]$, alors
$|g(a)-h(a)|\le\frac{1}{3}+\frac{1}{3}=\frac{2}{3}$ ;
si $g(a)\in[\frac{1}{3},1]$, de même.

\medskip\noindent
\textbf{Étape 2 : construction par approximations successives.}

Posons $g_0=f$ et appliquons le lemme d'approximation :
on obtient $h_0\colon X\to[-\frac{1}{3},\frac{1}{3}]$ telle que
$\|g_0-h_0|_A\|_\infty\le\frac{2}{3}$.

Posons $g_1=f-h_0|_A\colon A\to [-\frac{2}{3},\frac{2}{3}]$.
On normalise : $\frac{3}{2}g_1\colon A\to[-1,1]$.
Le lemme fournit $h_1'\colon X\to[-\frac{1}{3},\frac{1}{3}]$
approchant $\frac{3}{2}g_1$ à $\frac{2}{3}$ près.
Posons $h_1=\frac{2}{3}h_1'$; alors
$h_1\colon X\to[-\frac{2}{9},\frac{2}{9}]$ et
$\|g_1-h_1|_A\|_\infty\le\bigl(\frac{2}{3}\bigr)^2$.

Par récurrence, on construit des fonctions continues
$h_n\colon X\to\mathbf{R}$ telles que
\[
\|h_n\|_\infty \le \frac{1}{3}\Bigl(\frac{2}{3}\Bigr)^n,
\qquad
\Bigl\|f-\sum_{k=0}^n h_k\Big|_A\Bigr\|_\infty
  \le \Bigl(\frac{2}{3}\Bigr)^{n+1}.
\]

\medskip\noindent
\textbf{Étape 3 : convergence et conclusion.}

La série $\tilde{f}(x)=\sum_{n=0}^{\infty}h_n(x)$ converge
uniformément sur $X$, car
\[
\sum_{n=0}^{\infty}\|h_n\|_\infty
  \le \sum_{n=0}^{\infty}\frac{1}{3}\Bigl(\frac{2}{3}\Bigr)^n
  = \frac{1/3}{1-2/3}=1.
\]
Comme limite uniforme de fonctions continues, $\tilde{f}$ est continue.
De plus, $\tilde{f}(x)\in[-1,1]$ pour tout $x\in X$ (la somme de la
série est majorée par $1$ en valeur absolue). Enfin, pour $a\in A$ :
\[
\tilde{f}(a)=\lim_{n\to\infty}\sum_{k=0}^n h_k(a)=f(a),
\]
car le reste tend vers $0$. Donc $\tilde{f}|_A=f$.

\medskip
Pour le cas $f\colon A\to\mathbf{R}$, on compose avec un homéomorphisme
$\varphi\colon\mathbf{R}\to(-1,1)$ (par exemple
$\varphi(t)=\frac{t}{1+|t|}$), on étend
$\varphi\circ f\colon A\to(-1,1)\subset[-1,1]$ en
$\widetilde{\varphi\circ f}\colon X\to[-1,1]$, puis on pose
$C=(\widetilde{\varphi\circ f})^{-1}(\{-1,1\})$. Comme $C\cap A
=\varnothing$ (car $\varphi\circ f$ ne prend pas les valeurs $\pm 1$)
et $C$ est fermé, on peut ajuster par le lemme d'Urysohn pour éviter
les valeurs $\pm 1$, ou directement considérer
$\tilde{f}=\varphi^{-1}\circ\widetilde{\varphi\circ f}$ sur
$X\setminus C$ et montrer que $C$ peut être rendu vide par un choix
convenable. On renvoie le lecteur à l'exercice~\ref{exo:tietze-non-borne}
pour les détails de ce cas.
\end{proof}

\begin{remarque}
\label{rem:tietze-reciproque}
La réciproque est vraie : si tout fermé d'un espace $T_1$ admet la
propriété d'extension de Tietze, alors l'espace est normal. En effet,
si $A$ et $B$ sont fermés disjoints, la fonction $f\colon A\cup B\to
[0,1]$ définie par $f|_A\equiv 0$, $f|_B\equiv 1$ est continue sur le
fermé $A\cup B$ ; une extension $\tilde{f}$ fournit la fonction d'Urysohn.
\end{remarque}

%% -------------------------------------------------------
%%  9.4  Partitions de l'unité
%% -------------------------------------------------------
\section{Partitions de l'unité}
\label{sec:partitions-unite}
\index{partition de l'unité}

\begin{definition}[Partition de l'unité]
\label{def:partition-unite}
Soit $X$ un espace topologique et $(U_i)_{i\in I}$ un recouvrement
ouvert de $X$. Une \emph{partition de l'unité subordonnée} à
$(U_i)_{i\in I}$ est une famille $(\varphi_i)_{i\in I}$ de fonctions
continues $\varphi_i\colon X\to[0,1]$ telles que :
\begin{enumerate}[label=(\roman*)]
  \item $\operatorname{Supp}(\varphi_i)\subset U_i$ pour tout $i\in I$ ;
  \item la famille $(\operatorname{Supp}(\varphi_i))_{i\in I}$ est
    localement finie\index{localement fini} ;
  \item $\sum_{i\in I}\varphi_i(x)=1$ pour tout $x\in X$.
\end{enumerate}
\end{definition}

\begin{theoreme}
\label{thm:partition-unite-paracompact}
Tout espace paracompact\index{espace paracompact!partition de l'unité}
et de Hausdorff admet une partition de l'unité subordonnée à tout
recouvrement ouvert.
\end{theoreme}

\begin{proof}[Esquisse de preuve]
Soit $(U_i)_{i\in I}$ un recouvrement ouvert de $X$ paracompact $T_2$.
Comme $X$ est paracompact $T_2$, il est normal. On extrait un
raffinement localement fini $(V_j)_{j\in J}$ et, pour chaque $j$,
un rétrécissement fermé $F_j\subset V_j$ (existence garantie par
la normalité et le caractère localement fini). Le lemme d'Urysohn
fournit $\psi_j\colon X\to[0,1]$ avec $\psi_j|_{F_j}\equiv 1$
et $\operatorname{Supp}(\psi_j)\subset V_j$. La somme
$\Psi=\sum_j\psi_j$ est bien définie, continue et strictement positive.
On pose $\varphi_j=\psi_j/\Psi$.
\end{proof}

%% -------------------------------------------------------
%%  9.5  Théorème de métrisation d'Urysohn
%% -------------------------------------------------------
\section{Théorème de métrisation d'Urysohn}
\label{sec:metrisation-urysohn}
\index{théorème de métrisation d'Urysohn}\index{Urysohn!théorème de métrisation}

\begin{theoreme}[Plongement dans $\mathbf{R}^\omega$]
\label{thm:plongement-Rw}
Tout espace topologique $T_3$ à base dénombrable\index{base dénombrable}
se plonge dans $\mathbf{R}^\omega=\mathbf{R}^{\mathbf{N}}$ muni de la
topologie produit.
\end{theoreme}

\begin{proof}
Soit $(B_n)_{n\ge 1}$ une base dénombrable de $X$.
Puisque $X$ est $T_3$ (régulier $T_1$) à base dénombrable, il est
normal (tout espace de Lindel\"of r\'egulier est normal\index{Lindelöf!normal}).

Pour chaque couple $(m,n)$ tel que $\overline{B_m}\subset B_n$,
le lemme d'Urysohn fournit une fonction continue
$f_{m,n}\colon X\to[0,1]$ avec $f_{m,n}|_{\overline{B_m}}\equiv 1$
et $f_{m,n}|_{X\setminus B_n}\equiv 0$. L'ensemble de ces couples
est dénombrable ; on les énumère en $(f_k)_{k\ge 1}$.

L'application $F\colon X\to\mathbf{R}^\omega$ définie par
$F(x)=(f_1(x),f_2(x),\ldots)$ est continue (chaque composante l'est).
Elle est injective : si $x\neq y$, il existe $B_n$ contenant $x$
mais pas $y$, et $B_m$ avec $x\in B_m\subset\overline{B_m}\subset B_n$
(par régularité) ; alors $f_{m,n}(x)=1$ et $f_{m,n}(y)=0$.

Il reste à montrer que $F$ est un homéomorphisme sur son image.
Soit $U$ ouvert dans $X$ et $x\in U$. Choisissons $B_m,B_n$ avec
$x\in B_m\subset\overline{B_m}\subset B_n\subset U$. Alors
$f_{m,n}^{-1}((0,+\infty))\supset B_m\ni x$ et
$F(f_{m,n}^{-1}((0,+\infty)))\subset F(B_n)\subset F(U)$.
L'image $F(U)$ est donc ouvert dans $F(X)$ pour la topologie induite.
\end{proof}

\begin{theoreme}[Métrisation d'Urysohn]
\label{thm:metrisation-urysohn}
Un espace topologique est métrisable\index{espace métrisable} si et
seulement s'il est $T_3$ et à base dénombrable.
\end{theoreme}

\begin{proof}
$(\Leftarrow)$ Par le théorème~\ref{thm:plongement-Rw},
$X$ se plonge dans $\mathbf{R}^\omega$. Or $\mathbf{R}^\omega$ est
métrisable, par exemple par la métrique
\[
d\bigl((x_n),(y_n)\bigr)
  = \sum_{n=1}^{\infty}\frac{1}{2^n}\,
    \frac{|x_n-y_n|}{1+|x_n-y_n|}.
\]
Tout sous-espace d'un espace métrisable est métrisable.

$(\Rightarrow)$ Tout espace métrisable est $T_3$.
Si de plus il est à base dénombrable (hypothèse nécessaire), c'est
immédiat. En fait, la condition complète est : un espace métrisable
est à base dénombrable si et seulement s'il est séparable
(Lindelöf $\Leftrightarrow$ séparable $\Leftrightarrow$ base
dénombrable dans le cas métrique).
\end{proof}

%% -------------------------------------------------------
%%  9.6  Contre-exemple : espace non normal
%% -------------------------------------------------------
\section{Un contre-exemple fondamental}
\label{sec:contre-exemple-non-normal}

\begin{exemple}[Le plan de Sorgenfrey]
\label{ex:plan-sorgenfrey-non-normal}
\index{Sorgenfrey!plan de}\index{espace non normal}
Considérons la droite de Sorgenfrey $\mathbf{R}_\ell$
(cf.\ annexe~\ref{app:sorgenfrey}) et son carré
$S=\mathbf{R}_\ell\times\mathbf{R}_\ell$.
L'espace $S$ est séparable (les rationnels forment un ensemble dense),
de Hausdorff, mais \emph{non normal}.

En effet, les fermés $A=\{(x,-x):x\in\mathbf{Q}\}$ et
$B=\{(x,-x):x\in\mathbf{R}\setminus\mathbf{Q}\}$ (restreints à
l'anti-diagonale) sont fermés disjoints dans $S$, mais ne peuvent
être séparés par des ouverts disjoints. En particulier, le lemme
d'Urysohn ne s'applique pas : il n'existe pas de fonction continue
$f\colon S\to[0,1]$ séparant $A$ et $B$.
\end{exemple}

%% -------------------------------------------------------
%%  9.7  Exemples
%% -------------------------------------------------------
\section{Exemples développés}
\label{sec:urysohn-exemples}

\begin{exemple}[Extension d'une fonction sur un fermé de $\mathbf{R}^n$]
\label{ex:extension-Rn}
\index{extension!dans $\mathbf{R}^n$}
Soit $F=\{0\}\cup\{1/n:n\ge 1\}\subset\mathbf{R}$ et $f\colon F\to
\mathbf{R}$ définie par $f(0)=0$ et $f(1/n)=(-1)^n/n$. Comme
$\mathbf{R}$ est normal et $F$ est fermé, le théorème de Tietze garantit
l'existence d'une extension continue $\tilde{f}\colon\mathbf{R}\to
\mathbf{R}$. On peut construire explicitement $\tilde{f}$ par
interpolation linéaire entre les points de $F$.
\end{exemple}

\begin{exemple}[Fonction d'Urysohn dans un espace métrique]
\label{ex:urysohn-metrique}
Dans $(\mathbf{R}^2,\|\cdot\|)$, soient $A=\overline{B}(0,1)$
(boule fermée) et $B=\mathbf{R}^2\setminus B(0,2)$ (complémentaire
de la boule ouverte). La fonction
\[
f(x)=\frac{d(x,A)}{d(x,A)+d(x,B)}
\]
est une fonction d'Urysohn explicite : $f|_A\equiv 0$, $f|_B\equiv 1$,
$f$ continue et à valeurs dans $[0,1]$.
\end{exemple}

\begin{exemple}[Application du théorème de Tietze aux espaces compacts]
\label{ex:tietze-compact}
\index{Tietze!espace compact}
Soit $X$ un espace compact de Hausdorff (donc normal) et $K\subset X$
un fermé. Toute fonction continue $f\colon K\to\mathbf{R}$ se prolonge
en une fonction continue $\tilde{f}\colon X\to\mathbf{R}$. De plus,
on peut imposer $\|\tilde{f}\|_\infty=\|f\|_\infty$ (la preuve du
théorème de Tietze le garantit).
\end{exemple}

%% -------------------------------------------------------
%%  9.8  Exercices
%% -------------------------------------------------------
\section{Exercices}
\label{sec:exo-urysohn-tietze}

\begin{exercice}[$\star$]
\label{exo:urysohn-metrique-direct}
Soit $(X,d)$ un espace métrique et $A,B$ deux fermés disjoints.
Montrer directement (sans utiliser le lemme d'Urysohn) que la fonction
$f(x)=\frac{d(x,A)}{d(x,A)+d(x,B)}$ est une fonction d'Urysohn
pour $A$ et $B$.
\end{exercice}

\begin{exercice}[$\star$]
\label{exo:T3half-sous-espace}
Montrer qu'un sous-espace d'un espace complètement régulier est
complètement régulier.\index{complètement régulier!sous-espace}
\end{exercice}

\begin{exercice}[$\star\star$]
\label{exo:tietze-non-borne}
Compléter la preuve du théorème de Tietze~\ref{thm:tietze} dans le cas
d'une fonction $f\colon A\to\mathbf{R}$ non bornée. \textit{Indication :}
utiliser l'homéomorphisme $\varphi\colon\mathbf{R}\to(-1,1)$,
$\varphi(t)=t/(1+|t|)$, et montrer que l'ensemble
$(\widetilde{\varphi\circ f})^{-1}(\{-1,1\})$ est fermé et disjoint
de $A$.
\end{exercice}

\begin{exercice}[$\star\star$]
\label{exo:urysohn-fonction-lipschitz}
Soit $(X,d)$ un espace métrique et $A$ un fermé. Montrer que la
fonction $x\mapsto d(x,A)$ est $1$-lipschitzienne. En déduire
une preuve directe du lemme d'Urysohn pour les espaces métriques.
\end{exercice}

\begin{exercice}[$\star\star$]
\label{exo:partition-unite-Rn}
Soit $(U_i)_{i=1}^n$ un recouvrement ouvert fini de $\mathbf{R}^n$.
Construire explicitement une partition de l'unité subordonnée à
$(U_i)$.
\index{partition de l'unité!dans $\mathbf{R}^n$}
\end{exercice}

\begin{exercice}[$\star\star$]
\label{exo:plongement-compact-metrisable}
Soit $X$ un espace compact métrisable. Montrer que $X$ se plonge dans
le cube de Hilbert $[0,1]^{\mathbf{N}}$.
\end{exercice}

\begin{exercice}[$\star\star\star$]
\label{exo:sorgenfrey-non-normal}
Montrer en détail que le plan de Sorgenfrey
$\mathbf{R}_\ell\times\mathbf{R}_\ell$ n'est pas normal.
\textit{Indication :} considérer les fermés
$\{(x,-x):x\in\mathbf{Q}\}$ et $\{(x,-x):x\in\mathbf{R}\setminus
\mathbf{Q}\}$ et utiliser un argument de catégorie de Baire.
\index{Sorgenfrey!non normalité du plan}
\end{exercice}

\begin{exercice}[$\star\star\star$]
\label{exo:tietze-implique-normal}
Montrer la réciproque du théorème de Tietze : si pour tout fermé
$A$ d'un espace $T_1$ $X$, toute application continue
$f\colon A\to[0,1]$ admet une extension continue
$\tilde{f}\colon X\to[0,1]$, alors $X$ est normal.
\end{exercice}

%% -------------------------------------------------------
%%  9.9  Résumé du chapitre
%% -------------------------------------------------------
\section*{Résumé du chapitre~\thechapter}
\addcontentsline{toc}{section}{Résumé du chapitre~\thechapter}

\begin{itemize}
  \item Un espace est \emph{complètement régulier} ($T_{3\frac{1}{2}}$)
    si les points et les fermés sont séparés par des fonctions continues
    à valeurs dans $[0,1]$.
  \item Le \textbf{lemme d'Urysohn} affirme que dans un espace
    \emph{normal}, deux fermés disjoints quelconques sont séparés par
    une telle fonction. La construction utilise une famille d'ouverts
    indexée par les rationnels dyadiques.
  \item Le \textbf{théorème de Tietze} étend toute fonction continue
    définie sur un fermé d'un espace normal à l'espace tout entier,
    en préservant les bornes.
  \item Les partitions de l'unité, construites via le lemme d'Urysohn,
    sont un outil central en géométrie différentielle et en analyse.
  \item Le \textbf{théorème de métrisation d'Urysohn} caractérise les
    espaces métrisables à base dénombrable : ce sont exactement les
    espaces $T_3$ à base dénombrable.
\end{itemize}

%% ============================================================
%%  CHAPITRE 10 — Aperçu sur la Topologie Algébrique
%% ============================================================
\chapter{Aperçu sur la topologie algébrique}
\label{chap:topologie-algebrique}
\index{topologie algébrique}

\medskip
\noindent
\textit{La topologie générale permet de définir les notions de continuité,
de compacité, de connexité, mais elle peine à distinguer des espaces
qui sont pourtant « essentiellement différents ». Comment prouver
rigoureusement que $\mathbf{R}^2\not\cong\mathbf{R}^3$, ou que la
sphère $S^2$ n'est pas homéomorphe au tore $T^2$~? La topologie
algébrique répond à ces questions en associant à chaque espace
topologique des invariants algébriques~--- groupes, anneaux, modules~---
qui sont préservés par homéomorphisme.}

%% -------------------------------------------------------
%%  10.1  Homotopie
%% -------------------------------------------------------
\section{Homotopie de applications}
\label{sec:homotopie}
\index{homotopie}

\begin{definition}[Homotopie]
\label{def:homotopie}
Soient $X,Y$ des espaces topologiques et $f,g\colon X\to Y$ des
applications continues. On dit que $f$ et $g$ sont \emph{homotopes},
et l'on note $f\simeq g$, s'il existe une application continue
$H\colon X\times[0,1]\to Y$ telle que $H(x,0)=f(x)$ et $H(x,1)=g(x)$
pour tout $x\in X$. L'application $H$ est appelée une \emph{homotopie}
entre $f$ et $g$.
\end{definition}

\begin{proposition}
\label{prop:homotopie-equivalence}
La relation d'homotopie est une relation d'équivalence sur l'ensemble
$\mathcal{C}(X,Y)$ des applications continues de $X$ dans $Y$.
\end{proposition}

\begin{proof}
\emph{Réflexivité :} $H(x,t)=f(x)$. \emph{Symétrie :}
$H'(x,t)=H(x,1-t)$. \emph{Transitivité :} si $f\simeq g$ via $H_1$
et $g\simeq h$ via $H_2$, on pose
\[
H(x,t)=\begin{cases} H_1(x,2t) & \text{si } t\in[0,\tfrac{1}{2}],\\
  H_2(x,2t-1) & \text{si } t\in[\tfrac{1}{2},1]. \end{cases}
\]
La continuité résulte du lemme de recollement.
\end{proof}

\begin{definition}[Équivalence d'homotopie]
\label{def:equivalence-homotopie}
\index{équivalence d'homotopie}
Deux espaces $X$ et $Y$ sont dits \emph{homotopiquement équivalents}
(ou du même \emph{type d'homotopie}) s'il existe des applications
continues $f\colon X\to Y$ et $g\colon Y\to X$ telles que
$g\circ f\simeq\mathrm{id}_X$ et $f\circ g\simeq\mathrm{id}_Y$.
Un espace homotopiquement équivalent à un point est dit
\emph{contractile}\index{espace contractile}.
\end{definition}

\begin{exemple}
\label{ex:Rn-contractile}
L'espace $\mathbf{R}^n$ est contractile : l'homotopie
$H(x,t)=(1-t)x$ contracte $\mathbf{R}^n$ sur l'origine.
De même, tout espace convexe de $\mathbf{R}^n$ est contractile.
\end{exemple}

%% -------------------------------------------------------
%%  10.2  Groupe fondamental
%% -------------------------------------------------------
\section{Le groupe fondamental}
\label{sec:groupe-fondamental}
\index{groupe fondamental}\index{pi1@$\pi_1$}

\begin{definition}[Lacet, homotopie de lacets]
\label{def:lacet}
\index{lacet}
Soit $X$ un espace topologique et $x_0\in X$ (le \emph{point base}).
Un \emph{lacet} en $x_0$ est une application continue
$\gamma\colon[0,1]\to X$ telle que $\gamma(0)=\gamma(1)=x_0$.
Deux lacets $\gamma,\delta$ en $x_0$ sont \emph{homotopes relativement
aux extrémités} (on note $\gamma\simeq\delta\ \mathrm{rel}\;\{0,1\}$)
s'il existe une homotopie $H\colon[0,1]\times[0,1]\to X$ avec
$H(s,0)=\gamma(s)$, $H(s,1)=\delta(s)$, $H(0,t)=H(1,t)=x_0$ pour
tout $s,t$.
\end{definition}

\begin{definition}[Groupe fondamental]
\label{def:groupe-fondamental}
Le \emph{groupe fondamental}\index{groupe fondamental!définition}
de $X$ en $x_0$ est l'ensemble
\[
\pi_1(X,x_0)=\{\text{classes d'homotopie de lacets en }x_0\}
\]
muni de la loi de composition $[\gamma]\cdot[\delta]=[\gamma*\delta]$,
où la \emph{concaténation} est définie par
\[
(\gamma*\delta)(s) =
\begin{cases}
\gamma(2s) & \text{si } s\in[0,\tfrac{1}{2}],\\
\delta(2s-1) & \text{si } s\in[\tfrac{1}{2},1].
\end{cases}
\]
\end{definition}

\begin{theoreme}
\label{thm:pi1-groupe}
$(\pi_1(X,x_0),\cdot)$ est un groupe. L'élément neutre est la classe
du lacet constant $e(s)=x_0$, et l'inverse de $[\gamma]$ est
$[\bar{\gamma}]$ où $\bar{\gamma}(s)=\gamma(1-s)$.
\end{theoreme}

\begin{proof}[Esquisse]
L'associativité résulte d'une reparamétrisation continue : on montre
que $(\gamma*\delta)*\eta\simeq\gamma*(\delta*\eta)$ par une
homotopie qui ajuste les vitesses de parcours.
Le lacet constant est neutre : $\gamma*e\simeq\gamma$ (on « ralentit »
à la fin) et $e*\gamma\simeq\gamma$. Enfin,
$\gamma*\bar{\gamma}\simeq e$ : l'homotopie parcourt $\gamma$ jusqu'au
temps $t$, puis revient, et se contracte sur le point base quand $t\to 1$.
\end{proof}

\begin{proposition}
\label{prop:pi1-changement-base}
\index{groupe fondamental!changement de point base}
Si $X$ est connexe par arcs, alors $\pi_1(X,x_0)\cong\pi_1(X,x_1)$
pour tous $x_0,x_1\in X$. L'isomorphisme est donné par conjugaison
par un chemin reliant $x_0$ à $x_1$.
\end{proposition}

\begin{definition}[Espace simplement connexe]
\label{def:simplement-connexe}
\index{simplement connexe}
Un espace $X$ connexe par arcs est \emph{simplement connexe} si
$\pi_1(X,x_0)=\{e\}$ (le groupe trivial) pour un (donc tout) point
base $x_0$.
\end{definition}

\begin{exemple}
\label{ex:simplement-connexes}
Les espaces $\mathbf{R}^n$, $S^n$ pour $n\ge 2$, et tout espace
convexe sont simplement connexes. En revanche, le cercle $S^1$ ne
l'est pas.
\end{exemple}

%% -------------------------------------------------------
%%  10.3  Groupe fondamental du cercle
%% -------------------------------------------------------
\section{Groupe fondamental du cercle}
\label{sec:pi1-cercle}
\index{groupe fondamental!du cercle}

\begin{theoreme}
\label{thm:pi1-S1}
$\pi_1(S^1, 1)\cong\mathbf{Z}$.
\end{theoreme}

\begin{proof}[Esquisse via le revêtement $p\colon\mathbf{R}\to S^1$]
\index{revêtement!de $S^1$}
Considérons l'application de revêtement
$p\colon\mathbf{R}\to S^1$, $p(t)=e^{2\pi it}$.

\textbf{1. Relèvement des lacets.}
Tout lacet $\gamma\colon[0,1]\to S^1$ basé en $1=p(0)$ admet un unique
relèvement $\tilde{\gamma}\colon[0,1]\to\mathbf{R}$ tel que
$\tilde{\gamma}(0)=0$ et $p\circ\tilde{\gamma}=\gamma$.
Le point $\tilde{\gamma}(1)\in p^{-1}(1)=\mathbf{Z}$ est un entier
appelé le \emph{degré}\index{degré d'un lacet} du lacet.

\textbf{2. Invariance par homotopie.}
Si $\gamma\simeq\delta\ \mathrm{rel}\;\{0,1\}$, l'homotopie se relève
en une homotopie entre $\tilde{\gamma}$ et $\tilde{\delta}$, toutes
deux d'origine $0$. Par connexité de $[0,1]$ et le fait que
$p^{-1}(1)$ est discret, $\tilde{\gamma}(1)=\tilde{\delta}(1)$.
Le degré ne dépend donc que de la classe d'homotopie.

\textbf{3. Homomorphisme.}
L'application $\deg\colon\pi_1(S^1,1)\to\mathbf{Z}$,
$[\gamma]\mapsto\tilde{\gamma}(1)$ est un homomorphisme de groupes :
le relèvement de $\gamma*\delta$ est la concaténation des relèvements,
d'où $\deg(\gamma*\delta)=\deg(\gamma)+\deg(\delta)$.

\textbf{4. Bijectivité.}
Surjectivité : le lacet $\omega_n(s)=e^{2\pi ins}$ a pour degré $n$.
Injectivité : si $\deg(\gamma)=0$, alors $\tilde{\gamma}$ est un lacet
dans $\mathbf{R}$ (espace contractile), d'où
$\tilde{\gamma}\simeq \tilde{e}$ dans $\mathbf{R}$, et par projection,
$\gamma\simeq e$ dans $S^1$.
\end{proof}

%% -------------------------------------------------------
%%  10.4  Revêtements
%% -------------------------------------------------------
\section{Espaces de revêtement}
\label{sec:revetements}
\index{espace de revêtement}\index{revêtement}

\begin{definition}[Revêtement]
\label{def:revetement}
Un \emph{revêtement} d'un espace topologique $X$ est un couple
$(\tilde{X},p)$ où $\tilde{X}$ est un espace topologique et
$p\colon\tilde{X}\to X$ est une application continue surjective telle
que tout point $x\in X$ possède un voisinage ouvert $U$ dont la
préimage $p^{-1}(U)$ est une réunion disjointe d'ouverts de
$\tilde{X}$, chacun homéomorphe à $U$ via $p$.
\end{definition}

\begin{exemple}
\label{ex:revetements-classiques}
\begin{enumerate}[label=(\alph*)]
  \item $p\colon\mathbf{R}\to S^1$, $t\mapsto e^{2\pi it}$ est le
    revêtement universel de $S^1$.
  \item $p\colon S^n\to\mathbf{R}P^n$, $x\mapsto[x]$ est un
    revêtement à deux feuillets ($n\ge 1$).
  \item $p\colon\mathbf{R}^2\to T^2$, $(x,y)\mapsto(e^{2\pi ix},
    e^{2\pi iy})$ est le revêtement universel du tore.
\end{enumerate}
\end{exemple}

\begin{theoreme}[Propriété de relèvement des chemins]
\label{thm:relevement-chemins}
\index{relèvement!des chemins}
Soit $p\colon(\tilde{X},\tilde{x}_0)\to(X,x_0)$ un revêtement.
Pour tout chemin $\gamma\colon[0,1]\to X$ avec $\gamma(0)=x_0$,
il existe un unique chemin $\tilde{\gamma}\colon[0,1]\to\tilde{X}$
tel que $\tilde{\gamma}(0)=\tilde{x}_0$ et
$p\circ\tilde{\gamma}=\gamma$.
\end{theoreme}

\begin{theoreme}[Relèvement des homotopies]
\label{thm:relevement-homotopies}
\index{relèvement!des homotopies}
Soit $p\colon\tilde{X}\to X$ un revêtement, $H\colon[0,1]^2\to X$
une homotopie et $\tilde{H}_0\colon[0,1]\to\tilde{X}$ un relèvement
de $H(\cdot,0)$. Il existe une unique homotopie
$\tilde{H}\colon[0,1]^2\to\tilde{X}$ relevant $H$ avec
$\tilde{H}(\cdot,0)=\tilde{H}_0$.
\end{theoreme}

% --- TikZ : Revêtement de S^1 ---
\begin{figure}[ht]
\centering
\begin{tikzpicture}[>=stealth, scale=0.9]
  % Hélicoïde (R)
  \begin{scope}[shift={(0,0)}]
    \draw[thick, blue!70!black, domain=0:720, samples=200]
      plot ({1.2*cos(\x)}, {0.03*\x + 0.4*sin(\x)*0.15});
    % Axe vertical
    \draw[->, thin, gray] (0,-0.5) -- (0,5) node[right] {$\mathbf{R}$};
    \node[left] at (-1.4,2.5) {$\tilde{X}=\mathbf{R}$};
    % Points entiers
    \foreach \k in {0,1,2} {
      \fill[red] ({1.2*cos(360*\k)}, {0.03*360*\k + 0.4*sin(360*\k)*0.15})
        circle (2pt);
    }
  \end{scope}
  % Flèche p
  \draw[->, thick] (1.8, 1.5) -- (3.5, 1.5) node[midway, above] {$p$};
  % Cercle S^1
  \begin{scope}[shift={(5,1.5)}]
    \draw[thick] (0,0) circle (1.2);
    \fill[red] (1.2,0) circle (2pt) node[right] {$1$};
    \node[below] at (0,-1.5) {$X=S^1$};
  \end{scope}
\end{tikzpicture}
\caption{Le revêtement $p\colon\mathbf{R}\to S^1$, $t\mapsto e^{2\pi it}$.
Les points rouges sont les préimages de $1\in S^1$.}
\label{fig:revetement-S1}
\end{figure}

% --- TikZ : Lacets sur le tore ---
\begin{figure}[ht]
\centering
\begin{tikzpicture}[scale=1.2]
  % Tore schématique (carré identifié)
  \draw[thick] (0,0) rectangle (3,3);
  \draw[thick, ->] (0,0) -- (1.5,0);
  \draw[thick, ->] (3,0) -- (3,1.5);
  \draw[thick, ->] (0,0) -- (0,1.5);
  \draw[thick, ->] (0,3) -- (1.5,3);
  \node[below] at (1.5,-0.2) {$a$};
  \node[left] at (-0.2,1.5) {$b$};
  % Lacet a
  \draw[red, very thick, ->, domain=0:0.95, samples=50]
    plot ({3*\x}, {0.5});
  \node[red, above] at (1.5, 0.6) {$\alpha$};
  % Lacet b
  \draw[blue, very thick, ->, domain=0:0.95, samples=50]
    plot ({0.5}, {3*\x});
  \node[blue, right] at (0.6, 1.5) {$\beta$};
  % Point base
  \fill (0.5,0.5) circle (2pt) node[below right] {$x_0$};
  % Légende
  \node[right] at (4,2.5) {$\pi_1(T^2,x_0)\cong\mathbf{Z}^2$};
  \node[right] at (4,1.8) {engendré par $[\alpha]$ et $[\beta]$};
\end{tikzpicture}
\caption{Lacets générateurs du groupe fondamental du tore $T^2$,
représenté comme carré aux côtés identifiés.}
\label{fig:lacets-tore}
\end{figure}

%% -------------------------------------------------------
%%  10.5  Théorème du point fixe de Brouwer (dim 2)
%% -------------------------------------------------------
\section{Théorème du point fixe de Brouwer en dimension $2$}
\label{sec:brouwer-dim2}
\index{Brouwer!théorème du point fixe}\index{point fixe!Brouwer}

\begin{theoreme}[Brouwer, $n=2$]
\label{thm:brouwer-dim2}
Toute application continue $f\colon\overline{D}^2\to\overline{D}^2$
du disque fermé dans lui-même possède un point fixe.
\end{theoreme}

\begin{proof}
Supposons par l'absurde que $f(x)\neq x$ pour tout
$x\in\overline{D}^2$. Pour chaque $x$, considérons la demi-droite
issue de $f(x)$ passant par $x$ : elle intersecte le cercle $S^1$ en
un unique point $r(x)$. L'application $r\colon\overline{D}^2\to S^1$
ainsi définie est continue (par une formule explicite faisant
intervenir une équation du second degré dont on choisit la racine
correspondante de manière continue).

De plus, si $x\in S^1$, alors la demi-droite issue de $f(x)\in
\overline{D}^2$ passant par $x\in S^1$ rencontre $S^1$ en $x$ lui-même
(car $x$ est déjà sur le bord), d'où $r|_{S^1}=\mathrm{id}_{S^1}$.

Ainsi $r$ est une \emph{rétraction}\index{rétraction} de
$\overline{D}^2$ sur $S^1$. L'inclusion $\iota\colon S^1\hookrightarrow
\overline{D}^2$ vérifie $r\circ\iota=\mathrm{id}_{S^1}$, d'où par
fonctorialité de $\pi_1$ :
\[
r_*\circ\iota_* = \mathrm{id}\colon
\pi_1(S^1)\to\pi_1(S^1).
\]
Or $\iota_*\colon\pi_1(S^1)\to\pi_1(\overline{D}^2)$ envoie
$\mathbf{Z}$ dans le groupe trivial $\{0\}$ (car $\overline{D}^2$
est contractile), donc $r_*\circ\iota_*$ est l'application nulle.
Contradiction avec $r_*\circ\iota_*=\mathrm{id}_{\mathbf{Z}}$.
\end{proof}

\begin{theoreme}[Borsuk--Ulam, $n=1$]
\label{thm:borsuk-ulam-1}
\index{Borsuk--Ulam}
Pour toute application continue $f\colon S^1\to\mathbf{R}$, il existe
un point $x\in S^1$ tel que $f(x)=f(-x)$.
\end{theoreme}

\begin{proof}
La fonction $g(x)=f(x)-f(-x)$ est continue et vérifie $g(-x)=-g(x)$.
Par le théorème des valeurs intermédiaires (car $S^1$ est connexe),
$g$ s'annule.
\end{proof}

\begin{remarque}[Borsuk--Ulam en dimension $n$]
\label{rem:borsuk-ulam-n}
Le théorème de Borsuk--Ulam en dimension $n$ affirme que pour toute
application continue $f\colon S^n\to\mathbf{R}^n$, il existe $x\in S^n$
tel que $f(x)=f(-x)$. La preuve utilise des outils d'homologie ou de
cohomologie.\index{Borsuk--Ulam!dimension $n$}
\end{remarque}

%% -------------------------------------------------------
%%  10.6  Groupes d'homotopie supérieurs et homologie
%% -------------------------------------------------------
\section{Groupes d'homotopie supérieurs et homologie}
\label{sec:homotopie-superieure-homologie}

\begin{definition}[Groupes d'homotopie supérieurs]
\label{def:pin}
\index{groupes d'homotopie!supérieurs}\index{pin@$\pi_n$}
Pour $n\ge 1$, le \emph{$n$-ième groupe d'homotopie} de $(X,x_0)$ est
\[
\pi_n(X,x_0) = \bigl[(S^n,*),\,(X,x_0)\bigr],
\]
l'ensemble des classes d'homotopie (pointée) d'applications continues
de $(S^n,*)$ dans $(X,x_0)$, muni d'une structure de groupe
(abélien pour $n\ge 2$).
\end{definition}

\begin{remarque}
\label{rem:pi-n-exemples}
Le calcul des groupes $\pi_n$ est en général très difficile. Citons :
$\pi_n(S^n)\cong\mathbf{Z}$ pour tout $n\ge 1$ ;
$\pi_3(S^2)\cong\mathbf{Z}$ (fibration de Hopf\index{Hopf!fibration de}).
\end{remarque}

\begin{definition}[Homologie singulière --- aperçu]
\label{def:homologie-apercu}
\index{homologie singulière}
L'homologie singulière associe à tout espace topologique $X$ une suite
de groupes abéliens $(H_n(X;\mathbf{Z}))_{n\ge 0}$, définis à l'aide
de \emph{simplexes singuliers} (applications continues
$\sigma\colon\Delta^n\to X$) et de \emph{complexes de chaînes}.
L'homologie est un invariant d'homotopie plus calculable que les
groupes $\pi_n$.
\end{definition}

\begin{definition}[Caractéristique d'Euler]
\label{def:euler}
\index{caractéristique d'Euler}
Si $X$ est un espace dont les groupes d'homologie sont de type fini et
nuls sauf un nombre fini, la \emph{caractéristique d'Euler} est
\[
\chi(X)=\sum_{n=0}^{\infty}(-1)^n\operatorname{rang}(H_n(X;\mathbf{Z})).
\]
Par exemple, $\chi(S^n)=1+(-1)^n$, $\chi(T^2)=0$,
$\chi(\Sigma_g)=2-2g$ pour la surface orientable de genre $g$.
\end{definition}

%% -------------------------------------------------------
%%  10.7  Notes historiques
%% -------------------------------------------------------
\section{Notes historiques}
\label{sec:notes-historiques-topo-alg}

\begin{remarque}[Repères historiques]
\label{rem:histoire-topo-alg}
\index{Poincaré}\index{Brouwer!L.E.J.}\index{Hurewicz}
\begin{itemize}
  \item \textbf{Henri Poincaré} (1895) introduit l'\emph{Analysis Situs},
    fondant la topologie algébrique. Il définit le groupe fondamental,
    les nombres de Betti, et formule la conjecture de Poincaré
    (démontrée par Perelman en 2003).
  \item \textbf{L.\,E.\,J.\ Brouwer} (vers 1910--1912) démontre le
    théorème du point fixe, le théorème d'invariance du domaine et
    le théorème d'invariance de la dimension.
  \item \textbf{Witold Hurewicz} (1935--1936) introduit les groupes
    d'homotopie supérieurs $\pi_n$ et démontre le théorème de
    Hurewicz reliant $\pi_n$ et $H_n$.
  \item \textbf{Samuel Eilenberg} et \textbf{Norman Steenrod} (1945)
    axiomatisent la théorie de l'homologie.
\end{itemize}
\end{remarque}

%% -------------------------------------------------------
%%  10.8  Exercices
%% -------------------------------------------------------
\section{Exercices}
\label{sec:exo-topo-alg}

\begin{exercice}[$\star$]
\label{exo:homotopie-convexe}
Montrer que deux applications continues $f,g\colon X\to\mathbf{R}^n$
sont toujours homotopes. \textit{Indication :} homotopie linéaire.
\end{exercice}

\begin{exercice}[$\star$]
\label{exo:contractile-simplement-connexe}
Montrer que tout espace contractile est simplement connexe.
\end{exercice}

\begin{exercice}[$\star\star$]
\label{exo:pi1-produit}
Montrer que $\pi_1(X\times Y,(x_0,y_0))\cong\pi_1(X,x_0)\times
\pi_1(Y,y_0)$. En déduire $\pi_1(T^2)\cong\mathbf{Z}^2$.
\index{groupe fondamental!du tore}
\end{exercice}

\begin{exercice}[$\star\star$]
\label{exo:pi1-figure-huit}
En admettant le théorème de van Kampen\index{van Kampen}, calculer
$\pi_1(S^1\vee S^1)\cong\mathbf{Z}*\mathbf{Z}$ (groupe libre à deux
générateurs).
\end{exercice}

\begin{exercice}[$\star\star$]
\label{exo:retraction-Dn}
Montrer qu'il n'existe pas de rétraction continue
$r\colon\overline{D}^n\to S^{n-1}$ pour $n\ge 1$. \textit{Indication
pour $n=2$ :} utiliser $\pi_1$.
\end{exercice}

\begin{exercice}[$\star\star\star$]
\label{exo:brouwer-applications}
Déduire du théorème de Brouwer en dimension $2$ que tout champ de
vecteurs continu $v\colon\overline{D}^2\to\mathbf{R}^2$ pointant
vers l'intérieur sur $S^1$ possède un zéro.
\end{exercice}

%% -------------------------------------------------------
%%  10.9  Résumé du chapitre
%% -------------------------------------------------------
\section*{Résumé du chapitre~\thechapter}
\addcontentsline{toc}{section}{Résumé du chapitre~\thechapter}

\begin{itemize}
  \item L'\textbf{homotopie} est une déformation continue d'applications ;
    elle induit une relation d'équivalence et la notion d'équivalence
    d'homotopie entre espaces.
  \item Le \textbf{groupe fondamental} $\pi_1(X,x_0)$ est le groupe des
    classes d'homotopie de lacets basés en $x_0$. C'est un invariant
    topologique.
  \item $\pi_1(S^1)\cong\mathbf{Z}$ : la preuve utilise le revêtement
    $\mathbf{R}\to S^1$.
  \item Le théorème du \textbf{point fixe de Brouwer} (dim $2$) se
    déduit de l'impossibilité de rétracter le disque sur son bord,
    elle-même conséquence de $\pi_1(S^1)\cong\mathbf{Z}$.
  \item Les \textbf{groupes d'homotopie supérieurs} $\pi_n$ et
    l'\textbf{homologie singulière} $H_n$ sont des invariants
    algébriques plus fins, indispensables en topologie moderne.
\end{itemize}

%% ============================================================
%%  ANNEXE A — Répertoire de Contre-exemples
%% ============================================================
\appendix
\chapter{Répertoire de contre-exemples}
\label{app:contre-exemples}
\index{contre-exemples!répertoire}

\noindent
\textit{Nous rassemblons ici les principaux espaces et constructions
qui servent de contre-exemples en topologie générale. Pour chaque
espace, nous indiquons la définition, les propriétés de séparation,
la compacité, la connexité, et les énoncés classiques pour lesquels
il constitue un contre-exemple.}

\bigskip

%% --- Sorgenfrey ---
\section{Droite de Sorgenfrey}
\label{app:sorgenfrey}
\index{Sorgenfrey!droite de}

\textbf{Définition.} La \emph{droite de Sorgenfrey} $\mathbf{R}_\ell$
est l'ensemble $\mathbf{R}$ muni de la topologie engendrée par les
intervalles semi-ouverts $[a,b)$, $a<b$.

\textbf{Propriétés.}
\begin{itemize}
  \item Séparation : $T_{3\frac{1}{2}}$ (en fait parfaitement normal,
    donc $T_4$ et même $T_6$).
  \item Première dénombrabilité : oui.
  \item Seconde dénombrabilité : non (mais séparable et de Lindelöf).
  \item Compacité : non compact, non localement compact.
  \item Connexité : totalement séparé.
  \item Métrisabilité : non métrisable (séparable mais non à base
    dénombrable).
\end{itemize}

\textbf{Contre-exemples.}
\begin{itemize}
  \item Montre que séparable et de Lindelöf n'impliquent pas
    la seconde dénombrabilité.
  \item Montre que la normalité n'est pas préservée par produit
    (voir plan de Sorgenfrey ci-dessous).
\end{itemize}

%% --- Plan de Sorgenfrey ---
\section{Plan de Sorgenfrey}
\label{app:plan-sorgenfrey}
\index{Sorgenfrey!plan de}

\textbf{Définition.} $S=\mathbf{R}_\ell\times\mathbf{R}_\ell$,
produit de deux copies de la droite de Sorgenfrey.

\textbf{Propriétés.}
\begin{itemize}
  \item Séparation : $T_{3\frac{1}{2}}$ mais \emph{non normal}.
  \item Séparable : oui ($\mathbf{Q}^2$ est dense).
  \item De Lindelöf : non.
  \item Compacité : non compact.
\end{itemize}

\textbf{Contre-exemples.}
\begin{itemize}
  \item Contre-exemple canonique pour : le produit de deux espaces
    normaux n'est pas nécessairement normal.
  \item Contre-exemple pour : le produit de deux espaces de Lindelöf
    n'est pas nécessairement de Lindelöf.
  \item Le lemme d'Urysohn échoue dans cet espace.
\end{itemize}

%% --- Topologie cofinie ---
\section{Topologie cofinie}
\label{app:cofinie}
\index{topologie cofinie}

\textbf{Définition.} Sur un ensemble infini $X$, la topologie cofinie
est $\tau=\{\varnothing\}\cup\{U\subset X: X\setminus U\text{ est fini}\}$.

\textbf{Propriétés.}
\begin{itemize}
  \item Séparation : $T_1$ mais non $T_2$.
  \item Compacité : compact (tout recouvrement ouvert admet un
    sous-recouvrement fini).
  \item Connexité : connexe (tout ouvert non vide est cofiniment grand).
  \item Métrisabilité : non métrisable.
\end{itemize}

\textbf{Contre-exemples.}
\begin{itemize}
  \item Espace compact non séparé.
  \item Montre que compact et $T_1$ n'impliquent pas $T_2$.
\end{itemize}

%% --- Topologie codénombrable ---
\section{Topologie codénombrable}
\label{app:codenombrable}
\index{topologie codénombrable}

\textbf{Définition.} Sur un ensemble non dénombrable $X$, la topologie
codénombrable est
$\tau=\{\varnothing\}\cup\{U\subset X: X\setminus U\text{ est
dénombrable}\}$.

\textbf{Propriétés.}
\begin{itemize}
  \item Séparation : $T_1$ mais non $T_2$.
  \item Compacité : non compact en général.
  \item Connexité : connexe.
\end{itemize}

\textbf{Contre-exemples.}
\begin{itemize}
  \item Espace $T_1$, connexe, non séparé.
  \item Les suites ne suffisent pas à décrire la topologie.
\end{itemize}

%% --- Sierpiński ---
\section{Espace de Sierpiński}
\label{app:sierpinski}
\index{Sierpiński!espace de}

\textbf{Définition.} $S=\{0,1\}$ avec $\tau=\{\varnothing,\{1\},\{0,1\}\}$.

\textbf{Propriétés.}
\begin{itemize}
  \item Séparation : $T_0$ mais non $T_1$.
  \item Compacité : compact (espace fini).
  \item Connexité : connexe.
\end{itemize}

\textbf{Contre-exemples.}
\begin{itemize}
  \item Le plus petit espace $T_0$ non $T_1$.
  \item Joue un rôle universel : les ouverts de $X$ correspondent aux
    applications continues $X\to S$.
\end{itemize}

%% --- Cantor ---
\section{Ensemble de Cantor}
\label{app:cantor}
\index{Cantor!ensemble de}

\textbf{Définition.} L'ensemble triadique de Cantor est
$C=\bigcap_{n=0}^{\infty}C_n$ où $C_0=[0,1]$ et $C_{n+1}$ est obtenu
en retirant le tiers central ouvert de chaque intervalle de $C_n$.
Alternativement, $C=\{x\in[0,1]: x=\sum_{k=1}^{\infty}a_k/3^k,\;
a_k\in\{0,2\}\}$.

\textbf{Propriétés.}
\begin{itemize}
  \item Séparation : $T_4$ (métrisable).
  \item Compacité : compact.
  \item Connexité : totalement séparé.
  \item Homéomorphe à $\{0,1\}^{\mathbf{N}}$ (produit dénombrable).
  \item Non dénombrable, de mesure de Lebesgue nulle.
\end{itemize}

\textbf{Contre-exemples.}
\begin{itemize}
  \item Compact non dénombrable sans point isolé, de mesure nulle.
  \item Espace universel : tout espace métrique compact totalement
    séparé est homéomorphe à un fermé de $C$.
\end{itemize}

%% --- Courbe du topologiste ---
\section{Courbe du topologiste (\textit{topologist's sine curve})}
\label{app:courbe-topologiste}
\index{courbe du topologiste}

\textbf{Définition.} $T=\{(0,y):-1\le y\le 1\}\cup
\{(x,\sin(1/x)):0<x\le 1\}$ muni de la topologie induite de
$\mathbf{R}^2$.

\textbf{Propriétés.}
\begin{itemize}
  \item Séparation : $T_4$ (métrisable).
  \item Compacité : compact.
  \item Connexité : connexe mais \emph{non connexe par arcs}.
\end{itemize}

\textbf{Contre-exemples.}
\begin{itemize}
  \item Contre-exemple canonique : connexe n'implique pas connexe par arcs.
\end{itemize}

%% --- Longue droite ---
\section{Longue droite}
\label{app:longue-droite}
\index{longue droite}

\textbf{Définition.} Soit $\omega_1$ le premier ordinal non dénombrable.
La \emph{longue droite} est $L=\omega_1\times[0,1)$ muni de la
topologie de l'ordre lexicographique, privé de son point minimal.

\textbf{Propriétés.}
\begin{itemize}
  \item Séparation : $T_5$ (complètement normal).
  \item Compacité : non compact, localement compact, séquentiellement
    compact, dénombrablement compact.
  \item Connexité : connexe, localement connexe par arcs.
  \item Non métrisable, non séparable, non à base dénombrable.
\end{itemize}

\textbf{Contre-exemples.}
\begin{itemize}
  \item Séquentiellement compact mais non compact.
  \item Dénombrablement compact mais non compact.
  \item Localement métrisable mais non métrisable globalement.
\end{itemize}

%% --- Double origine ---
\section{Droite à double origine (\textit{bug-eyed line})}
\label{app:double-origine}
\index{droite à double origine}\index{bug-eyed line}

\textbf{Définition.} Soit $X=(\mathbf{R}\setminus\{0\})\cup\{0^+,0^-\}$.
La topologie est définie en gardant la topologie usuelle sur
$\mathbf{R}\setminus\{0\}$ et en prenant comme voisinages de $0^\pm$
les ensembles $(-\varepsilon,0)\cup\{0^\pm\}\cup(0,\varepsilon)$.

\textbf{Propriétés.}
\begin{itemize}
  \item Séparation : $T_1$, localement homéomorphe à $\mathbf{R}$,
    mais \emph{non $T_2$} (les points $0^+$ et $0^-$ ne sont pas
    séparés par des ouverts disjoints).
  \item Première dénombrabilité : oui. Seconde dénombrabilité : oui.
  \item Connexité : connexe.
\end{itemize}

\textbf{Contre-exemples.}
\begin{itemize}
  \item Variété topologique (localement euclidienne) non séparée.
  \item Montre que la condition de Hausdorff est indépendante des
    conditions de dénombrabilité et de la structure locale.
\end{itemize}

%% --- Point particulier ---
\section{Topologie du point particulier}
\label{app:point-particulier}
\index{topologie du point particulier}

\textbf{Définition.} Soit $X$ un ensemble et $p\in X$ un point
distingué. La \emph{topologie du point particulier} est
$\tau=\{\varnothing\}\cup\{U\subset X: p\in U\}$.

\textbf{Propriétés.}
\begin{itemize}
  \item Séparation : $T_0$ mais non $T_1$ (les points $\neq p$ ne
    sont pas fermés).
  \item Compacité : compact (tout ouvert non vide contient $p$, donc
    un seul ouvert suffit).
  \item Connexité : connexe et connexe par arcs.
\end{itemize}

\textbf{Contre-exemples.}
\begin{itemize}
  \item Espace compact, connexe par arcs, non $T_1$.
  \item Illustre la distinction entre $T_0$ et $T_1$.
\end{itemize}

%% --- Moore (Niemytzki) ---
\section{Espace de Moore (plan de Niemytzki)}
\label{app:niemytzki}
\index{Moore!espace de}\index{Niemytzki!plan de}

\textbf{Définition.} Soit $X=\{(x,y)\in\mathbf{R}^2:y\ge 0\}$
(demi-plan supérieur fermé). Pour $(x,y)$ avec $y>0$, les voisinages
sont les disques ouverts habituels. Pour $(x,0)$ sur l'axe, les
voisinages de base sont les ensembles $\{(x,0)\}\cup D((x,\varepsilon),
\varepsilon)$ où $D((x,\varepsilon),\varepsilon)$ est le disque ouvert
de centre $(x,\varepsilon)$ et de rayon $\varepsilon$, tangent à
l'axe des abscisses en $(x,0)$.

\textbf{Propriétés.}
\begin{itemize}
  \item Séparation : $T_{3\frac{1}{2}}$ (complètement régulier) mais
    \emph{non normal}.
  \item Première dénombrabilité : oui.
  \item Séparable : oui ($\mathbf{Q}^2\cap X$ est dense).
  \item Seconde dénombrabilité : non.
  \item Compacité : non compact.
\end{itemize}

\textbf{Contre-exemples.}
\begin{itemize}
  \item Espace séparable, complètement régulier, non normal.
  \item Montre que $T_{3\frac{1}{2}}$ n'implique pas $T_4$.
\end{itemize}

%% --- Planche de Tychonoff ---
\section{Planche de Tychonoff}
\label{app:planche-tychonoff}
\index{Tychonoff!planche de}

\textbf{Définition.} Soit $\Omega=\omega_1+1=[0,\omega_1]$ et
$\Sigma=\omega+1=[0,\omega]$ munis de la topologie de l'ordre.
La \emph{planche de Tychonoff} est l'espace
$T=(\Omega\times\Sigma)\setminus\{(\omega_1,\omega)\}$, c'est-à-dire
le produit $\Omega\times\Sigma$ privé de son coin supérieur droit.

\textbf{Propriétés.}
\begin{itemize}
  \item Séparation : $T_{3\frac{1}{2}}$ (complètement régulier) mais
    \emph{non normal}.
  \item Compacité : non compact (mais localement compact).
  \item Connexité : connexe.
\end{itemize}

\textbf{Contre-exemples.}
\begin{itemize}
  \item Sous-espace d'un produit de deux espaces compacts qui n'est
    pas normal.
  \item Illustre que la normalité est une propriété fragile vis-à-vis
    des sous-espaces.
\end{itemize}

%% ============================================================
%%  BIBLIOGRAPHIE
%% ============================================================
\begin{thebibliography}{99}

\bibitem{munkres}
\textsc{Munkres}, J.\,R.,
\emph{Topology}, 2\textsuperscript{e} éd.,
Prentice Hall, Upper Saddle River, 2000.
\index{Munkres}

\bibitem{bourbaki}
\textsc{Bourbaki}, N.,
\emph{Topologie Générale}, chapitres 1 à 4,
Springer, Berlin, 2007 (réimpression).
\index{Bourbaki}

\bibitem{willard}
\textsc{Willard}, S.,
\emph{General Topology},
Dover Publications, Mineola, 2004 (réimpression de l'éd.\ 1970).
\index{Willard}

\bibitem{kelley}
\textsc{Kelley}, J.\,L.,
\emph{General Topology},
Springer, Graduate Texts in Mathematics \textbf{27}, 1975
(réimpression).
\index{Kelley}

\bibitem{dugundji}
\textsc{Dugundji}, J.,
\emph{Topology},
Allyn and Bacon, Boston, 1966.
\index{Dugundji}

\bibitem{dixmier}
\textsc{Dixmier}, J.,
\emph{Topologie Générale},
Presses Universitaires de France, collection « Mathématiques »,
Paris, 1981.
\index{Dixmier}

\bibitem{steen-seebach}
\textsc{Steen}, L.\,A. et \textsc{Seebach}, J.\,A.,
\emph{Counterexamples in Topology},
Dover Publications, Mineola, 1995 (réimpression de l'éd.\ 1978).
\index{Steen}\index{Seebach}

\bibitem{hatcher}
\textsc{Hatcher}, A.,
\emph{Algebraic Topology},
Cambridge University Press, Cambridge, 2002.
Disponible en ligne : \texttt{https://pi.math.cornell.edu/\~{}hatcher/AT/ATpage.html}.
\index{Hatcher}

\end{thebibliography}

%% ============================================================
%%  INDEX
%% ============================================================
\printindex

\end{document}
